\begingroup
\clearpage
\part{Absolute values and original-metric control}
\noindent This continuation integrates the nine new notes and their complete
receivers. The preceding editions remain recoverable byte for byte. Source
corrections and original-minor scope are stated at their exact use.
\providecommand{\tightlist}{\setlength{\itemsep}{0pt}\setlength{\parskip}{0pt}}
\providecommand{\C}{\mathbb C}\providecommand{\R}{\mathbb R}
\providecommand{\Z}{\mathbb Z}\providecommand{\Q}{\mathbb Q}
\providecommand{\N}{\mathbb N}
\providecommand{\im}{\operatorname{im}}
\providecommand{\diag}{\operatorname{diag}}
\providecommand{\Gr}{\operatorname{Gram}}
\providecommand{\Gram}{\operatorname{Gram}}
\providecommand{\Res}{\operatorname{Res}}
\providecommand{\cond}{\operatorname{cond}}
\providecommand{\ExtensionTok}[1]{#1}\providecommand{\NormalTok}[1]{#1}
\providecommand{\StringTok}[1]{#1}\providecommand{\BuiltInTok}[1]{#1}
\providecommand{\ControlFlowTok}[1]{#1}

\clearpage
\begingroup
\def\R{\mathbb R}
\def\C{\mathbb C}
\def\PP{\mathcal P}
\def\BB{\mathcal B}
\def\dd{\,d}
\def\Gram{\operatorname{Gram}}

\setcounter{equation}{0}
\renewcommand{\theequation}{A03AG.\arabic{equation}}
\section*{Absolute Gamma baseline and the complete canonical action}
\addcontentsline{toc}{section}{Absolute Gamma baseline and the complete canonical action}
% BEGIN PRESERVED INPUT 01_ABSOLUTE_GAMMA_AND_ACTION.tex SHA256 7532d29e183ee5f080d920d7b2f13520d953a39294b6d8a2ca6a39c730554250


\noindent The incoming note \texttt{07\_ABSOLUTE\_GAMMA\_BASELINE.md}
identified the previously unevaluated absolute Gamma coefficient. This proof
retains its complete root polynomial, original monic frames and source mass.
It strengthens its remainder to $O_h(\log q)$ for every fixed multiplicity.
The original four blocks are expressed using a single fixed potential and
three analytic factors, and a finite convexity argument controls the remaining
bounded root parameter. The previously proved arithmetic correction is then
returned to the absolute baseline and the complete canonical action.
The human analytic input is Borot--Guionnet \cite{a03ag:BG}; it is applied after
a comparison of the complete quadratic forms on both complementary regions.

\section{Original objects and the determinant map}
Fix $0<\delta<1/2$, $\gamma>2$ and $m_0\in\mathbb N$. For the arithmetic
receivers these are the stipulated complete quartet data of the preceding
programme, with $h(s)$ its actual quartet factor in $2\xi(s)$.
The Gamma calculation itself requires only these fixed parameters. Put
\begin{equation}\tag{AG1}
\begin{gathered}
 k=4\ell+1,\quad e=1+k(m_0-1),\quad q=e(k+1)^2,\quad c=k/2,\\
 \chi(S)=\prod_{a,b=0}^k
 [S-c-(2a-k)\delta-i(2b-k)\gamma]^e.
\end{gathered}
\end{equation}
The monic polynomial $Q(y)=i^{-q}\chi(c+iy)$ has its full roots
\begin{equation}\tag{AG2}
 \omega_{ab}=(2b-k)\gamma-i(2a-k)\delta
 \quad\hbox{with multiplicity }e.
\end{equation}
Conjugation and negation preserve the root multiset. Since $k$ is odd,
no root is real. Thus $Q$ is real, even and strictly positive on $\R$.
Here $q$ is divisible by four. Retain
\[
 \sigma(y)=\frac{|\Gamma(1/4+iy/2)|^2}{2\pi},\qquad
 \int_\R\sigma(y)\dd y=\sqrt{2\pi},\qquad \alpha=\pi/2.
\]
Let $H_N$ be the monic source Gram through degree $N$, $J_N$ the remainder
map into $\C[y]/Q$, and $G_N=(J_NH_N^{-1}J_N^*)^{-1}$ for $N\ge q-1$.
Its definition is the squared minimum norm in each original remainder fibre.
The monic source norms are
\begin{equation}\tag{AG3}
 \gamma_n=\sqrt{2\pi}\,n!(1/2)_n
 =\sqrt{2\pi}\frac{\Gamma(2n+1)}{4^n}.
\end{equation}
These are the original Meixner--Pollaczek specialization; the generating
function is $(1+t^2)^{-1/4}\exp(y\arctan t)$, with coefficients $p_n(y)/n!$.
The general polynomial convention is recorded in DLMF \cite{a03ag:DLMF}.
Put $R_n=\det[\int y^{i+j}|Q|^2\sigma\dd y]_{0\le i,j<n}$ and $R_0=1$.
The frame $(1,\ldots,y^{q-1},Q,yQ,\ldots,y^{N-q}Q)$ has determinant one.
Taking its Schur complement, with every source/relation cross term present,
proves $\det G_N=\det H_N/R_{N-q+1}$. Consequently
\begin{equation}\tag{AG4}
 \BB_k[\sigma]=\log R_q+\log R_{q+1}-\log R_1
       -\sum_{n=q}^{2q}c_{n-q}\log\gamma_n,
 \quad c_0=c_q=1,\quad c_r=2\ (0<r<q).
\end{equation}
The affine substitution $S=c+iy$ has triangular coefficient matrix with
diagonal $1,i,\ldots,i^N$, of determinant modulus one. Its inverse is
$y=(S-c)/i$; it intertwines the relation injections and remainder maps.
Thus (AG4) computes the original $S$-coordinate return, without a new mass.

\section{The fixed equilibrium measure and its inverse moment}
Use the EIQ measure for $V(x)=\pi\sqrt x-2\log x$, with support
$[u^2,v^2]$. In the original endpoint convention
\begin{equation}\tag{AG5}
\begin{gathered}
 \kappa^2=1-u^2/v^2,\qquad uK(\kappa)=2,\qquad vE(\kappa)=4,\\
 F=-\min\mathcal E_{2,\pi},\quad
 L=\int\log x\dd\rho(x),\quad M=\int x^{-1}\dd\rho(x).
\end{gathered}
\end{equation}
All integrals use the original probability density, rather than a scaled
endpoint model. Its complete proof is EIQ1--26 in
\texttt{EIQ\_GEL\_COMPLETE.tex}. For application of the analytic theorem
we recall its decisive exact formulas. For general $a=u^2,b=v^2$,
$R_s=\sqrt{(a+s)(b+s)}$, and parameters $A,B>0$, they are
\begin{equation}\tag{AG6}
 h_{A,B}(x)=\frac{B}{2\pi x}\int_0^\infty
 \frac{\sqrt s}{(x+s)R_s}\dd s>0,\qquad
 \rho_{A,B}(x)=\frac{\sqrt{(b-x)(x-a)}}{2\pi}h_{A,B}(x).
\end{equation}
The density has mass one, positive endpoint coefficients, and its effective
potential derivative off the support is
$\sqrt{(x-a)(x-b)}h_{A,B}(x)$, with the square root negative below $a$.
This proves the strict exterior gap on every closed interval outside the
support. The functions extend holomorphically near every compact positive
interval. These are precisely the offcriticality and exterior properties
used below. Variation under $x\mapsto tx$ at its minimum gives
$B\int\sqrt x\dd\rho_{A,B}=2(A+1)$; hence
\begin{equation}\tag{AG7}
 \int\sqrt x\dd\rho=6/\pi.
\end{equation}
Indeed the energy of the pushed-forward measure is
$B\sqrt t\int\sqrt x\dd\rho-A\log t-A\int\log x\dd\rho
-\log t-\iint\log|x-y|\dd\rho\dd\rho$; its derivative at $t=1$ is zero.

Here is an explicit derivation of the required new moment, including its sign.
Set $r_0=\sqrt{ab}$, $m=(a+b)/2$, and
$R(z)=\sqrt{(z-a)(z-b)}$ with $R(z)/z\to1$. EIQ's elementary transform is
\[
 H_s(z)=1-\frac{R_s+R(z)}{z+s},\qquad
 G(z)=\frac{A}{2r_0}H_0(z)-\frac B{4\pi}
       \int_0^\infty\frac{H_s(z)}{\sqrt s R_s}\dd s.
\]
Since $R(0)=-r_0$, $R'(0)=m/r_0$, and
$\int_0^\infty (\sqrt sR_s)^{-1}\dd s=2\pi A/(Br_0)$,
evaluation at zero, followed by an integration by parts, gives
\begin{align}\tag{AG8}
 \int x^{-1}\dd\rho_{A,B}
 &=\frac{A(a+b)}{4ab}
   -\frac B{4\pi}\int_0^\infty
           \frac{1-r_0/R_s}{s^{3/2}}\dd s\\
 &=\frac{A(a+b)}{4ab}
   -\frac{Br_0}{4\pi}\int_0^\infty
         \frac{a+b+2s}{\sqrt s[(a+s)(b+s)]^{3/2}}\dd s.\notag
\end{align}
The boundary terms vanish: $1-r_0/R_s=O(s)$ at zero and is bounded at infinity.
Differentiating the absolutely convergent identity
$\int (\sqrt sR_s)^{-1}\dd s=2K(\kappa)/v$ with the operator
$-2(\partial_a+\partial_b)$ evaluates the last integral as $2E(\kappa)/(u^2v)$.
This derivative also follows directly from the defining elliptic integrals,
and every differentiated integral is integrable at both ends. Therefore
\begin{equation}\tag{AG9}
 M=\frac{u^2+v^2-4uv}{2u^2v^2}>0.
\end{equation}
Its positivity follows from the original positive probability measure.

\section{Full-form bounds on both complementary regions}
Write $r_k=k\sqrt{\delta^2+\gamma^2}/q$, so $r_k\to0$.
Pair opposite roots before estimating. The normalized factor
$E_q(x)=|Q(q\sqrt x)|^2/q^{2q}$ satisfies for $x>r_k^2$
\begin{equation}\tag{AG10}
 (x-r_k^2)^q\le E_q(x)\le(x+r_k^2)^q;
 \quad E_q(x)\le(x+r_k^2)^q\quad(x\ge0).
\end{equation}
Stirling's formula and positivity on compact real intervals give fixed
$c_\sigma,C_\sigma'>0$ such that, for all $y\ge0$,
\[
 c_\sigma(1+y)^{-1/2}e^{-\alpha y}
 \le\sigma(y)\le C_\sigma'(1+y)^{-1/2}e^{-\alpha y}.
\]
For $\varepsilon=0,1$ use the exact rescaled density
$g_{q,\varepsilon}(x)=q^{1/2}x^{\varepsilon-1/2}E_q(x)\sigma(q\sqrt x)$.
Every relevant polynomial has degree $d\le q/2$. Its shifted Legendre
expansion on $[1,2]$ proves the elementary pointwise bounds
\begin{equation}\tag{AG11}
 |P(x)|^2\le(d+1)^2
 \begin{cases}7^{2d},&0\le x\le1,\\(5x)^{2d},&x\ge2\end{cases}
 \int_1^2|P(t)|^2\dd t.
\end{equation}
For completeness, the normalized basis is
$\sqrt{2j+1}P_j(2x-3)$; Cauchy--Schwarz and
$|P_j(t)|\le(2|t|+1)^j$ outside $[-1,1]$ give (AG11), since
$\sum_{j=0}^d(2j+1)=(d+1)^2$. The polynomial bound follows from the
integral representation of $P_j$ or its real three-term recurrence.
On $[1,2]$, with $r_k^2\le1/2$ and $q\ge1$, the full density is at least
$c_0e^{-C_0q}$, where one may take
$C_0=\log2+\alpha\sqrt2$ and
$c_0=c_\sigma/[\sqrt2\sqrt{1+\sqrt2}]$.

Choose fixed $0<\epsilon<\min(1/2,u^2/2)$ with
$14\epsilon e^{C_0}<e^{-2}$. Impose the finite guard $r_k^2\le\epsilon$.
On $[0,\epsilon]$, (AG10)--(AG11) bound the integral of
$|P|^2g_{q,\varepsilon}$, relative to its $[1,2]$ integral, by
\begin{equation}\tag{AG12}
 \frac{C_\sigma'}{c_0}(d+1)^2
 (14\epsilon e^{C_0})^q
 \frac{\epsilon^{\varepsilon+1/4}}{\varepsilon+1/4}.
\end{equation}
The bound uses $q^{1/2}\sigma(q\sqrt x)\le
C_\sigma'x^{-1/4}e^{-\alpha q\sqrt x}$, valid for every $x>0$.
In particular it includes the entire region $q\sqrt x\le1$.

Choose $M_0>\max(2,2v^2,(8/\alpha)^2)$ such that
$\log10+C_0+2\log x\le\alpha\sqrt x/2$ for $x\ge M_0$.
It suffices to check this at $M_0$ and the displayed derivative guard.
For $x\ge M_0$, the same formulas bound the upper-tail ratio by
\begin{equation}\tag{AG13}
 \frac{C_\sigma'}{c_0}(d+1)^2
 \int_{M_0}^\infty x^{\varepsilon-3/4}
                  e^{-\alpha q\sqrt x/2}\dd x.
\end{equation}
Here $(5x)^{2d}(2x)^q\le(10x^2)^q$ and $d\le q/2$.
Substitution $z=\sqrt x$ shows that (AG13) is bounded by a fixed
polynomial in $q$ times $e^{-\alpha q\sqrt{M_0}/4}$.
Thus the sum $b_q$ of (AG12)--(AG13) is $O_h(q^2e^{-c_1q})$ for a
fixed $c_1>0$. These are inequalities on every vector before minimization:
\begin{equation}\tag{AG14}
 H_{[\epsilon,M_0]}\preceq H_{(0,\infty)}
 \preceq(1+b_q)H_{[\epsilon,M_0]},\qquad
 0\le\log\frac{\det H_{(0,\infty)}}{\det H_{[\epsilon,M_0]}}
 \le n\log(1+b_q).
\end{equation}
The identical argument with $E_q=x^q$ gives the comparison for the power
reference. No polynomial root, minimum or complementary integral is omitted.

\section{Exact fixed-potential block and a bounded-parameter lemma}
On the fixed interval, the real part of Stirling's expansion gives
\begin{equation}\tag{AG15}
 \sigma(q\sqrt x)=\sqrt2\,q^{-1/2}x^{-1/4}
 e^{-\pi q\sqrt x/2}(1+O(q^{-2})),
\end{equation}
uniformly in $x$. Its inverse-$y$ contribution has zero real part;
the next term is $O(y^{-2})$. This relative bound passes through every
Gram in (AG14), costing $O(nq^{-2})$ in its logarithmic determinant.
The finite root sums give
\begin{equation}\tag{AG16}
 \sum\omega_{ab}^2=\frac{qk(k+2)}3(\gamma^2-\delta^2),\qquad
 \log\frac{E_q(x)}{x^q}=-\frac{\tau_k}{x}+r_q(x),
 \quad \tau_k=\frac{\eta_hk(k+2)}q,\quad
 \eta_h=\frac{\gamma^2-\delta^2}3.
\end{equation}
All sums include multiplicity. Indeed
$\sum_{a=0}^k(2a-k)=0$ and
$\sum_{a=0}^k(2a-k)^2=(k+1)k(k+2)/3$.
In a fixed complex neighborhood of $[\epsilon,M_0]$ the logarithmic series
converges once $r_k^2/|x|\le1/2$; opposite roots remove odd terms and the
sum of the fourth and higher terms is bounded by
$C_h k^4/q^3$. Hence $\|r_q\|_\infty\le C_hk^4/q^3$.
Replacing $e^{r_q}$ in a complete determinant costs at most
$n\|r_q\|_\infty=O_h(k^4/q^2)=O_h(1)$, since $q\ge(k+1)^2$.

Let $q=2N$, $n=N+s$, and use $(s,\varepsilon)=(0,0),(1,0),(0,1)$.
After its exact powers of $q$ and $\sqrt2$ are extracted, the block density is
\begin{equation}\tag{AG17}
 x^{q+\varepsilon-3/4}e^{-\pi q\sqrt x/2-\tau_k/x}
 =e^{-nV(x)+W_{s,\varepsilon,\tau_k}(x)},\qquad
 W_{s,\varepsilon,t}=(\varepsilon-2s-3/4)\log x+\pi s\sqrt x-t/x.
\end{equation}
This equality, not a limiting replacement, fixes every parity increment.
The potential $V$ is now independent of $n$ and $k$.

Borot--Guionnet's soft-edge expansion and analytic linear-statistic formula
\cite[Propositions 1.2 and 5.2, equations (1-3)--(1-6)]{a03ag:BG}, applied to
(AG6), give, for fixed analytic real $W$ on this interval,
\begin{equation}\tag{AG18}
 \log D_n[e^{-nV+W}]
 =n^2F+n\log(2\pi)+n\int W\dd\rho+O_W(\log n).
\end{equation}
Here $D_n$ is the monic Gram determinant. Its Vandermonde integral is the
partition function divided by $n!$. At ensemble index two the potential's
order-$n$ correction vanishes. The Gaussian reference is
$(2\pi)^{n/2}n^{-n^2/2}\prod_{j=0}^{n-1}j!$; Stirling summation gives
$-3n^2/4+n\log(2\pi)+O(\log n)$. This fixes the stated mass and factorial.
For the ratio with analytic factor the error after its order-$n$ term is
$O_W(1)$ by the cited linear-statistic formula. The support lies strictly
inside the actual integration interval, the potential is analytic there,
and (AG6) verifies its positive density and exterior gap; analyticity at
$x=0$ is not used.

We need no unproved uniformity in the varying $\tau_k$. Choose
$T>1+|\eta_h|$ and fix one of the three $W_{s,\varepsilon,0}$.
The function
\[
 f_n(t)=\log\frac{D_n[e^{-nV+W_{s,\varepsilon,t}}]}{D_n[e^{-nV}]}
                  -n\int W_{s,\varepsilon,t}\dd\rho
\]
is convex in $t$, by differentiating the positive squared-Vandermonde
integral twice (its second derivative is a variance).
Formula (AG18)'s ratio statement bounds $|f_n(-2T)|,|f_n(0)|,|f_n(2T)|$
by a common constant $C$. Convexity bounds $f_n(t)\le C$ on $[-T,T]$.
For $0\le t\le T$, write $0=\lambda t+(1-\lambda)(-2T)$,
$\lambda=2T/(t+2T)\ge2/3$; then $f_n(t)\ge-2C$.
The negative half is identical. This proves the uniform ratio error $O_h(1)$
for every actual $\tau_k$, since $|\tau_k|\le|\eta_h|$.
It includes fixed higher multiplicities without a fractional-power potential
expansion or a differentiated asymptotic remainder.

\section{Evaluate every block, source norm and root term}
The even block consists of $1,y^2,\ldots,y^{2n-2}$, and the odd block of
$y,y^3,\ldots,y^{2n-1}$. The substitution $y=q\sqrt x$, including both
halflines, contributes $q x^{-1/2}\dd x$.
Factoring the row and column powers and using (AG7), (AG14)--(AG18) proves
\begin{align}\tag{AG19}
 \log H_n^{q+\varepsilon}
 ={}&[2(q+\varepsilon)n+2n^2-\tfrac32n]\log q
       +\tfrac n2\log2+n^2F+n\log(2\pi)\\
 &+n[(\varepsilon-2s-\tfrac34)L+6s-\tau_kM]
       +O_h(\log q).\notag
\end{align}
The complete parity identity is
\begin{equation}\tag{AG20}
 \log R_q+\log R_{q+1}
 =\log H_N^q+\log H_{N+1}^q+2\log H_N^{q+1}.
\end{equation}
Their dimensions sum to $2q+1$. Their squared dimensions sum to
$q^2+q+1$; their $s$-weighted dimensions sum to $q/2+1$;
and their $\varepsilon$-weighted dimensions sum to $q$.
Keeping these exact sums until the end gives
\begin{align}\tag{AG21}
 \log R_q+\log R_{q+1}
 ={}&(6q^2+3q+\tfrac12)\log q+Fq^2\\
 &+q[F+3-\tfrac32L+3\log2+2\log\pi]
       -2\eta_h M k(k+2)+O_h(\log q).\notag
\end{align}
The extra $-\tau_kM$ from $2q+1$ is bounded and included in the displayed
error. Every constant dropped here is bounded independently of $k$.

For $R_1$ the reference integral is exactly
$2\sqrt2\,\Gamma(2q+1/2)/\alpha^{2q+1/2}$.
On $[\epsilon,M_0]$ the full root multiplier has logarithm bounded by
$|\tau_k|/\epsilon+\|r_q\|_\infty$, and (AG15) has relative error $O(q^{-2})$.
Both references satisfy (AG14), so their full integral ratio has bounded
logarithm. Stirling now gives
\begin{equation}\tag{AG22}
 \log R_1=2q\log q+(4\log2-2\log\pi-2)q+O_h(\log q).
\end{equation}
For the source norms, (AG3) yields uniformly for $q\le n\le2q$
$\log\gamma_n=2n\log n-2n+\tfrac12\log n+\tfrac32\log2+\log\pi+O(n^{-1})$.
The coefficients in (AG4) are twice the trapezoidal weights. Integrating
$2x\log x-2x$ and $\tfrac12\log x$ from $q$ to $2q$, and bounding the
trapezoidal remainder by its second derivative integral, gives
\begin{equation}\tag{AG23}
 \sum_{n=q}^{2q}c_{n-q}\log\gamma_n
 =6q^2\log q+(8\log2-9)q^2
  +q\log q+(5\log2+2\log\pi-1)q+O(\log q).
\end{equation}
For example $2\int_q^{2q}(2x\log x-2x)\dd x
=6q^2\log q+(8\log2-9)q^2$ and
$\int_q^{2q}\log x\dd x=q\log q+(2\log2-1)q$;
the constant part contributes $(3\log2+2\log\pi)q$.

\begin{theorem}[Absolute Gamma value on every fixed multiplicity]
With all original definitions retained, put
\begin{equation}\tag{AG24}
 C_B=9-8\log2+F,\qquad
 C_\sigma=F+6-\tfrac32L-6\log2+2\log\pi.
\end{equation}
Then
\begin{equation}\tag{AG25}
 \boxed{\BB_k[\sigma]=C_Bq^2+C_\sigma q
       -2\eta_hM k(k+2)+O_h(\log q).}
\end{equation}
\end{theorem}
\begin{proof}
Substitute (AG21)--(AG23) into the exact identity (AG4).
The $q\log q$ coefficient is $3-2-1=0$.
The remaining order-$q$ coefficient is
$F+3-\tfrac32L+3\log2+2\log\pi
-(4\log2-2\log\pi-2)-(5\log2+2\log\pi-1)$,
which is exactly $C_\sigma$. All preceding bounds are uniform in the
bounded $\tau_k$ and use $k^4/q^2\le1$; they apply to every fixed $m_0$.
\end{proof}

\section{Backwards and forwards through the original action}
Retain $w_h=|(2\xi/h)(1/2+iy)|^2/(2\pi)$, $m_k=w_h^{*k}$,
$I_h=\int_0^{\pi/2}[(\log M_h)'(\theta)]^2\dd\theta$,
$d=4m_0-2$, $c_\zeta=2\gamma_{\rm E}-\log(2\pi)$ and
$C_\Gamma=2L-8\log2+4$. The already proved SR1--48 receiver gives
\begin{equation}\tag{AG26}
 \BB_k[m_k]-\BB_k[\sigma]
 =-dC_\Gamma q+\frac{\log2}{\pi}I_hk^2
  +\frac{C_\Gamma q}{2(\log q+c_\zeta)}+R_{h,k},
 \qquad R_{h,k}=o_h(q/\log q).
\end{equation}
Its full finite remainder, annulus guards, inner and outer integrals and
the eight positive Grams are retained in
\texttt{SLOW\_TAIL\_AND\_NONLINEAR\_RECEIVER.tex}; its result is used
without differentiating $R_{h,k}$ or assigning its constituent signed pieces.
Adding (AG25) to that exact scalar difference proves
\begin{align}\tag{AG27}
 \BB_k[m_k]={}&C_Bq^2+(C_\sigma-dC_\Gamma)q
 +\frac{\log2}{\pi}I_hk^2-2\eta_hMk(k+2)\\
 &+\frac{C_\Gamma q}{2(\log q+c_\zeta)}
 +R_{h,k}+O_h(\log q).\notag
\end{align}
This evaluates the absolute baseline through the $q/\log q$ scale.
It strengthens the incoming $o(q)$ statement; no new arithmetic estimate
is presumed. The simple-packet root term contributes at order $q$.
At higher fixed multiplicity its displayed $k^2$ value is retained even
though the arithmetic error does not yet resolve all terms of that size.

The original CAI15 identity is
$J_k^{\rm action}=\BB_k[m_k]+\mathcal W_k^{\rm ar}-\mathcal P_--\mathcal P_+$.
Its finite guards and penalties give (CAI29--35)
\begin{gather}\tag{AG28}
 -E_{\rm hi}\le\mathcal W_k^{\rm ar}-\mathcal W_q^\sigma\le E_{\rm lo},
 \quad0\le\mathcal P_-+\mathcal P_+\le\Pi_{h,k},\\
 \mathcal W_q^\sigma=2\log\frac{(4q)!}{(2q)!}-4q\log2
       +\log\frac{2q-1}{2(4q-1)},\qquad
 E_{\rm hi}+E_{\rm lo}+\Pi_{h,k}=O_h(k+\log q).\notag
\end{gather}
In particular its exact factorial window is
$4q\log q+(8\log2-4)q-\log2+O(q^{-1})$.
Returning (AG27) through (AG28) proves
\begin{equation}\tag{AG29}
 \boxed{\begin{aligned}J_k^{\rm action}={}&C_Bq^2+4q\log q
  +[C_\sigma+8\log2-4-dC_\Gamma]q\\
 &+\frac{\log2}{\pi}I_hk^2-2\eta_hMk(k+2)
  +\frac{C_\Gamma q}{2(\log q+c_\zeta)}+o_h(q/\log q).
 \end{aligned}}
\end{equation}
Indeed $(k+\log q)/(q/\log q)\to0$ on the original sequence. At finite $k$
the exact interval (AG28) and the complete SR error remain the bounds,
rather than an assigned coefficient for any separate penalty or native loss.
The constant identity
\begin{equation}\tag{AG30}
 C_\sigma=C_B-\tfrac34 C_\Gamma+2\log(\pi/4)
\end{equation}
follows by substituting the original definitions; no mass is changed.
For the original benchmark $L_{h,k}=\delta q(k+1)/2$, subtracting its
actual logarithm gives the further receiver
\begin{align}\tag{AG31}
 J_k^{\rm action}-4q\log L_{h,k}
 ={}&C_Bq^2-4q\log(k+1)
 +[C_\sigma+8\log2-4-dC_\Gamma-4\log(\delta/2)]q\\
 &+\frac{\log2}{\pi}I_hk^2-2\eta_hMk(k+2)
 +\frac{C_\Gamma q}{2(\log q+c_\zeta)}+o_h(q/\log q).\notag
\end{align}
The full observation kernel still enters through its original inclusion
$I$ and quotient $\Lambda$: $H_{K,N}=I^*G_NI$ and
$Q_{B,N}=(\Lambda G_N^{-1}\Lambda^*)^{-1}$.
For any fixed section $S_*$ of $\Lambda$ the exact map
$[I,G_N^{-1}\Lambda^*Q_{B,N}]$ has determinant equal to $[I,S_*]$ and
orthogonal columns between its two blocks. Thus
\begin{equation}\tag{AG32}
 \det H_{K,N}\det Q_{B,N}=|\det[I,S_*]|^2\det G_N,
 \qquad F_K+F_B=\BB_k[m_k].
\end{equation}
Equation (AG27) replaces the previously unevaluated total in this equality;
it does not allocate that total to either original block. The independent
proper-source return and the original projected-current response ratios
retain their own exact maps and remain subsequent calculations.



% END PRESERVED INPUT 01_ABSOLUTE_GAMMA_AND_ACTION.tex
\endgroup

\clearpage
\begingroup
\def\C{\mathbb C}
\def\R{\mathbb R}
\def\PP{\mathcal P}
\def\norm#1{\left\lVert#1\right\rVert}
\def\Tr{\operatorname{Tr}}
\def\rank{\operatorname{rank}}
\def\im{\operatorname{im}}
\def\diag{\operatorname{diag}}
\def\spec{\operatorname{spec}}

\setcounter{equation}{0}
\renewcommand{\theequation}{A03PT.\arabic{equation}}
\section*{Full-primary positive-trace minimum and fixed-metric control}
\addcontentsline{toc}{section}{Full-primary positive-trace minimum and fixed-metric control}
% BEGIN PRESERVED INPUT 02_FULL_PRIMARY_POSITIVE_TRACE.tex SHA256 babad5ddc0b76ffb4fa4c27b9240c1c7539497adebdb6bf2f1ce34e2669b5adc



\begin{abstract}
For the unchanged arithmetic multiplication packet, the minimum of the
positive Hermitian weight trace over all positive metrics is evaluated,
attained, and its entire minimizing set is parametrized by finite
Lyapunov sums. All primary jets and all same-half cross terms are
retained. A rank-one observable forcing on each complete spectral half
constructs an attaining metric whose weight has precisely two nonzero
eigenvalues. The canonical attained polynomial metric is then evaluated
through its actual boundary rows: its exact rank-two scalar exceeds the
minimum by an explicit positive spectral-projection expression. The
positive-primary determinant line has exactly the same growth exponent
for every metric. Finally an original finite polynomial section realizes
every minimizing metric up to a specified scalar, and its complete
residual in each smaller polynomial source quotient is computed.
\end{abstract}

\section{Original packet, source and coordinate maps}

The incoming note \cite{a03pt:incoming} proposes the minimum and its complete
Lyapunov parametrization. The present text proves them, adds the
rank-two attaining construction, and evaluates their precise receivers.
The fixed source and cutoff conventions are GMF1--3 of
\cite{a03pt:gmf}; the earlier exact metric and source maps are TC1--29d of
\cite{a03pt:tc}. These source references record provenance. Every new finite
assertion below is proved here.

Keep the original hypothetical actual quartet with full multiplicity:
\begin{equation}\tag{PT1}
\begin{gathered}
 \rho=\tfrac12+\delta+i\gamma,\quad
 0<\delta<\tfrac12,\quad\gamma>2,\quad m_0\ge1,\\
 h(s)=\prod_{\zeta\in\{\rho,\bar\rho,1-\rho,1-\bar\rho\}}
                     (s-\zeta)^{m_0},\quad
 v_h=2\xi/h,\quad w_h(y)=|v_h(1/2+iy)|^2/(2\pi),\\
 m_k=w_h^{*k},\quad k\equiv1\pmod4,\quad k\ge5,\quad
 e=1+k(m_0-1),\quad q=e(k+1)^2,\quad c=k/2,\\
 \lambda_{ab}=c+(2a-k)\delta+i(2b-k)\gamma,\quad0\le a,b\le k,\\
 \chi(S)=\prod_{a,b=0}^k(S-\lambda_{ab})^e,\quad
 E=\C[S]/(\chi),\quad A=M_S:E\to E,\quad B=A-cI.
\end{gathered}
\end{equation}
The original basis of \(E\) is \(1,S,\ldots,S^{q-1}\).
The numbers \(\lambda_{ab}\) are pairwise distinct: equality first
forces the real indices \(a\), and then the imaginary indices \(b\),
to coincide. Thus \(\chi\) has exactly the displayed complete primary
factors. None is divided out in this paper.

For \(G>0\) in this original basis define
\begin{equation}\tag{PT2}
 W(G)=A^*G+GA-kG=B^*G+GB,\quad
 H_G=G^{-1/2}W(G)G^{-1/2},\quad
 T_+(G)=\Tr(H_G)_+.
\end{equation}
The symbol \(X_+\) means the positive part of a Hermitian matrix,
obtained by replacing each negative eigenvalue by zero; write
\(X_-=(-X)_+\). Positive square roots always refer to the fixed
standard coordinate Hermitian form. The physical coordinate in the
source remains \(S=c+iy\); its measure is exactly \(m_k(y)\,dy\).

Here is the complete confluent coordinate map and its inverse. Define
\begin{equation}\tag{PT3}
 C:E\longrightarrow\bigoplus_{a,b=0}^k\C^e,\qquad
 (CP)_{(\lambda,r)}=P^{(r)}(\lambda)/r!,\qquad
 C_{(\lambda,r),j}=\begin{cases}\binom jr\lambda^{j-r},&j\ge r,\\0,&j<r.
 \end{cases}
\end{equation}
For each \(\lambda\) put \(h_\lambda=\chi/(S-\lambda)^e\) and
write the actual Taylor series
\(h_\lambda(\lambda+z)^{-1}=\sum_{n\ge0}a_{\lambda,n}z^n\).
Its constant denominator \(h_\lambda(\lambda)\) is the full product
of the other root differences, each with exponent \(e\), and is
nonzero. Set
\begin{equation}\tag{PT4}
 p_{\lambda,r}(S)=h_\lambda(S)(S-\lambda)^r
       \sum_{n=0}^{e-r-1}a_{\lambda,n}(S-\lambda)^n,
 \qquad0\le r<e.
\end{equation}
This polynomial has degree at most \(q-1\). At \(\lambda\) its
Taylor class modulo \((S-\lambda)^e\) is \((S-\lambda)^r\);
at every other root its first \(e\) Taylor coefficients vanish.
Consequently these polynomials are the columns of \(C^{-1}\).
This proves both inverse laws of (PT3), without replacing any primary
factor by a simple root. Multiplication by \(S\) now gives
\begin{equation}\tag{PT5}
 CAC^{-1}=\bigoplus_\lambda(\lambda I_e+J_e),\qquad
 (J_e)_{r+1,r}=1\quad(0\le r<e-1),\quad J_e^e=0.
\end{equation}
The direct sum and the subsequent maps are coordinate morphisms;
the recorded operator remains the original \(A\).

Let \(E_+\) be the sum of the complete primary spaces with
\(a\ge(k+1)/2\), and \(E_-\) its complementary sum. Both have
dimension \(r=q/2\). Formulae (PT3)--(PT4) give explicit column
frames for their inclusions. For a column frame \(X_+:\C^r\to E_+\)
define the actual Euclidean orthogonal projection
\begin{equation}\tag{PT6}
 P_G=G^{1/2}X_+(X_+^*GX_+)^{-1}X_+^*G^{1/2}
       :\C^q\longrightarrow G^{1/2}E_+.
\end{equation}
It is self-adjoint and idempotent by matrix multiplication, and its
image is precisely the stated subspace. It is independent of the
choice of frame within \(E_+\).

\section{The trace bound, exact defect and equality}

The finite projection argument below is the positive-part form of
Ky Fan's variational trace principle \cite{a03pt:fan}. We give the full
argument rather than use it as a black box.

\begin{theorem}
For every original-coordinate \(G>0\),
\begin{equation}\tag{PT7}
 L:=\frac{\delta q(k+1)}2=\Tr(P_GH_G),\qquad
 T_+(G)-L
 =\Tr((I-P_G)(H_G)_+)+\Tr(P_G(H_G)_-)\ge0.
\end{equation}
The two terms in the defect are exactly
\begin{equation}\tag{PT8}
 T_+(G)-L
 =\norm{(H_G)_+^{1/2}(I-P_G)}_F^2
   +\norm{(H_G)_-^{1/2}P_G}_F^2.
\end{equation}
Equality holds precisely when \(E_+\) and \(E_-\) are
\(G\)-orthogonal and their actual restricted metrics satisfy
\begin{equation}\tag{PT9}
 B_+^*G_++G_+B_+\succeq0,\qquad
 C_-^*G_-+G_-C_-\succeq0,
 \quad B_+=A|_{E_+}-cI,\quad C_-=cI-A|_{E_-}.
\end{equation}
\end{theorem}
\begin{proof}
Put \(K=G^{1/2}BG^{-1/2}\), so \(H_G=K+K^*\).
The subspace \(G^{1/2}E_+\) is \(K\)-invariant, and the
compression of \(K\) to it is similar to \(B_+\). Therefore
\[
 \Tr(P_GH_G)=2\Re\Tr B_+
 =e(k+1)2\delta\sum_{a=(k+1)/2}^k(2a-k).
\]
The summands in the last sum are the odd integers
\(1,3,\ldots,k\); their sum is \(((k+1)/2)^2\). This proves
the first equality in (PT7). Substitute
\(H_G=(H_G)_+-(H_G)_-\) and use linearity of trace to obtain
the second equality. For a positive matrix \(X\) and an orthogonal
projection \(P\),
\(\Tr(PX)=\Tr(PX^{1/2}X^{1/2}P)=\norm{X^{1/2}P}_F^2\).
Apply this twice to get (PT8) and nonnegativity, without assuming
that either projection commutes with either positive part.

The defect vanishes exactly when
\((H_G)_+^{1/2}(I-P_G)=0\) and \((H_G)_-^{1/2}P_G=0\).
These equalities say that every positive eigenvector of \(H_G\)
belongs to \(\im P_G\), and every negative eigenvector belongs
to its orthogonal complement. The zero eigenspace may be split in
any way. In particular \(P_GH_G=H_GP_G\).
In an orthonormal frame adapted to \(\im P_G\), invariance gives
\[
 K=\begin{pmatrix}K_+&D\\0&K_-\end{pmatrix},\qquad
 H_G=\begin{pmatrix}K_++K_+^*&D\\D^*&K_-+K_-^*\end{pmatrix}.
\]
Commutation forces \(D=0\). Thus the orthogonal complement is
also \(K\)-invariant; its spectrum is precisely the negative-real
part of the original spectrum. To identify this complement exactly,
apply the complete spectral idempotents supplied by (PT4) with
\(r=0\). A vector in this invariant complement is killed by the
product of the negative primary factors, so all positive idempotent
components vanish. Dimension \(r\) then proves that the complement
is \(G^{1/2}E_-\). The restrictions of \(H_G\) on the two halves
are respectively nonnegative and nonpositive, which is exactly
(PT9) by congruence with their positive square roots.
Conversely, orthogonality and (PT9) give those two signs on the
two invariant halves. The positive trace is then their positive
half trace \(2\Re\Tr B_+=L\). This proves both directions,
including all possible zero eigenvalues within either half.
\end{proof}

\section{An explicit full-primary attainer}

Put \(t_0=\delta/2\), \(D_e=\diag(1,t_0,\ldots,t_0^{e-1})\),
and define in the original frame
\begin{equation}\tag{PT10}
 T=\left(\bigoplus_\lambda D_e\right)C,\qquad G_*=T^*T.
\end{equation}
The map \(T:(E,G_*)\to(\C^q,I)\) is an exact isometry.
Indeed \(U=TG_*^{-1/2}\) is unitary, and directly
\[
 UH_{G_*}U^*=TAT^{-1}+(TAT^{-1})^*-kI.
\]
Since \(D_eJ_eD_e^{-1}=t_0J_e\), each block on the right is
\(2(2a-k)\delta I_e+t_0(J_e+J_e^*)\). The vector with
\(r\)-th entry \(\sin((r+1)j\pi/(e+1))\), \(0\le r<e\),
is an eigenvector of \(J_e+J_e^*\) with eigenvalue
\(2\cos(j\pi/(e+1))\): the first and last missing entries are
zero, and the interior recurrence is the sine addition identity.
These \(e\) eigenvalues are distinct, giving a full eigenbasis.
Consequently
\begin{equation}\tag{PT11}
 \spec H_{G_*}
 =\left\{2(2a-k)\delta+\delta\cos\frac{j\pi}{e+1}:
 0\le a,b\le k,\ 1\le j\le e\right\}.
\end{equation}
Every one of these values has the sign of \(2a-k\), since
\(|2a-k|\ge1\) and \(|\cos|<1\). The cosine terms sum to
zero in each block by pairing \(j\) with \(e+1-j\). Thus
\begin{equation}\tag{PT12}
 \min_{G>0}T_+(G)=T_+(G_*)=L,\quad
 \rank(H_{G_*})_+=\rank(H_{G_*})_-=q/2,\quad
 \Tr H_{G_*}=0,\quad\Tr|H_{G_*}|=2L.
\end{equation}
In particular \(\norm{H_{G_*}}=2k\delta\) when \(e=1\),
and \(2k\delta+\delta\cos(\pi/(e+1))\) when \(e>1\).
These are norms of this displayed metric. They do not replace
the positive trace of any other original metric.

\section{Every minimizing metric and every nilpotent term}

The method here is the classical Lyapunov matrix equation method,
named for A.~M.~Lyapunov \cite{a03pt:lyapunov}. All convergence,
positivity and inverse statements needed here are proved below.
Use the fixed primary bases (PT3) on each half and put
\begin{equation}\tag{PT13}
 D_+=\bigoplus_{a>(k/2),b}( (\lambda_{ab}-c)I_e+J_e),\qquad
 D_-=\bigoplus_{a<(k/2),b}( (c-\lambda_{ab})I_e-J_e).
\end{equation}
Every eigenvalue of \(D_+\) and \(D_-\) has strictly positive
real part. Notice the minus sign on every negative-half nilpotent.

\begin{theorem}
All minimizing metrics, with no duplicates suppressed, are obtained
by choosing \(Q_+,Q_-\succeq0\) such that
\begin{equation}\tag{PT14}
 \bigcap_{j=0}^{r-1}\ker(Q_\pm^{1/2}D_\pm^j)=\{0\},
 \qquad r=q/2,
\end{equation}
and setting
\begin{equation}\tag{PT15}
 \Gamma_\pm=\int_0^\infty e^{-tD_\pm^*}Q_\pm e^{-tD_\pm}\,dt,
 \qquad G=C^*\diag(\Gamma_+,\Gamma_-)C.
\end{equation}
In this formula \(C\) is ordered first by the positive half and
then by the negative half, with the same primary coordinates
within each half. Every integral is a finite rational block sum.
For blocks \(D_\alpha=b_\alpha I+N_\alpha\) its exact value is
\begin{equation}\tag{PT16}
 \Gamma_{\alpha\beta}
 =\sum_{i,j=0}^{e-1}
 \frac{(-1)^{i+j}(i+j)!}
      {i!j!(\bar b_\alpha+b_\beta)^{i+j+1}}
 (N_\alpha^*)^i Q_{\alpha\beta}N_\beta^j.
\end{equation}
All blocks within each half, including off-diagonal blocks, occur.
\end{theorem}
\begin{proof}
The full block exponential is the finite sum
\(e^{-tb_\alpha}\sum_{j=0}^{e-1}(-t)^jN_\alpha^j/j!\).
Every entry of the integrand in (PT15), and its derivative, is
therefore a polynomial times an exponential of strictly negative
real exponent. The integrals converge absolutely and the endpoint
at infinity is zero. Differentiation and integration give
\[
 D_\pm^*\Gamma_\pm+\Gamma_\pm D_\pm=Q_\pm.
\]
The quadratic form of \(\Gamma_\pm\) on \(v\) is the integral
of \(\norm{Q_\pm^{1/2}e^{-tD_\pm}v}^2\). It vanishes
exactly when this continuous nonnegative integrand is identically
zero. Its entire power series is zero exactly when all
\(Q_\pm^{1/2}D_\pm^jv\) vanish. The characteristic polynomial
of \(D_\pm\) has degree \(r\); division by it expresses every
power \(j\ge r\) as a linear combination of lower powers.
This proves the finite test (PT14) and strict positivity.

Conversely let \(\Gamma>0\) on one half satisfy
\(Q=D^*\Gamma+\Gamma D\succeq0\). Differentiating
\(e^{-tD^*}\Gamma e^{-tD}\), and integrating between zero
and infinity, recovers (PT15) exactly. The preceding positivity
argument proves (PT14). The same argument shows that the
homogeneous Lyapunov equation has only the zero solution.
The characterization (PT9) now proves that (PT15) gives every
minimizer and only minimizers.

For (PT16), expand both finite exponentials. The term with powers
\(i,j\) has scalar integral
\(\int_0^\infty t^{i+j}e^{-zt}\,dt=(i+j)!/z^{i+j+1}\),
where \(z=\bar b_\alpha+b_\beta\) has positive real part.
The integral identity starts with \(1/z\) and follows by
successive integration by parts, whose endpoints vanish by the
same exponential decay. This proves every coefficient and sign.
On the negative half \(N_\alpha=-J_e\), so its two powers
supply an additional \((-1)^{i+j}\); these signs remain
explicit in (PT13)--(PT16). No inverse of a nilpotent is used.
\end{proof}

\section{A minimizing metric with rank-two weight}

There is an attaining family with many zero weight directions,
even when \(e>1\). On each half define the single literal row
\begin{equation}\tag{PT17}
 \ell_\pm(v)=\sum_{\lambda\text{ in that half}}v_{\lambda,e-1},
 \qquad Q_\pm=\ell_\pm^*\ell_\pm.
\end{equation}
The row extracts the final divided-derivative coefficient in every
complete primary block, with coefficient exactly one in each.

\begin{proposition}
The two choices (PT17) satisfy (PT14). The metric
\(G_{\mathrm{two}}\) obtained from (PT15)--(PT16) has
\begin{equation}\tag{PT18}
 \spec H_{G_{\mathrm{two}}}=\{L,-L,0,\ldots,0\},\qquad
 T_+(G_{\mathrm{two}})=\norm{H_{G_{\mathrm{two}}}}=L.
\end{equation}
Hence the positive-trace minimum is attained within the rank-two
weight stratum as well as within the full-rank weight stratum.
\end{proposition}
\begin{proof}
Write the blocks of either half as \(b_\lambda I+sJ_e\), with
the fixed common sign \(s=1\) on the positive half and \(s=-1\)
on the negative half. Then
\[
 \ell e^{-tD}v
 =\sum_\lambda e^{-tb_\lambda}
      \sum_{j=0}^{e-1}\frac{(-st)^j}{j!}v_{\lambda,e-1-j}.
\]
Distinct \(\lambda\)'s give distinct \(b_\lambda\)'s on
either half. The functions \(e^{-tb_\lambda}p_\lambda(t)\),
with \(\deg p_\lambda<e\), are linearly independent. Here is
an explicit proof: apply
\(\prod_{\mu\ne\lambda}(\partial_t+b_\mu)^e\) to a
vanishing sum. It kills every summand except the \(\lambda\)
summand, where it gives \(e^{-tb_\lambda}\) times
\(\prod_{\mu\ne\lambda}(\partial_t+b_\mu-b_\lambda)^e
p_\lambda\). On the polynomials of degree less than \(e\),
each operator \(\partial_t+a\) with \(a\ne0\) is invertible:
its matrix is triangular with diagonal \(a\), or its inverse
is the finite geometric sum in \(-\partial_t/a\).
Therefore every \(p_\lambda\) is zero. Their displayed
coefficients force every coordinate of \(v\) to be zero.
Thus \(\ell e^{-tD}v=0\) for all \(t\) implies \(v=0\),
and (PT14) follows as in the preceding theorem.

In each half the Lyapunov weight is \(Q_\pm\), of rank one.
Whitening by the positive \(\Gamma_\pm\) preserves this rank
and gives a nonnegative rank-one matrix on the positive half
and a nonpositive rank-one matrix on the negative half. Its
nonzero eigenvalue is its trace. Cyclicity of trace gives
\[
 \Tr(\Gamma_+^{-1}Q_+)=2\Re\Tr D_+=L,
 \qquad\Tr(\Gamma_-^{-1}Q_-)=2\Re\Tr D_-=L.
\]
The two complete halves are orthogonal in the constructed metric,
so these are exactly the eigenvalues in (PT18). All vectors of
all primary blocks remain present; zero Hermitian weight does not
remove those vectors or their original operator action.
\end{proof}

The rank-two metric has fully evaluated entries as well. On either
half put \(\sigma=1\) for \(D_+\), \(\sigma=-1\) for \(D_-\),
and, for indices \(0\le i,j<e\), put \(u=e-1-i\), \(v=e-1-j\).
In (PT16), \(Q_{\alpha\beta}\) is the matrix with its only nonzero
entry, equal to one, at \((e-1,e-1)\). Thus
\begin{equation}\tag{PT18a}
 (\Gamma_\pm)_{\alpha i,\beta j}
 =\frac{(-\sigma)^{u+v}(u+v)!}
        {u!v!(\bar b_\alpha+b_\beta)^{u+v+1}}.
\end{equation}
For \(u+v>0\), substituting these entries into the Lyapunov
equation reduces its numerator to
\[
 \frac{(u+v)!}{u!v!}
 -\mathbf1_{u>0}\frac{(u+v-1)!}{(u-1)!v!}
 -\mathbf1_{v>0}\frac{(u+v-1)!}{u!(v-1)!}=0.
\]
For \(u=v=0\) the single term gives one. This directly verifies
every row and both nilpotent signs independently of the integral.

Every positive/negative weight-rank pair
\(1\le r_+,r_-\le q/2\) occurs among minimizers. In each half
choose a positive semidefinite matrix of the prescribed rank whose
range contains \(\ell_\pm^*\). Its kernel is contained in
\(\ker\ell_\pm\), so its finite observability kernel is zero
by the proof of (PT17). Formula (PT15) gives a positive metric,
and invertible congruence preserves the prescribed weight rank.
This proves the claim, including all intermediate zero-weight
dimensions, with every original primary vector present.

\section{The original positive-primary exterior line}

Equip every exterior power with the Gram-determinant metric induced
by its stated \(G\). The inclusion of the positive-primary line is
the actual map
\(\Lambda^rX_+:\Lambda^r\C^r\to\Lambda^r E\), not a
choice of a positive eigenspace of \(H_G\). On this line the
generator induced by \(B\) is multiplication by \(\Tr B_+\).
Its imaginary part vanishes because for each fixed \(a\) the
sum of \(2b-k\) over \(0\le b\le k\) is zero. Its real
part is \(L/2\) by (PT7). Thus for any nonzero element
\(\omega_+\in\Lambda^r E_+\) and every real \(t\),
\begin{equation}\tag{PT19}
 \norm{\Lambda^r e^{tB}\omega_+}_{\Lambda^rG}^{2}
       =e^{Lt}\norm{\omega_+}_{\Lambda^rG}^{2},\qquad
 \norm{\Lambda^r e^{tA}\omega_+}_{\Lambda^rG}^{2}
       =e^{(kr+L)t}\norm{\omega_+}_{\Lambda^rG}^{2}.
\end{equation}
To verify the claimed action directly, triangularize each full
primary block by its fixed basis (PT5). The determinant of
\(e^{tB_+}\) is the product of its diagonal entries,
\(e^{t\Tr B_+}\); all nilpotent exponential diagonals are one.
Squaring the scalar modulus gives (PT19). In the original column
frame the derivative has the equally exact expression
\begin{equation}\tag{PT20}
 \left.\frac{d}{dt}\right|_{t=0}
 \log\det\bigl(X_+^*e^{tB^*}Ge^{tB}X_+\bigr)
 =\Tr\bigl((X_+^*GX_+)^{-1}X_+^*W(G)X_+\bigr)
 =\Tr(P_GH_G)=L.
\end{equation}
The derivative of \(\log\det M(t)\) follows by multilinear
expansion of \(\det(I+tM(0)^{-1}M'(0))\). Substituting (PT6)
and moving finite factors around the trace proves the second
equality. This supplies the exact receiver to the original exterior
benchmark for every \(G\), including the fixed attained metric.

\section{The canonical attained metric and its exact rank-two gap}

Use the original source \(m_k(y)\,dy\) and physical coordinate
\(S=c+iy\) from (PT1). Its positivity and moments are the original
source data of \cite{a03pt:gmf}; in particular it is positive on a set
of positive measure and all polynomial moments are finite. For
\(\PP_N=\C[S]_{\le N}\), \(N\ge q-1\), put
\begin{equation}\tag{PT21}
 \begin{gathered}
 \langle f,g\rangle=\int_\R\overline{f(c+iy)}g(c+iy)m_k(y)\,dy,
 \quad H_N=(\langle S^i,S^j\rangle)_{0\le i,j\le N},\\
 J_N:\PP_N\twoheadrightarrow E,\quad f\mapsto[f]_\chi,\qquad
 G_N=(J_NH_N^{-1}J_N^*)^{-1},\quad
 R_N=H_N^{-1}J_N^*G_N.
 \end{gathered}
\end{equation}
A nonzero polynomial has only finitely many zeros on the vertical
line; thus \(H_N>0\). Surjectivity of \(J_N\) gives \(G_N>0\).
Matrix multiplication gives \(J_NR_N=I\) and
\(R_N^*H_NR_N=G_N\). For \(z\in\ker J_N\),
\(z^*H_NR_N=z^*J_N^*G_N=0\). Therefore \(R_Nu\) is the
actual attained minimum representative of \(u\); every other
representative has squared norm \(u^*G_Nu+\norm z^2\).

Define the following original rows and polynomial, with no omitted
leading coefficient:
\begin{equation}\tag{PT22}
 \ell_N(u)=[S^N]R_Nu,\qquad
 r_N(S)=\chi(S)S^{N+1-q},\qquad
 s_N(u)=\langle r_N,R_Nu\rangle.
\end{equation}
The exponent \(N+1-q\) is nonnegative even at \(N=q-1\).
Set \(T_Nf=Sf-[S^N]f\,r_N\). As \(\chi\) is monic,
the degree-\(N+1\) coefficient cancels, so \(T_N:\PP_N\to\PP_N\).
It satisfies the exact quotient intertwining relation
\(J_NT_N=AJ_N\). Thus \(R_NA\) is the orthogonal projection
of \(T_NR_N\) onto \((\ker J_N)^\perp\).
The physical identity \(\bar S+S=k\) on the vertical line gives
\(\langle Sf,g\rangle+\langle f,Sg\rangle=k\langle f,g\rangle\).
Taking \(f=R_Nu,g=R_Nv\), and subtracting the two displayed
leading-coefficient corrections, proves
\begin{equation}\tag{PT23}
 W(G_N)=-(\ell_N^*s_N+s_N^*\ell_N),\qquad
 \rank W(G_N)\le2,\qquad
 \Tr H_{G_N}=2\Re\Tr A-kq=0.
\end{equation}
The last equality follows by summing \(2a-k\) over all \(a,b\).
Since (PT7) gives a strictly positive trace lower bound, the rank
is exactly two and its two nonzero eigenvalues are
\(+\epsilon_N,-\epsilon_N\), with \(\epsilon_N>0\).

Define, in the original packet coordinates,
\begin{equation}\tag{PT24}
 u_N=s_NG_N^{-1}s_N^*,\qquad
 v_N=\ell_NG_N^{-1}\ell_N^*,\qquad
 z_N=s_NG_N^{-1}\ell_N^*.
\end{equation}
The trace identity in (PT23) gives \(\Re z_N=0\). Put
\(a=G_N^{-1/2}\ell_N^*\) and \(b=G_N^{-1/2}s_N^*\),
so \(H_{G_N}=-(ab^*+ba^*)\). Direct multiplication yields
\[
 \Tr(H_{G_N}^2)=2u_Nv_N+2\Re(z_N^2)
                 =2u_Nv_N-2(\Im z_N)^2.
\]
The eigenvalues in (PT23) give the exact certificate and bound
\begin{equation}\tag{PT25}
 \epsilon_N^2=u_Nv_N-(\Im z_N)^2
             =u_Nv_N-|z_N|^2\ge L^2.
\end{equation}
No coefficient of \(R_N\), no root difference in \(\chi\),
and no term in either original source region was changed to obtain
this identity.

Choose orthonormal eigenvectors \(p_N,n_N\) of the whitened
canonical matrix with eigenvalues \(\epsilon_N,-\epsilon_N\).
Then the full exact receiver of (PT7) is
\begin{equation}\tag{PT26}
 \epsilon_N-L
 =\epsilon_N\left(1-\norm{P_{G_N}p_N}^2
                         +\norm{P_{G_N}n_N}^2\right).
\end{equation}
This follows by substituting
\((H_{G_N})_+=\epsilon_Np_Np_N^*\) and
\((H_{G_N})_-=\epsilon_Nn_Nn_N^*\) in (PT7).
In particular equality means precisely that \(p_N\) lies in
\(G_N^{1/2}E_+\) and \(n_N\) lies in its orthogonal complement;
equivalently (PT9) holds for the fixed \(G_N\).
Equations (PT22)--(PT26) evaluate the comparison through its actual
minimum map. They do not assert that this minimum sequence happens
to lie in the minimizing set (PT15).

\section{Exact metric realization in the original finite source}

Use \(N=2q-1\), the unmodified \(H_N,J_N,R_N,G_N\) in
(PT21), and the complete relation map
\begin{equation}\tag{PT27}
 K_\chi:\PP_{q-1}\hookrightarrow\PP_{2q-1},\quad
 f\mapsto\chi f,\qquad
 Q_\chi=K_\chi^*H_NK_\chi>0,\quad
 B_\chi=K_\chi Q_\chi^{-1/2}.
\end{equation}
Injectivity follows from the integral domain \(\C[S]\).
Polynomial division gives \(\im K_\chi=\ker J_N\), exactly,
and this space has dimension \(q\). Hence
\(J_NB_\chi=0\), \(B_\chi^*H_NB_\chi=I\), and
\(B_\chi^*H_NR_N=0\).

For every displayed minimizing metric \(G_0\) from (PT15),
including \(G_*\) and \(G_{\mathrm{two}}\), set
\begin{equation}\tag{PT28}
 \begin{gathered}
 \alpha=1+\norm{G_0^{-1/2}G_NG_0^{-1/2}},\qquad
 \Delta=\alpha G_0-G_N,\qquad
 F=Q_\chi^{-1/2}\Delta^{1/2}:E\longrightarrow\PP_{q-1},\\
 R=R_N+K_\chi F.
 \end{gathered}
\end{equation}
The norm is the ordinary largest eigenvalue of the positive matrix
shown. Consequently
\(G_0^{-1/2}\Delta G_0^{-1/2}\succeq I\), so \(\Delta>0\)
and \(F\) is invertible. The orthogonal cross terms vanish and give
\begin{equation}\tag{PT29}
 J_NR=I_E,\qquad R^*H_NR=G_N+F^*Q_\chi F
                         =\alpha G_0,\qquad
 H_{\alpha G_0}=H_{G_0}.
\end{equation}
The last equality follows directly from (PT2), because a positive
scalar commutes with the specified positive square root. This is
an actual finite polynomial section realizing the trace minimum.
It retains the full relation \(\chi\), not a root-value ideal.

Its residual at every smaller polynomial source is also exact.
For \(W\subseteq\PP_{q-1}\), write
\[
 Q_W=\PP_{2q-1}/K_\chi W,\qquad
 \eta_W:\PP_{q-1}/W\hookrightarrow Q_W,\quad
 [f]\mapsto[K_\chi f],\qquad
 \pi_W:Q_W\twoheadrightarrow E.
\]
The first map is well-defined and injective because \(K_\chi\)
is injective, and the second is induced by \(J_N\).
Their image and kernel agree by \(\ker J_N=\im K_\chi\).
Therefore this is the explicit exact sequence and section comparison
\begin{equation}\tag{PT30}
 \begin{gathered}
 0\longrightarrow\PP_{q-1}/W\xrightarrow{\eta_W}Q_W
   \xrightarrow{\pi_W}E\longrightarrow0,\\
 [R(\cdot)]_W-[R_N(\cdot)]_W=\eta_W([F(\cdot)]_W),\qquad
 \ker([F]_W)=F^{-1}W,\qquad
 \rank([F]_W)=q-\dim W.
 \end{gathered}
\end{equation}
Surjectivity of \([F]_W\) follows from invertibility of \(F\).
Thus every proper finite polynomial source retains a nonzero
residual map of the computed rank. The arithmetic quotient class
is nevertheless identical, because \(\pi_W[R]_W=I_E\).

For completeness, the entire change in the multiplication boundary
is a relation as well, with its original degree recorded. Interpret
the identity below in \(\PP_{2q}\), and define
\(\widehat F:E\to\PP_q\) by \(\widehat F=SF-FA\).
Then direct substitution gives
\begin{equation}\tag{PT31}
 (SR-RA)-(SR_N-R_NA)=K_\chi^{(q)}\widehat F,
 \quad K_\chi^{(q)}:\PP_q\hookrightarrow\PP_{2q},\ f\mapsto\chi f.
\end{equation}
For a target source \(W'\subseteq\PP_q\) this difference has
precisely the class \([\widehat F]_{W'}\) under the injective
map induced by \(K_\chi^{(q)}\). Thus a full-quotient equality
has not erased the smaller-source generator residual.

\section{Backward receivers for the earlier theta source}

The original TC source has weight one and a packet of actual zeros;
the present cyclic sum packet has weight \(k\). We give the exact
receiving operations at those two types. For any positive Gram and
operator pair \((G,A_0)\) of weight \(w\), the operation is
\((G,A_0,w)\mapsto B_0=A_0-wI/2\), followed by the trace
projection proof (PT7)--(PT9) on its complete positive-primary sum.
The polynomial coordinate change to (PT1), when used, is the
original full-primary sum realization, whose explicit CRT map is
(PT3). This operation does not replace an ordered tensor space by
its cyclic image: TC21 remains on its original ordered tensor space.

In the original TC22 coordinates, write
\(W_m=-z_mb_m-b_m^*z_m^*\), where \(z_m\) is its actual
column and \(b_m\) its actual row. The receiver to our row form
is exactly
\begin{equation}\tag{PT32}
 \ell=b_m,qquad s=z_m^*,\qquad
 sG_m^{-1}s^*=u_m,quad \ell G_m^{-1}\ell^*=v_m,quad
 sG_m^{-1}\ell^*=\bar c_m.
\end{equation}
Complex conjugation changes the sign of the imaginary part but
not its square; (PT25) is therefore precisely the original scalar
certificate TC26 and lower bound TC28 in these coordinates. The
new projection formula (PT26) strengthens its description by
evaluating the excess relative to the unchanged positive-primary
exterior line. The known rank-one Gram update TC27 remains exactly
as stated; it is not assigned to the distinct sequence (PT21).
TC18 and TC21 are operator-norm infima. Equations (PT12) and
(PT18) add the positive-trace minimum and a minimizing rank-two
stratum, without changing those previously proved infima.

Here is the exact smaller-source map for the original theta
realization TC29a--d \cite{a03pt:tc}, including its finite correction.
Retain
\[
 V=\{\phi\in\mathcal S(\R):\phi\text{ even},\ \phi(0)=0,
                 \ \int_\R\phi=0\},\quad D=-x\partial_x,
 \quad\Theta\phi=2\sum_{n\ge1}\phi(nx),
\]
and the original space \(\mathscr B\) of smooth functions rapidly
decreasing at both ends on \(\R_{>0}\). The map \(\Theta:V\to
\mathscr B\) is injective. Indeed on \(\Re s>1\) absolute
convergence and the substitution \(u=nx\) give
\(\mathcal M\Theta\phi(s)=2\zeta(s)\mathcal M\phi(s)\).
The Euler product is nonzero there, and uniqueness of the Mellin
transform on a vertical line is Fourier uniqueness after \(x=e^y\).
Thus \(\Theta\phi=0\) forces \(\phi=0\).
For each \(W\subseteq V\) define the actual maps
\[
 \eta_W:V/W\longrightarrow\mathscr B/\Theta W,quad
 [\phi]\longmapsto[\Theta\phi],\qquad
 \pi_W:\mathscr B/\Theta W\longrightarrow\mathscr B/\Theta V.
\]
Injectivity of \(\Theta\) proves injectivity of \(\eta_W\),
and its image is exactly \(\ker\pi_W\). For the original
reference section \(s_h\) and its finite primitive correction
\(T=\Phi_LC:E_h\to V\), put \(R=s_h+\Theta T\). Subtraction
of actual functions gives
\begin{equation}\tag{PT33}
 [R]_W-[s_h]_W=\eta_W([T]_W),\qquad
 \pi_W[R]_W=[s_h]_V,qquad
 \ker([R-s_h]_W)=T^{-1}W.
\end{equation}
The last equality uses the proved injectivity. This is the
precise map from a full-source correction to its smaller-source
class, for the original theta packet \(E_h\).

The original reference relation is
\(Ds_h-s_hA_h=\Theta(\phi_*\ell_h)\). Since termwise
differentiation gives \(D\Theta=\Theta D\), substitution proves
\begin{equation}\tag{PT34}
 DR-RA_h=\Theta K_R,\qquad
 K_R=\phi_*\ell_h+DT-TA_h:E_h\longrightarrow V.
\end{equation}
For \(DW\subseteq W'\), the induced map
\(\bar D:\mathscr B/\Theta W\to\mathscr B/\Theta W'\)
is well-defined because \(D\Theta W\subseteq\Theta W'\).
Its defect on (PT33) is exactly \(\eta_{W'}([K_R]_{W'})\),
which vanishes on a vector precisely when that original primitive
belongs to \(W'\). This proves the typed generator receiver.

For \(a>0\) let \(\mathcal R_aF(x)=F(x/a)\) and let
\(A_{a,h}=e^{(\log a)A_h}\) be the original arithmetic action.
Retain the reference primitive \(\Psi_{a,h}\) defined by
\(\mathcal R_as_h-s_hA_{a,h}=\Theta\Psi_{a,h}\).
The equality \(\mathcal R_a\Theta=\Theta\mathcal R_a\)
and substitution give
\begin{equation}\tag{PT35}
 \begin{gathered}
 c_{R,a}=\Psi_{a,h}+\mathcal R_aT-TA_{a,h},\qquad
 \mathcal R_aR-RA_{a,h}=\Theta c_{R,a},\\
 \bar{\mathcal R}_a^{W\to\mathcal R_aW}[R]_W
       -[R]_{\mathcal R_aW}A_{a,h}
       =\eta_{\mathcal R_aW}([c_{R,a}]_{\mathcal R_aW}),\\
 c_{R,ab}=\mathcal R_a c_{R,b}+c_{R,a}A_{b,h}.
 \end{gathered}
\end{equation}
For the last identity, expand the preceding two definitions and
cancel the terms \(\mathcal R_aT A_{b,h}\). The reference
primitives satisfy the same identity because the difference has
zero \(\Theta\)-image and \(\Theta\) is injective. The target
source remains \(\mathcal R_aW\), so no source invariance is
silently inserted. Equations (PT30)--(PT31) give the complete
polynomial counterpart on the unchanged sum packet, and
(PT33)--(PT35) are the actual TC proper-source receiver.

\section{Quantitative excess and the original cross metric}

The trace defect also measures the original metric interaction
between the positive and negative primary halves. In the fixed
CRT ordering let
\begin{equation}\tag{PT36}
 \begin{gathered}
 \mathcal G=(C^{-1})^*GC^{-1}
       =\begin{pmatrix}P&X\\X^*&Y\end{pmatrix},\quad
 K=P^{-1}X,\quad Z=Y-X^*P^{-1}X>0,\\
 F=\begin{pmatrix}P^{-1/2}&-KZ^{-1/2}\\0&Z^{-1/2}\end{pmatrix},
 \qquad\mathcal F=C^{-1}F,\qquad\mathcal F^*G\mathcal F=I.
 \end{gathered}
\end{equation}
Positivity of \(P\) is restriction positivity; positivity of
\(Z\) follows by testing \(\mathcal G\) on \((-Kv,v)\).
Multiplication verifies the displayed isometry. Write the original
CRT operator as \(CBC^{-1}=\diag(B_+,B_-)\), with
\(B_-=-D_-\). Its matrix in this isometry is
\begin{equation}\tag{PT37}
 \begin{gathered}
 \mathcal F^{-1}B\mathcal F
   =\begin{pmatrix}\widehat B_+&D\\0&\widehat B_-\end{pmatrix},
 \quad\widehat B_+=P^{1/2}B_+P^{-1/2},\quad
 \widehat B_-=Z^{1/2}B_-Z^{-1/2},\\
 D=P^{1/2}(KB_--B_+K)Z^{-1/2},\qquad
 \mathcal H=\mathcal F^*W(G)\mathcal F
   =\begin{pmatrix}D_+^{\rm H}&D\\D^*&D_-^{\rm H}\end{pmatrix},\\
 D_\pm^{\rm H}=\widehat B_\pm+\widehat B_\pm^*,\quad
 \Tr D_+^{\rm H}=L,\quad\Tr D_-^{\rm H}=-L.
 \end{gathered}
\end{equation}
The matrix \(G^{1/2}\mathcal F\) is unitary, so \(\mathcal H\)
and \(H_G\) have exactly the same eigenvalues. Let
\(\mathcal E=T_+(G)-L\), and let \(\norm D_1\) denote the
sum of the singular values of \(D\). Then
\begin{equation}\tag{PT38}
 \norm D_1^2\le\mathcal E(2L+\mathcal E),\qquad
 \mathcal E\ge\sqrt{L^2+\norm D_1^2}-L.
\end{equation}
To prove this, take a singular-value decomposition
\(D=U\Sigma V^*\), extending the singular frames to full unitary
bases, and put \(W=UV^*\). The matrices
\[
 J_0=\begin{pmatrix}I&0\\0&-I\end{pmatrix},\qquad
 T_0=\begin{pmatrix}0&W\\W^*&0\end{pmatrix}
\]
square to \(I\) and anticommute. Hence
\(Z_\theta=\cos\theta J_0+\sin\theta T_0\) is a Hermitian
contraction for every real \(\theta\). Diagonalizing
\(\mathcal H\), every diagonal entry of \(Z_\theta\) lies
between \(-1\) and \(1\), so
\(\Tr(Z_\theta\mathcal H)\le\Tr|\mathcal H|\).
The two sides are respectively
\(2L\cos\theta+2\norm D_1\sin\theta\) and
\(2(L+\mathcal E)\), since \(\Tr\mathcal H=0\).
Maximizing over \(\theta\) proves (PT38).

There is a complete finite inverse from this measured cross block
to the original \(X\). For positive-half block parameters
\(b_\alpha=\lambda_\alpha-c\) and negative-half parameters
\(\beta_\eta=\lambda_\eta-c\), the original operator blocks
are \(b_\alpha I+J_e\) and \(\beta_\eta I+J_e\).
Set \(R=P^{-1/2}DZ^{1/2}\). Then
\begin{equation}\tag{PT39}
 K_{\alpha\eta}
 =-\sum_{u,v=0}^{e-1}
 \frac{(-1)^u(u+v)!}
      {u!v!(b_\alpha-\beta_\eta)^{u+v+1}}
 J_e^uR_{\alpha\eta}J_e^v,\qquad X=PK.
\end{equation}
Indeed \(R=KB_--B_+K\) by (PT37). The absolutely convergent
integral \(-\int_0^\infty e^{-tB_+}Re^{tB_-}\,dt\) solves
this equation by integration of its derivative. Each real
exponent is at most \(-2\delta t\), multiplied by a finite
polynomial. Expansion of its two full nilpotent exponentials
gives (PT39), with sign \((-1)^u\) because the right exponential
is \(e^{+tB_-}\). Uniqueness follows from the homogeneous
equation: \(e^{-tB_+}Ke^{tB_-}=K\) then holds for all \(t\),
while its left side tends to zero. Thus the integral and the
finite sum invert the original cross map exactly.

For an explicit bound in the same full primary coordinates set
\begin{equation}\tag{PT40}
 \begin{gathered}
 \kappa=\max_{\alpha,\eta}
       \sum_{u,v=0}^{e-1}
       \frac{(u+v)!}{u!v!|b_\alpha-\beta_\eta|^{u+v+1}},\\
 \norm{P^{-1}X}_F^2
 \le\kappa^2\norm{P^{-1/2}}^2\norm{Z^{1/2}}^2
                  \mathcal E(2L+\mathcal E).
 \end{gathered}
\end{equation}
Each power of \(J_e\) has norm at most one. Applying the triangle
inequality to every block in (PT39), then summing squared block
Frobenius norms, gives \(\norm K_F\le\kappa\norm R_F\).
The multiplication bound for \(R=P^{-1/2}DZ^{1/2}\), followed
by \(\norm D_F\le\norm D_1\) and (PT38), proves (PT40).
All original metric factors are displayed; no uniform bound on
their condition numbers has been inserted.
For the canonical metric (PT21) one substitutes exactly
\(G=G_N\) and \(\mathcal E=\epsilon_N-L\). Thus the same
scalar which occurs in (PT25)--(PT26) controls the actual
positive/negative primary metric cross block through (PT36)--(PT40).
In particular the original scalar certificate itself strengthens to
\begin{equation}\tag{PT41}
 u_Nv_N-(\Im z_N)^2\ \ge\ L^2+\norm{D_N}_1^2,
 \qquad
 D_N=P_N^{1/2}(K_NB_--B_+K_N)Z_N^{-1/2},
\end{equation}
where \(P_N,X_N,Y_N,K_N,Z_N\) are obtained by substituting
the unchanged \(G_N\) in (PT36). Indeed
\(\mathcal E(2L+\mathcal E)=(\epsilon_N-L)
(\epsilon_N+L)=\epsilon_N^2-L^2\), so (PT38) and (PT25)
give precisely (PT41). This is an estimate on the actual attained
metric; every factor in its cross block has its original coordinate
map and its finite inverse in (PT36)--(PT39).

\section{Scope of the finite result}

The positive-trace minimum (PT12), every minimizing metric
(PT14)--(PT16), the rank-two attainer (PT18), the exterior growth
(PT19)--(PT20), the canonical scalar and defect (PT25)--(PT26),
the finite source and proper-source residuals (PT28)--(PT35),
and the complete cross-metric receiver (PT36)--(PT41)
are established by the displayed proofs. The fixed attained
metrics remain the matrices in (PT21). The variational minimum
does not evaluate their separate projected kernel and boundary
allocations; their exact trace gap is (PT26). No value for an
unevaluated projection is assigned in this result.



% END PRESERVED INPUT 02_FULL_PRIMARY_POSITIVE_TRACE.tex
\endgroup

\clearpage
\begingroup


\setcounter{equation}{0}
\renewcommand{\theequation}{A03CM.\arabic{equation}}
\section*{Every original coefficient section and native--arithmetic metric return}
\addcontentsline{toc}{section}{Every original coefficient section and native--arithmetic metric return}
% BEGIN PRESERVED INPUT 03_COEFFICIENT_METRIC_COMPLETE.tex SHA256 fb3e18774eaf5f8686e6a69a52c12f7f4a1f32b47f3864edcb30b58046fa60d6
% Exact original-chart calculation; no new arithmetic injection is substituted.
\section{The prescribed conductor minor and the exact arithmetic image}

This derives the coefficient-image continuation supplied in
\cite{a03cm:program:arrival03-01}, retaining its exact residue map.
All parameters and full polynomial factors are those of CAL1.
Let $\mathbb V=(\beta_\alpha^n)_{0\le n<q,\alpha}$ be the
original upper Vandermonde, with the original root labels, and let
$P_v$ insert zeros in its first $v$ derivative coordinates.
Write $q'=(a-7)^2$ and retain the original conductor
\[
 E_A(z)=\sum_{r,t=0}^8a_{rt}e^{\beta_{8,r,t}z},\quad
 \mu_j=E_A^{(j)}(0),\quad
 \mu_0=\cdots=\mu_{v-1}=0,\quad\mu_v\ne0.
 \tag{COI1}
\]
The coefficient map $C_a^{\mathsf T}:\mathbb C^{q'}\to\mathbb C^q$
is exactly the unscaled convolution by $E_A$:
\[
 \mathcal N_{C_a^{\mathsf T}\eta}(z)
       =E_A(z)\sum_{\lambda\in\mathrm{lower}}\eta_\lambda e^{\lambda z}.
 \tag{COI2}
\]
The original DNE2--3, NSO12 and NCT6--7 specify
$J_v=P_v^{\mathsf T}\mathbb VC_a^{\mathsf T}$ and
$Z_I=\mathbb V^{-1}P_vE_{J_I}$, where $I$ is its selected invertible
$q'$-row minor and $J_I$ is the complementary set. These sources are
pinned in \cite{a03cm:program:DNE,a03cm:program:NCT}.

\subsection{The first consecutive block is evaluated completely}

Leibniz differentiation of COI2 gives, for $r\ge0$,
\[
 (J_v)_{r,\lambda}
       =\sum_{j=v}^{v+r}\binom{v+r}{j}\mu_j\lambda^{v+r-j}.
 \tag{COI3}
\]
This is a degree-$r$ polynomial in the actual lower root $\lambda$,
with leading coefficient $\binom{v+r}{v}\mu_v$.
Let $I_0=(0,1,\ldots,q'-1)$ in the increasing derivative order,
and let $V_-$ be the original increasing-power lower Vandermonde.
Then $(J_v)_{I_0}=HV_-$, where the complete lower-triangular factor is
\[
 H_{rt}=\binom{v+r}{v+r-t}\mu_{v+r-t}\quad(0\le t\le r<q'),
 \qquad H_{rt}=0\quad(t>r).
 \tag{COI4}
\]
It follows that
\[
 \det(J_v)_{I_0}
  =\mu_v^{q'}\prod_{r=0}^{q'-1}\binom{v+r}{v}
       \prod_{\lambda<\lambda'}(\lambda'-\lambda)\ne0.
 \tag{COI5}
\]
All signs are with the original ordered lower-root columns and the
increasing rows. The Vandermonde determinant follows by alternation,
divisibility by every root difference, and its leading monomial of
coefficient one; no absolute value or ordering change is made.
The inverse is equally explicit: solve
$H_{rr}u_r=b_r-\sum_{t<r}H_{rt}u_t$ successively and then
apply $V_-^{-1}$. Thus every original conductor derivative in this
block remains in its inverse.

The exact consequence for a first-minor rule needs its actual order.
In increasing lexicographic enumeration of the increasing $q'$-tuples
of row indices, $I_0$ is the very first tuple, so COI5 evaluates that
rule before any search for a later nonzero determinant.
The original frozen DNE3, however, states a ``fixed ordering of all
subsets''; NSO12 and NCT7 refer to the first ordered nonzero minor.
Their printed lexicographic convention is for upper-root coefficient
columns. It does not specify an enumeration of the row subsets.
Accordingly COI5 proves the consecutive determinant but does not
alter an independently retained original choice $I$.
The exact calculation for that original choice is COI9--17 below.
This records a source-specification distinction and supplies the full
morphism; it does not discard the relation between the two charts.

\subsection{Every residue column and its full inverse}

Write $\chi(S)=\sum_{n=0}^{q}c_nS^n$, $c_q=1$, and define the
original residue map
\[
 \mathcal L(\nu)=\sum_\alpha\nu_\alpha\frac{\chi(S)}{S-\beta_\alpha},
 \qquad
 \mathfrak p_n(S)=\sum_{t=n+1}^{q}c_tS^{t-n-1}.
 \tag{COI6}
\]
The exact partial-fraction identity
$\mathcal L(\nu)/\chi=\sum_\alpha\nu_\alpha/(S-\beta_\alpha)$
has Laurent coefficients
$[S^{-n-1}]\mathcal L(\nu)/\chi=\sum_\alpha\nu_\alpha\beta_\alpha^n$.
The first $q$ coefficients determine its degree-below-$q$ numerator:
after multiplication by monic $\chi$, any remaining tail is strictly
negative in $S$ and has no polynomial part. This proves
$\mathcal L\mathbb V^{-1}e_n=\mathfrak p_n$ for every $0\le n<q$,
including all coefficients and phases.
For the consecutive block, $J_0=\{q',\ldots,q-v-1\}$, whence
\[
 \begin{gathered}
 L_0=\mathcal L\mathbb V^{-1}P_vE_{J_0},\qquad
 L_0e_j=\mathfrak p_{q-m+j}
       =\sum_{t=q-m+j+1}^{q}c_tS^{t-q+m-j-1},\quad0\le j<m,\\
 \operatorname{im}L_0=\mathcal P_{m-1},\qquad
 (L_0^{-1}P)_j=[S^{-q+m-j-1}]\frac{P(S)}{\chi(S)}.
 \end{gathered}                                                   \tag{COI7}
\]
The first equality has one polynomial of every degree
$m-1,m-2,\ldots,0$, each monic. Its coefficient matrix in reversed
degree order is triangular with diagonal one, proving the asserted
image and invertibility. The displayed Laurent extraction is its
inverse by the coefficient identity just proved. If $\deg P<m$,
all earlier coefficients $[S^{-n-1}]P/\chi$, $n<q-m$, are zero
because the rational function starts at order $S^{m-1-q}$.
For an arbitrary original $I$ the exact image instead is
\[
 L_I=\mathcal L\mathbb V^{-1}P_vE_{J_I},\qquad
 \operatorname{im}L_I
       =\operatorname{span}\{\mathfrak p_{v+j}:j\in J_I\}.
 \tag{COI8}
\]
These are original injections, not isometric identifications of
their arithmetic pullbacks with native minimum metrics.
In fact COI8 equals $\mathcal P_{m-1}$ exactly for $I=I_0$.
If a selected index $v+j$ is less than $q-m$, its individual
column has degree $q-v-j-1\ge m$ and cannot belong to that
low-degree space. Otherwise all $m$ selected indices lie in
the $m$-element set $\{q-m,\ldots,q-1\}$ and hence equal it;
their complement is precisely $I_0$.

\subsection{The complete transition for every original minor}

Retain the actually selected $I$ and define, with ordered coordinate
extractions in $\mathbb C^{q-v}$,
\[
 D_{0I}=(J_v)_{I_0}^{-1}E_{I_0}^{\mathsf T}E_{J_I},\qquad
 M_{0I}=E_{J_0}^{\mathsf T}E_{J_I}-(J_v)_{J_0}D_{0I}.
 \tag{COI9}
\]
Extraction in the $I_0$ rows and the $J_0$ rows verifies exactly
\[
 E_{J_I}=J_vD_{0I}+E_{J_0}M_{0I}.
 \tag{COI10}
\]
The induced $M_{0I}$ is invertible. Indeed both
$[J_v,E_{J_I}]$ and $[J_v,E_{J_0}]$ are square isomorphisms;
their transition is the block triangular matrix
$\left(\begin{smallmatrix}I&D_{0I}\\0&M_{0I}\end{smallmatrix}\right)$.
To see the first isomorphism directly, extraction in $I$ of a
relation between its columns annihilates its first coefficient,
and then its complementary coefficient vanishes; the second uses
COI5. Consequently no determinant or inverse is merely presumed.
Using the literal convolution identity gives the full exact
arithmetic transition
\[
 Z_I=C_a^{\mathsf T}D_{0I}+Z_0M_{0I},\qquad
 \boxed{L_I=F_CD_{0I}+L_0M_{0I}},\qquad
 F_C=\mathcal LC_a^{\mathsf T}.
 \tag{COI11}
\]
The term $F_CD_{0I}$ is an actual arithmetic vector. It remains
under the injective full-unit theta realization and cannot be
removed there as a relation. On the native side, dividing its
numerator by $E_A$ gives exactly the lower exponential
$\sum_\lambda(D_{0I}z)_\lambda e^{\lambda z}$.
That vector is among the very same lower evaluations minimized
in the original native metric. Translation of the full affine
fibre by it is a bijection of that fibre. Therefore, at every
original cutoff and source order, the exact native congruence is
\[
 D_{I,a,N-v}^{(s)}=M_{0I}^*D_{0,a,N-v}^{(s)}M_{0I},
             \qquad s=1,a.
 \tag{COI12}
\]
This proves the native comparison by its actual minimum, rather
than imposing it on the different arithmetic metric.

For the full arithmetic quotient forms put
$G_{00,N}=L_0^*G_NL_0$, $G_{CC,N}=F_C^*G_NF_C$ and
$G_{0C,N}=L_0^*G_NF_C$, $G_{C0,N}=G_{0C,N}^*$.
Multiplication of COI11 gives the complete receiver
\[
 \boxed{\begin{aligned}
 L_I^*G_NL_I={}&M_{0I}^*G_{00,N}M_{0I}
             +D_{0I}^*G_{CC,N}D_{0I}\\
 &+M_{0I}^*G_{0C,N}D_{0I}
             +D_{0I}^*G_{C0,N}M_{0I}.
 \end{aligned}}                                                   \tag{COI13}
\]
The identical expansion holds with $G_N$ replaced by the full
centered action form $A^*G_N+G_NA-aG_N$, and with every boundary
difference $G_{N_l}-G_{N_h}$. Thus no complete-action or boundary
cross block is lost in transport.

For a more quantitative bound on a genuine original chart, put
$\Delta G_N=G_{q-1}-G_N\succeq0$.
For $u=L_0M_{0I}z$ and $w=F_CD_{0I}z$, CAL13 gives
$u^*\Delta G_Nu\le\varepsilon_a\|u\|_{G_{q-1}}^2$.
The complete minimum also gives
$w^*\Delta G_Nw\le\|w\|_{G_{q-1}}^2$.
The Cauchy--Schwarz inequality for the positive form $\Delta G_N$
follows by minimizing $(u+tw)^*\Delta G_N(u+tw)$ over complex $t$;
it remains true at its zero directions. Expanding the square proves
\[
 \begin{split}
 0\le z^*L_I^*\Delta G_NL_Iz
 \le\left(\sqrt{\varepsilon_a}
       \|L_0M_{0I}z\|_{G_{q-1}}
       +\|F_CD_{0I}z\|_{G_{q-1}}\right)^2.
 \end{split}                                                     \tag{COI14}
\]
The exact cross block remains COI13; COI14 is not substituted for
the last conductor norm. For example applying
$2xy\le x^2+y^2$ to the displayed nonnegative terms gives
the valid whole-space inequality
\[
 L_I^*\Delta G_NL_I
 \preceq2\varepsilon_a M_{0I}^*G_{00,q-1}M_{0I}
             +2D_{0I}^*G_{CC,q-1}D_{0I}.
 \tag{COI15}
\]
The sharper unexpanded form COI14 can be used on each actual
direction. For four copies use $I_4\otimes M_{0I}$ and
$I_4\otimes D_{0I}$ in all these expressions. A native coefficient
flag $X$ passes to $(I_4\otimes M_{0I})X$ in the consecutive
coordinates, while its arithmetic image still carries the entire
conductor term in COI11. These are the exact typed maps needed to
apply a metric estimate to the original flags.

In the specified lexicographic row-subset enumeration, $I=I_0$
has $D_{0I}=0$ and $M_{0I}=I_m$ by COI9. This is the exact
specialization of all preceding formulas; COI13 then becomes the
low-degree form controlled in CAL13--18. For any other retained
enumeration, COI9 computes its actual conductor matrix explicitly,
and COI13--15 are its full receiving formulas.

\subsection{Primary intersections and the actual invariant closure}

For a monic divisor $g$ of $\chi$ of degree $J$, polynomial
cancellation in the integral domain $\mathbb C[S]$ proves
\[
 \ker g(A)=(\chi/g)\mathcal P_{J-1},\qquad
 \mathcal P_{m-1}\cap\ker g(A)
      =(\chi/g)\mathcal P_{m-q+J-1}.
 \tag{COI16}
\]
Indeed $\chi\mid gP$ is equivalent to $\chi/g\mid P$;
the degree-below-$q$ representative then gives the first degree
bound, and membership in $\mathcal P_{m-1}$ gives the second.
Multiplication by $\chi/g$ is injective on that literal polynomial
space. The intersection dimension is therefore
$\max(0,m+J-q)$, with negative-degree polynomial spaces equal
to zero. This proves terminal ($J=1$) and reflected-pair ($J=2$)
exclusion on the original $m\le q-2$ domain.
Since $a$ is odd, exactly half of the grid has positive real
displacement; its corresponding divisor has $J=q/2$.
For $a\ge29$, the elementary inequality
$16a-48-v\le(a+1)^2/2$ shows that the intersection is zero.
The original full-unit arithmetic injection preserves these
intersection equalities because it is injective.

The exact increasing invariant hull is
\[
 \sum_{r=0}^{d}A^r\mathcal P_{m-1}
      =\mathcal P_{\min(q-1,m+d-1)},\qquad d\ge0.
 \tag{COI17}
\]
Before the right side reaches degree $q-1$, multiplication by $S$
is literal and successively produces every next monomial. Once
the full quotient basis has appeared, the left side equals $E$
and remains there. This proves both inclusions and every rank,
including the four-copy rank $4\min(q,m+d)$.
The corresponding full-space action and its original metric cross
terms are retained in the companion action calculation.


% Independent finite derivation of the complete attained low-degree return.
\section{Complete-source localization before any relation minimum}

This supplies the complete proof of the arithmetic concentration
calculation in \cite{a03cm:program:arrival03-08}, with a stronger finite
factorial tail and an explicit original-chart receiver.
Keep the original actual simple quartet, its entire unit $v_h=2\xi/h$,
the original conductor and branch, and the original degree
$a=k+4$, $k=4\ell+1\ge9$. Set
\[
 \begin{gathered}
 c=a/2,\quad q=(a+1)^2,\quad q'=(a-7)^2,\quad
 m=q-q'-v=16a-48-v,\\
 \beta_{rt}=c+(2r-a)\delta+i(2t-a)\gamma,\quad
 \chi(S)=\prod_{r,t=0}^{a}(S-\beta_{rt}),\\
 0<\delta<1/2,\quad\gamma>2,\quad0\le v\le80,\\
 w_h(y)=|v_h(1/2+iy)|^2/(2\pi),\quad
 d\mu_a(y)=w_h^{*a}(y)\,dy,\quad
 d\sigma(y)=\frac{|\Gamma(1/4+iy/2)|^2}{2\pi}\,dy.
 \end{gathered}                                                   \tag{CAL1}
\]
The full masses are $(\int w_h)^a$ and $\sqrt{2\pi}$.
Every polynomial norm in this section is evaluated at the literal
physical coordinate $S=c+iy$, on both real half-lines. The arithmetic
source order is $a$. The two native Gamma orders $s=1,a$ enter the
separate native comparison, and are not substituted into $d\mu_a$.

\subsection{The original source envelope and a sharper finite polynomial bound}

The complete TW/BT bounds, retained in RTM2 and GRA3--4
\cite{a03cm:program:RTM,a03cm:program:GRA}, use the original positive constants
$c_h,C_h^{\rm tilt}$ and, on the simple branch, are
\[
 \begin{gathered}
 \alpha_*=\pi/2,\quad
 \vartheta_h=\int_{-1}^1w_h(y)\,dy>0,\quad
 M_{\alpha_*}=\int e^{\alpha_*y}w_h(y)\,dy\in(0,\infty),\\
 w_h^{*3}(y)\ge c_he^{-\alpha_*|y|}(1+|y|)^{-50},\quad
 e^{\alpha_*y}w_h(y)\le C_h^{\rm tilt}(1+|y|)^{-7/2},\\
 c_\Gamma e^{-\alpha_*|y|}(1+|y|)^{-1/2}
 \le\sigma(y)\le
 C_\Gamma e^{-\alpha_*|y|}(1+|y|)^{-1/2},\\
 c_\Gamma=\sqrt2e^{-7/6},\qquad
 C_\Gamma=\sqrt{2\sqrt5}e^{2/3}.
 \end{gathered}                                                   \tag{CAL2}
\]
Here $c_h$ is a three-factor convolution constant, not a pointwise
positive lower bound for the one-factor zeta numerator.
For clarity, the passage to the all-factor envelope is reproduced.
In the last $a-3$ convolution variables, restrict each to $[-1,1]$.
Their sum $z$ has $|z|\le a-3$ and their retained product mass is
$\vartheta_h^{a-3}$. Positivity and CAL2 therefore give
\[
 w_h^{*a}(y)\ge c_h\vartheta_h^{a-3}
 e^{-\alpha_*(|y|+a-3)}(a-2+|y|)^{-50}.
\]
For the upper envelope put $f(y)=e^{\alpha_*y}w_h(y)$.
On a convolution fibre whose coordinates sum to $y\ge0$, at least
one coordinate is at least $y/a$. Cover by these $a$ choices,
apply the bound on that coordinate, and integrate every other
coordinate with its full mass $M_{\alpha_*}$. Evenness supplies
the negative half-line. This proves
\[
 w_h^{*a}(y)\le
 aC_h^{\rm tilt}M_{\alpha_*}^{a-1}e^{-\alpha_*|y|}
                    (1+|y|/a)^{-7/2}.
\]
Consequently the following complete constants satisfy the literal
pointwise inequalities:
\[
 \begin{gathered}
 b_a=\frac{c_h\vartheta_h^{a-3}e^{-\alpha_*(a-3)}}
            {C_\Gamma2^{25}(a-2)^{50}},\qquad
 B_a=\frac{C_h^{\rm tilt}}{c_\Gamma}
                 a^{9/2}M_{\alpha_*}^{a-1},\\
 b_a(1+y^2)^{-25}\sigma(y)
       \le w_h^{*a}(y)\le B_a\sigma(y).
 \end{gathered}                                                   \tag{CAL3}
\]
For the lower bound use
$a-2+|y|\le(a-2)(1+|y|)$ and
$(1+|y|)^{50}\le2^{25}(1+y^2)^{25}$, dropping only the positive
factor $(1+|y|)^{1/2}\ge1$. For the upper bound use
$1+|y|/a\ge(1+|y|)/a$; the displayed $a^{9/2}$ is exactly
$a\cdot a^{7/2}$.

The original Gamma orthonormal recurrence is
\[
 ye_n=\sqrt{(n+1)(n+1/2)}e_{n+1}
                  +\sqrt{n(n-1/2)}e_{n-1}.
\]
Its Meixner--Pollaczek constants are those of
\cite[18.19.8--9,18.22.8]{a03cm:human:dlmf18}. The raising and lowering
shifts on degree at most $D$ have norms at most $D+1$ and $D$,
respectively. The triangle inequality therefore proves
$\|yp\|_\sigma\le2(D+1)\|p\|_\sigma$.
There is a sharper way to use this one-step estimate than iterating
twenty-five times. Set
$z_0=\|yp\|_\sigma^2/\|p\|_\sigma^2$ for $p\ne0$.
The second derivative of $f(z)=(1+z)^{-25}$ is positive on $z\ge0$,
so $f(z)\ge f(z_0)+f'(z_0)(z-z_0)$ by integrating its increasing
first derivative. Multiply this inequality at $z=y^2$ by
$|p(y)|^2d\sigma(y)$ and integrate. The linear term cancels exactly.
Using $z_0\le4(D+1)^2$ gives, also trivially for $p=0$,
\[
 \frac{b_a}{[1+4(D+1)^2]^{25}}\|P\|_\sigma^2
       \le\|P\|_{\mu_a}^2\le B_a\|P\|_\sigma^2,
 \quad\deg P\le D,\qquad
 K_{a,D}=\frac{B_a}{b_a}[1+4(D+1)^2]^{25}.
 \tag{CAL4}
\]
Here $p(y)=P(c+iy)$ is the actual polynomial. No source measure or
polynomial vector was replaced. Taking logarithms of the explicit
positive constants proves
$\log K_{a,D}=O_h(a+\log(D+2))$ uniformly in $a,D$.

\subsection{Every full relation is concentrated outside the inner interval}

Keep precisely the original finite quantities and guards
\[
 R=a\sqrt{\delta^2+\gamma^2},\qquad
 Y=2^{-16}q,\qquad d=\lfloor Y/16\rfloor,\qquad
 R\le Y/4,\quad m+1\le Y/16,\quad d\ge1.
 \tag{CAL5}
\]
The physical roots of $i^{-q}\chi(c+iy)$ have modulus at most $R$.
Thus $|\chi(c+iy)|\le(Y+R)^q$ when $|y|\le Y$, whereas
$|\chi(c+iy)|\ge(q-R)^q$ on $[q,2q]$. Every factor of the full
polynomial is used in these triangle inequalities.

We give the complete finite polynomial estimate used between these
two regions. Let $L_n$ denote the Legendre polynomial on $[-1,1]$.
Rodrigues' formula
$L_n(x)=(2^nn!)^{-1}(d/dx)^n(x^2-1)^n$ gives orthogonality by
$n$ integrations by parts, with no endpoint terms. Its leading
coefficient is $(2n)!/(2^n(n!)^2)$. Moreover
\[
 \int_{-1}^1x^nL_n(x)\,dx
  =2^{-n}\int_{-1}^1(1-x^2)^n\,dx
  =\frac{2^{n+1}(n!)^2}{(2n+1)!}.
\]
The last elementary integral follows inductively by integration by
parts, starting with $2$ at $n=0$. Subtracting the lower-degree
part of $L_n$ in the first factor now gives
$\int L_n^2=2/(2n+1)$. These are the standard classical formulas
\cite[18.3,18.5.5]{a03cm:human:dlmf18}, with their required proof included.
Leibniz applied to $(x-1)^n(x+1)^n$ also gives
\[
 L_n(x)=2^{-n}\sum_{j=0}^n
    \binom nj^2(x-1)^{n-j}(x+1)^j.
\]
For $|x|\le4$ it follows that
$|L_n(x)|\le(5/2)^n\binom{2n}{n}\le10^n$; the sum of squared
binomial coefficients is the coefficient of $z^n$ in
$(1+z)^n(1+z)^n$, and $\binom{2n}{n}\le4^n$.

On $[q,2q]$ the functions
$\sqrt{(2n+1)/q}\,L_n(2y/q-3)$ are therefore an orthonormal
polynomial basis. For $|y|\le Y$ their argument has modulus less
than $4$. Expanding any polynomial $u$ of degree at most $q$ in
that basis and applying complex Cauchy--Schwarz proves
\[
 |u(y)|^2\le\frac{(q+1)^2}{q}100^q
                        \int_q^{2q}|u(t)|^2\,dt\quad(|y|\le Y).
 \tag{CAL6}
\]
The exact sum of the coefficients $2n+1$ is $(q+1)^2$.
Apply CAL6 to $u(y)=Q(c+iy)$ for an arbitrary
$F=\chi Q\in\chi\mathcal P_q$. The full Gamma mass bounds its
inner integral; on $[q,2q]$, CAL2 gives
$\sigma(y)\ge c_\Gamma e^{-\pi q}/\sqrt{1+2q}$.
Consequently
\[
 \begin{split}
 \int_{|y|\le Y}|F(c+iy)|^2d\mu_a
       &\le\eta_a\|F\|_{\mu_a}^2,\\
 \eta_a&=K_{a,2q}\frac{\sqrt{2\pi}}{c_\Gamma}
      \frac{(q+1)^2\sqrt{1+2q}}q\,
       e^{\pi q}100^q
       \left(\frac{Y+R}{q-R}\right)^{2q}.
 \end{split}                                                     \tag{CAL7}
\]
In the final transfer the degree of $F$ is at most $2q$, so
CAL4 applies with exactly that degree. Since
$(Y+R)/(q-R)\le2^{-15}$ by CAL5, this proves
\[
 \log\eta_a\le
  -c_\eta q+O_h(a+\log q),\qquad
 c_\eta=30\log2-\pi-\log100>0.
 \tag{CAL8}
\]
Both the central interval and the exterior source remain in the
complete norms; no relation row or cross pairing was dropped.

\subsection{Every actual low-degree polynomial is concentrated inside}

For $P\in\mathcal P_m$, successive applications of the Gamma
multiplication bound at the actual intermediate degrees give
\[
 \|y^d P(c+iy)\|_\sigma
 \le 2^d\frac{(m+d)!}{m!}\|P\|_\sigma.
\]
On $|y|>Y$, the factor $|y|^{2d}$ is at least $Y^{2d}$.
Using the upper and lower inequalities of CAL4 on this same
$P$ therefore proves the sharpened finite bound
\[
 \begin{gathered}
 \int_{|y|>Y}|P(c+iy)|^2d\mu_a\le\zeta_a^*\|P\|_{\mu_a}^2,\\
 \zeta_a^*=K_{a,m}
       \left(\frac{2^d(m+d)!}{m!Y^d}\right)^2
       \le K_{a,m}4^{-2d}\le K_{a,m}3^{-2d}.
 \end{gathered}                                                   \tag{CAL9}
\]
For the first comparison, every factor
$2(m+j)/Y$, $1\le j\le d$, is at most $1/4$, since
$m+1+d\le Y/8$ by CAL5. This improves the supplied note's
$3^{-2d}$ without changing its auxiliary integer or any original
space. Since $d\ge2^{-20}q-1$,
\[
 \log\zeta_a^*\le
       -c_\zeta q+O_h(a+\log q),\qquad
 c_\zeta=2^{-19}\log4>0.
 \tag{CAL10}
\]
The factorial expression in CAL9, rather than its weaker displayed
upper estimates, can be used in every subsequent finite receiver.

Let $\mathcal R=\chi\mathcal P_q$ with its original arithmetic
inner product, and let $\Pi_{\mathcal R}$ be its attained
orthogonal projector. For every unit $F\in\mathcal R$, split
$\langle F,P\rangle$ at $|y|=Y$. Apply Cauchy--Schwarz separately
to the two full regions, using CAL7 on the first and CAL9 on the
second. Taking the supremum over unit $F$ gives the exact
whole-relation estimate
\[
 \|\Pi_{\mathcal R}P\|_{\mu_a}^2\le\varepsilon_a\|P\|_{\mu_a}^2,
 \qquad
 \varepsilon_a=(\sqrt{\eta_a}+\sqrt{\zeta_a^*})^2,
 \qquad P\in\mathcal P_m.
 \tag{CAL11}
\]
The supremum equals the projection norm: Cauchy--Schwarz bounds it
above, and its nonzero projected vector divided by its own norm
attains the bound. This argument incurs no row count or unestimated
relation cross block. The literal finite values CAL7 and CAL9 give
\[
 \varepsilon_a\le
       \exp\{-c_0q+O_h(a+\log q)\},\qquad
 c_0=\min\{c_\eta,c_\zeta\}>0.
 \tag{CAL12}
\]
This is an upper estimate. It is not an assertion of equality with
that exponential for every smaller choice of decay constant.
The condition $\varepsilon_a<1$, when needed for a logarithm, is a
direct finite check on CAL7 and CAL9 and holds eventually because
$q=(a+1)^2$ and all original data in their constants are fixed.

\subsection{Attained metrics, all fixed flags, and both boundary bands}

Let $E=\mathbb C[S]/(\chi)$ in its original increasing-power
representatives, and let $G_N$ be the minimum arithmetic form over
$\mathcal P_N\to E$. For $q-1\le N\le2q$, the fibre above a
literal $P\in\mathcal P_m$ is precisely
$P+\chi\mathcal P_{N-q}$, with a zero relation space at $N=q-1$.
Write $\Pi_N$ for the projector onto that entire relation space.
Expansion into the orthogonal sum gives
\[
 G_N[P]=\|P-\Pi_NP\|_{\mu_a}^2
       =\|P\|_{\mu_a}^2-\|\Pi_NP\|_{\mu_a}^2,
 \qquad
 (1-\varepsilon_a)G_{q-1}|_{\mathcal P_m}
       \preceq G_N|_{\mathcal P_m}
       \preceq G_{q-1}|_{\mathcal P_m}.
 \tag{CAL13}
\]
The smaller relation space is contained in $\mathcal R$, so its
projection norm is no larger. This proves CAL13 before choosing
any coefficient flag or taking any further minimum.

For the consecutive coefficient map $L_0$ whose image is proved
to be $\mathcal P_{m-1}$, put $M_0=I_4\otimes L_0$ and
$\mathcal G_N^0=M_0^*\operatorname{diag}(G_N,4)M_0$.
The map $M_0$ is precisely specified; the following statement does
not silently replace another actual minor by this one.
CAL13 gives
\[
 (1-\varepsilon_a)\mathcal G_{q-1}^0
       \preceq\mathcal G_N^0\preceq\mathcal G_{q-1}^0.
 \tag{CAL14}
\]
Every fixed restriction, and every successive attained quotient of
such a restriction, satisfies the same two bounds. For a restriction
substitute its original columns in the form inequality. For a
quotient, every representative of one value obeys the inequality;
minimizing both sides in the identical full affine fibre preserves
both bounds. Successive minima range over the union of those fibres,
so this proof also covers every original nested flag, including
$B,S,W,J,I,K_c,K_i,T,W_i$, after its actual coefficient map has
been specified. Empty value spaces have determinant one.

For any such attained form $H_{X,N}$ of dimension $r$, at the four
original cutoffs its signed arithmetic return satisfies
\[
 \begin{split}
 0\le F_X^{\rm ar}:={}&
 \log\det H_{X,q-1}+\log\det H_{X,q}
 -\log\det H_{X,2q-1}-\log\det H_{X,2q}\\
 &\le-2r\log(1-\varepsilon_a).
 \end{split}                                                     \tag{CAL15}
\]
Indeed set $l_N=\log\det H_{X,N}-\log\det H_{X,q-1}$.
Nested source fibres make these forms and logarithms decrease,
so $0\ge l_q\ge l_{2q-1}\ge l_{2q}$. The return is
$l_q-l_{2q-1}-l_{2q}\ge0$; CAL14 gives each high logarithm
at least $r\log(1-\varepsilon_a)$, while $l_q\le0$.
This proves the stated coefficient $2r$.
For example the full four-copy form has $r=4m$, giving $8m$.
For $\varepsilon_a\le1/2$, the identity
$-\log(1-x)=\int_0^x(1-t)^{-1}dt\le2x$ proves
$F_X^{\rm ar}\le4r\varepsilon_a$. Since $r\le4m=O(a)$,
CAL12 gives exponential smallness, in particular zero at every
fixed polynomial scale, on this precisely specified image.

The two exact original boundary differences pulled through $M_0$ are
\[
 \Gamma_0^0=M_0^*\operatorname{diag}(G_{q-1}-G_{2q-1},4)M_0,
 \quad
 \Gamma_1^0=M_0^*\operatorname{diag}(G_q-G_{2q},4)M_0.
 \tag{CAL16}
\]
Positivity of successive projection losses and CAL13 give
\[
 0\preceq\Gamma_i^0\preceq
       \varepsilon_a\mathcal G_{q-1}^0,\qquad i=0,1.
 \tag{CAL17}
\]
Here the strict sign can also be proved directly on the stated
simple odd-degree domain, rather than imported. The physical
polynomial $\chi(c+iy)$ is strictly positive: its roots in the
physical coordinate have nonzero imaginary parts and occur in
conjugate pairs, and $q=(a+1)^2$ is divisible by four. For a
nonzero $P$ of degree below $m<q$, the first band contains $\chi P$,
and $\langle\chi P,P\rangle=\int\chi(c+iy)|P|^2d\mu_a>0$.
Thus the first projection loss is positive. For the second band,
let $e=\chi$, let $b=\langle e,P\rangle/\langle e,e\rangle$
and put $P_1=P-be$. This is the attained representative at $q$.
The polynomial $\chi P_1$ has degree at most $2q$ and lies in
the high relation space. Since $P_1\perp e$, subtracting the
$e$ component from $\chi P_1$ preserves its pairing with $P_1$.
That pairing equals $\int\chi|P_1|^2d\mu_a>0$; $P_1\ne0$
because a polynomial of degree below $q$ cannot be a nonzero
multiple of $\chi$. The second projection loss is therefore
positive as well. Applied to the nonzero components of four copies,
this proves $\Gamma_i^0\succ0$.
Consequently the exact inverse comparison is
\[
 \lambda_{\min}\!\left((\Gamma_i^0)^{-1/2}
       \mathcal G_{q-1}^0(\Gamma_i^0)^{-1/2}\right)
       \ge\varepsilon_a^{-1}.
 \tag{CAL18}
\]
This is a statement about the displayed original boundary comparison.
It does not compare $\Gamma_0^0$ with $\Gamma_1^0$, does not replace
the native Gamma metric, and does not evaluate any individual
native allocation. The exact original-minor return, including every
conductor correction when that minor differs from the consecutive
one, is calculated in the companion coefficient-image section.


% Additive mathematical derivation, arrival 03, note 01 sections 4--7.
% Citation keys are supplied in ACTION_REFERENCES.tex.
\section{Physical compression, retained metric defects, and invariant enlargement}

\subsection{The consecutive comparison image and the original source metric}

Retain the original simple-quartet domain
\[
\begin{gathered}
 a=k+4,\quad k=4\ell+1\geq9,\quad
 0<\delta<\tfrac12,\quad \gamma>2,\\
 q=(a+1)^2,\quad q'=(a-7)^2,\quad
 0\leq v\leq80,\quad m=q-q'-v>0,\quad c=\tfrac a2.
\end{gathered}
\]
In particular \(a\) is odd and \(m<q\).
Put
\begin{equation}
 \beta_{rt}=c+(2r-a)\delta+i(2t-a)\gamma,\qquad
 \chi(S)=\prod_{r,t=0}^{a}(S-\beta_{rt})
        =\sum_{j=0}^q c_jS^j,\qquad c_q=1 .
 \tag{CJA1}
\end{equation}
Let \(E=\mathbb C[S]/(\chi)\), with the literal increasing-power
representative of degree below \(q\).  Its action \(A\) is multiplication
by the original \(S\), followed by reduction modulo the complete \(\chi\).
The exact minor correction COI3--7 evaluates the consecutive
comparison injection \(L_0\), retaining the supplied note's provenance
\cite{a03cm:program:arrival03-01}.
Throughout CJA2--19a, \(L\) denotes \(L_0\), and
\(M=I_4\otimes L_0\).  The originally selected injection \(L_I\)
is retained separately in CJA31--37.
Its full conductor transition is used there; no row-subset ordering
is inferred from the order of root labels.
The consecutive injection is
\begin{equation}
 L:\mathbb C^m\longrightarrow E,\qquad
 Le_j=\ell_j(S)=
 \sum_{t=q-m+j+1}^{q}c_tS^{t-q+m-j-1},\quad 0\leq j<m .
 \tag{CJA2}
\end{equation}
The polynomial \(\ell_j\) is monic of degree \(m-1-j\).
Thus these columns are a basis of the literal subspace
\(\mathcal P_{m-1}\subset E\): elimination in descending degree proves
both spanning and independence.  The Jacobi calculation uses this
consecutive map, including every coefficient of \(\chi\).
It applies directly to the original map when the retained row-subset
rule selects this consecutive block.

The arithmetic measure, with its complete mass, is
\begin{equation}
 d\mu(y)=w_h^{*a}(y)\,dy,\qquad
 w_h(y)=\frac{|v_h(1/2+iy)|^2}{2\pi},\qquad
 v_h=2\xi/h,\qquad S=c+iy .
 \tag{CJA3}
\end{equation}
It is the original source of NCRP8 \cite{a03cm:prog:NCRP}.
Its density is even, positive almost everywhere, and has all polynomial
moments finite.  Evenness follows from conjugation of the complete
quartet and of \(v_h\), then from convolution; the complete quartet and
all convolution factors are retained.
Every nonzero polynomial has positive squared integral.
Write \(G=G_{q-1}\) for the \(q\)-by-\(q\) Gram in the literal
\(1,S,\ldots,S^{q-1}\) coordinates.
At this degree the remainder map is the identity and its fibre is a
singleton.  Consequently
\begin{equation}
 \|P\|_G^2=\int_{\mathbb R}|P(c+iy)|^2\,d\mu(y),
 \qquad \deg P<q.
 \tag{CJA4}
\end{equation}
This is precisely the original attained arithmetic metric at \(q-1\).
The coefficient metric used below is \(G_L=L^*GL\).
The native conductor minimum and the original boundary forms
\(\Gamma_i\) keep their own values; no equality of those metrics with
\(G_L\) is used.

\subsection{The exact physical Jacobi compression and its nonzero residual}

Apply orthogonal elimination to \(1,y,y^2,\ldots\) in the measure
(CJA3).  Let \(\pi_r(y)\) be the resulting monic real polynomials,
\(\omega_r=\int|\pi_r|^2\,d\mu>0\), and
\(\varphi_r=\pi_r/\sqrt{\omega_r}\).
The full source mass is retained: in particular
\(\varphi_0=1/\sqrt{\int d\mu}\).
Uniqueness of monic orthogonal polynomials follows because the
difference of two candidates has lower degree and is orthogonal to
itself.  Evenness of the measure and uniqueness give
\(\pi_r(-y)=(-1)^r\pi_r(y)\).
Set
\begin{equation}
 a_r=\sqrt{\omega_r/\omega_{r-1}}>0\quad(r\geq1),\qquad
 a_0=0,\qquad\varphi_{-1}=0.
 \tag{CJA5}
\end{equation}
These original norm ratios are not assigned external values.

The exact recurrence is
\begin{equation}
 y\varphi_r=a_{r+1}\varphi_{r+1}+a_r\varphi_{r-1}.
 \tag{CJA6}
\end{equation}
Here is the complete argument in this measure.
For \(s\leq r-2\),
\(\langle\varphi_s,y\varphi_r\rangle
=\langle y\varphi_s,\varphi_r\rangle=0\) by degree.
The coefficient along \(\varphi_r\) vanishes by parity.
Comparison of leading coefficients gives \(a_{r+1}\) on
\(\varphi_{r+1}\).  Symmetry of multiplication by the real \(y\),
and the same leading-coefficient calculation at \(r-1\), give
\(a_r\) on \(\varphi_{r-1}\).
This proves every coefficient.
The general recurrence and its orthogonal-polynomial provenance are
recorded in NIST DLMF, section 18.2(iv)
\cite{a03cm:human:dlmf18}; the proof here supplies its specialization
without changing the arithmetic measure.

Define the unique matrix \(T\in\mathbb C^{m\times m}\) by
\begin{equation}
 (Lz)(c+iy)=\sum_{r=0}^{m-1}(Tz)_r\varphi_r(y),\qquad
 T_{rj}=\int_{\mathbb R}\varphi_r(y)\ell_j(c+iy)\,d\mu(y).
 \tag{CJA7}
\end{equation}
Both systems in this identity are bases of \(\mathcal P_{m-1}\);
therefore \(T\) is invertible.
Integrating the identity against itself proves the exact pullback
\begin{equation}
 T^*T=L^*GL=G_L .
 \tag{CJA8}
\end{equation}
Thus \(T\) is an isometry from the coefficient space with its specified
\(G_L\) metric into orthonormal coordinates.  It is not a replacement
of the metric on that coefficient space.

For four copies set
\[
\begin{gathered}
 H=(\mathbb C^m)^4,\quad
 \mathbf E=E^4,\quad
 M=I_4\otimes L,\quad
 \mathbf G=I_4\otimes G,\\
 \mathbf A=I_4\otimes A,\quad
 \mathcal T=I_4\otimes T,\quad
 G_H=M^*\mathbf GM=\mathcal T^*\mathcal T.
\end{gathered}
\]
Let \(J_m\) have zero diagonal and off-diagonals
\(a_1,\ldots,a_{m-1}\), and put \(\mathbf J=I_4\otimes J_m\).
The source-metric compression of this consecutive map and its residual are
\begin{equation}
 A_H=G_H^{-1}M^*\mathbf G\mathbf AM,\qquad
 R=\mathbf AM-MA_H .
 \tag{CJA9}
\end{equation}
By definition \(M^*\mathbf GR=0\).
Since \(m<q\), multiplication of a polynomial of degree below \(m\)
by \(S\) has degree at most \(m<q\).  No reduction by \(\chi\)
occurs.  Applying (CJA6) to that multiplication and dropping only its
orthogonal degree-\(m\) component proves
\begin{equation}
 \boxed{\mathcal TA_H\mathcal T^{-1}=cI+i\mathbf J,\qquad
 R\mathcal T^{-1}=i a_m\mathcal U_m.}
 \tag{CJA10}
\end{equation}
Here \(\mathcal U_m:\mathbb C^{4m}\to\mathbf E\) is completely specified:
in each copy it sends a coordinate column to its last entry times the
polynomial \(\varphi_m((S-c)/i)\), and has no other component.
Those four target polynomials are orthonormal in the direct-sum
metric \(\mathbf G\).  Hence \(R\) has rank exactly four, and its four
nonzero singular values from \((H,G_H)\) to
\((\mathbf E,\mathbf G)\) are exactly \(a_m\).
More explicitly, if \(P_{\rm last}\) projects onto the last coordinate
of each of the four copies, then
\begin{equation}
 R^*\mathbf GR=a_m^2\mathcal T^*P_{\rm last}\mathcal T .
 \tag{CJA11}
\end{equation}
This retains the domain and codomain norms of the residual.

The determinant \(D_r(y)=\det(yI_r-J_r)\) obeys
\(D_0=1\), \(D_1=y\), and
\(D_{r+1}=yD_r-a_r^2D_{r-1}\), by expansion on the last row.
Multiplying (CJA6) by \(\sqrt{\omega_r}\) gives exactly the same
monic recurrence for \(\pi_r\).
Induction proves \(D_r=\pi_r\).
The real symmetric \(J_m\) has an orthonormal eigenbasis and real
eigenvalues: maximizing its Rayleigh quotient gives an eigenvector,
its orthogonal complement is invariant, and induction diagonalizes it.
An eigenvector is determined by its first entry, because its
successive equations have nonzero off-diagonal coefficients.
If that entry vanishes, all entries vanish.  Each eigenspace thus has
dimension at most one; the orthonormal diagonalization proves all
eigenvalues are simple.
Writing these eigenvalues as \(t_1,\ldots,t_m\), we have proved
\begin{equation}
 \operatorname{spec}A_H=\{c+it_j:1\leq j\leq m\},
 \quad\text{each repeated four times},\qquad
 A_H^*G_H+G_HA_H-aG_H=0.
 \tag{CJA12}
\end{equation}
The latter identity also follows immediately by substituting
(CJA10) and \(2c=a\).
It is accompanied by the exact nonzero residual (CJA10).

\subsection{Every boundary-metric and action cross term}

We first give the literal COB endpoint maps \cite{a03cm:prog:COB}.
For \(n\geq q-1\), let \(J_n^\chi:\mathcal P_n\to E\) be remainder,
let \(H_n\) be the original moment Gram, and let
\[
 \mathscr R_n=H_n^{-1}(J_n^\chi)^*
     \bigl(J_n^\chi H_n^{-1}(J_n^\chi)^*\bigr)^{-1}.
 \]
The identity \(J_n^\chi\mathscr R_n=I\), together with
\(\mathscr R_n^*H_n\ker J_n^\chi=0\), proves its exact attained
minimum property by the Pythagorean identity on each affine fibre.
Define the two distinct original windows
\[
 \Delta_0=\mathscr R_{2q-1}-\mathscr R_{q-1},\qquad
 \Delta_1=\mathscr R_{2q}-\mathscr R_q .
 \]
All these polynomial-valued maps are embedded in the same original
polynomial space before they are subtracted.
On the physical line put \(Q(y)=\chi(c+iy)\).
Since \(a\) is odd, the real displacements \(2r-a\) are nonzero.
Pairing \(r\) with \(a-r\), and using that \(q\) is divisible by
four, writes \(Q\) as the product of positive factors
\(((y-(2t-a)\gamma)^2+((2r-a)\delta)^2)\), one for each pair.
Thus \(Q\) is real, even, and strictly positive.

Let \(\pi_r^\nu\) be the real monic orthogonal polynomials for
\(d\nu=Q^2d\mu\), with squared norms \(\nu_r>0\), and put
\[
 e_r(S)=\frac{\chi(S)\pi_r^\nu((S-c)/i)}{\sqrt{\nu_r}},
 \qquad c_r(x)=-\langle e_r,x\rangle_\mu,\qquad
 C_0=(c_0,\ldots,c_{q-1})^{\mathsf T},\quad
 C_1=(c_1,\ldots,c_q)^{\mathsf T}.
 \]
The argument used in (CJA6) applies to the even positive measure
\(d\nu\).  These \(e_r\) are orthonormal in the original measure
\(d\mu\), and they span the successive original relation spaces.
Projecting the literal representative \(x\in\mathcal P_{q-1}\)
against those spaces proves
\begin{equation}
 \Delta_0x=\sum_{r=0}^{q-1}e_rc_r(x),\qquad
 \Delta_1x=\sum_{r=1}^{q}e_rc_r(x),\qquad
 M_i^{\rm one}=C_i^*C_i .
 \tag{CJA12a}
\end{equation}
For the second equality, the common \(e_0c_0(x)\) term cancels
between the two minima at \(q\) and \(2q\).
Both maps are injective.
For \(\Delta_0x=0\), the common minimum is \(P=x\), and its
orthogonality to the relation \(\chi P\) at degree \(2q-1\)
contradicts \(\int Q|P|^2d\mu>0\) unless \(x=0\).
For \(\Delta_1x=0\), use the common minimum
\(P=\mathscr R_qx\) of degree at most \(q\), and its relation
\(\chi P\) of degree at most \(2q\), to obtain the same conclusion.
Therefore \(C_i\) is invertible and \(M_i^{\rm one}>0\).

Write \(\boldsymbol\Delta_i=I_4\otimes\Delta_i\),
\(\mathsf M_i=I_4\otimes M_i^{\rm one}\).
For each original endpoint retain the pullback of its actual COB form
along the consecutive map:
\[
 \Gamma_i=M^*\mathsf M_iM,\qquad
 Q_i=\mathcal T^{-*}\Gamma_i\mathcal T^{-1}.
\]
Here \(\mathsf M_i\) is the full original form on \(\mathbf E\)
just defined by both original attained minima.
Direct substitution of (CJA10), with no commutation assumption, gives
\begin{equation}
 \boxed{\mathcal T^{-*}
 (A_H^*\Gamma_i+\Gamma_iA_H-a\Gamma_i)\mathcal T^{-1}
       =i(Q_i\mathbf J-\mathbf JQ_i).}
 \tag{CJA13}
\end{equation}
Thus this compressed defect vanishes exactly when
\([Q_i,\mathbf J]=0\).
Its arithmetic-relative trace is
\[
 \operatorname{tr}\{G_H^{-1}
 (A_H^*\Gamma_i+\Gamma_iA_H-a\Gamma_i)\}
   =i\operatorname{tr}(Q_i\mathbf J-\mathbf JQ_i)=0 .
\]
When \(\Gamma_i>0\), its own relative trace is also zero:
\(\operatorname{tr}(Q_i^{-1}i[Q_i,\mathbf J])=0\).
The first assertion does not require inverting \(\Gamma_i\).

The full pulled-back ambient action contains both additional terms:
\begin{equation}
 \begin{split}
 M^*(\mathbf A^*\mathsf M_i+\mathsf M_i\mathbf A-a\mathsf M_i)M
 ={}&A_H^*\Gamma_i+\Gamma_iA_H-a\Gamma_i\\
   &+R^*\mathsf M_iM+M^*\mathsf M_iR .
 \end{split}
 \tag{CJA14}
\end{equation}
To prove it, substitute \(\mathbf AM=MA_H+R\) into the two
action terms and multiply out.  Orthogonality
\(M^*\mathbf GR=0\) applies to \(\mathbf G\);
it does not remove either of the displayed \(\mathsf M_i\) cross terms.
In the physical orthonormal coordinates, put
\(M_0=M\mathcal T^{-1}\), \(R_0=R\mathcal T^{-1}=ia_m\mathcal U_m\).
The last two terms become, exactly,
\begin{equation}
 R_0^*\mathsf M_iM_0+M_0^*\mathsf M_iR_0
 =ia_m(M_0^*\mathsf M_i\mathcal U_m
              -\mathcal U_m^*\mathsf M_iM_0).
 \tag{CJA15}
\end{equation}
All four degree-\(m\) boundary directions remain in this expression.

For any specified transported coefficient cone subspace \(X=B,S,W\)
in these consecutive coordinates,
take a full independent column frame \(I_X\) and its original
arithmetic projector
\[
 P_X^{G_H}=I_X(I_X^*G_HI_X)^{-1}I_X^*G_H .
\]
The inverse is the complete inverse on that frame.
Equation (CJA8) gives
\(\mathcal TP_X^{G_H}\mathcal T^{-1}=P_{\mathcal TX}\),
the Euclidean orthogonal projector onto the actual transformed
subspace.  Hence
\begin{equation}
 \mathcal T(I-P_X^{G_H})A_HP_X^{G_H}\mathcal T^{-1}
     =i(I-P_{\mathcal TX})\mathbf JP_{\mathcal TX}.
 \tag{CJA16}
\end{equation}
This is an exact preservation test for the compression.
The complete failure of the specified inclusion \(M\) to intertwine
\(\mathbf A\) is still \(R\), as given in (CJA10).

The proper-source action receives the same residual.
Put \(\mathfrak a_r=\sqrt{\nu_r/\nu_{r-1}}>0\) for \(r\geq1\).
Multiplying the relation recurrence by \(\chi\) gives
\[
 Se_r=ce_r+i\mathfrak a_{r+1}e_{r+1}
                      +i\mathfrak a_re_{r-1}.
 \]
Let \(J_0^{\rm rel}\) be its \(q\)-by-\(q\) Jacobi matrix on
indices \(0,\ldots,q-1\), and \(J_1^{\rm rel}\) the one on
indices \(1,\ldots,q\).  Their off-diagonals are respectively
\(\mathfrak a_1,\ldots,\mathfrak a_{q-1}\) and
\(\mathfrak a_2,\ldots,\mathfrak a_q\).
Define the exact maps
\[
 \begin{split}
 A_i^{\rm tan}&=C_i^{-1}(cI+iJ_i^{\rm rel})C_i,\qquad
 C_i^{\rm act}=A_i^{\rm tan}-A,\\
 N_0&=i\mathfrak a_qe_qc_{q-1},\\
 N_1&=i\mathfrak a_1e_0c_1+
                         i\mathfrak a_{q+1}e_{q+1}c_q .
 \end{split}
 \]
Retain their four-copy maps
\(\boldsymbol C_i=I_4\otimes C_i^{\rm act}\) and
\(\boldsymbol N_i=I_4\otimes N_i\).
The displayed recurrence, projected onto the stated band and its
orthogonal complement, proves
\begin{equation}
 S\boldsymbol\Delta_i-\boldsymbol\Delta_i\mathbf A
    =\boldsymbol\Delta_i\boldsymbol C_i+\boldsymbol N_i,
 \qquad
 \boldsymbol\Delta_i^*\mathsf H\boldsymbol N_i=0.
 \tag{CJA16a}
\end{equation}
Here \(\mathsf H=I_4\otimes H_{2q+1}\) is the original moment
metric on the polynomial representatives through degree \(2q+1\).
Both \(\boldsymbol\Delta_i\) and \(\boldsymbol N_i\) are embedded
in that same coefficient space.
Their degrees are at most \(2q\) and \(2q+1\), respectively.
They and their action edges inject into
\((\mathbb C[S]/\chi^{a+1})^4\), because
\((a+1)q>2q+1\).
Thus (CJA16a) is an identity of the original proper-source maps,
with a specified norm on the displayed polynomial bands and edge.
No norm on the entire quotient is being introduced.

In this notation,
\[
 \beta_i=\boldsymbol\Delta_iM,\qquad
 \mathsf M_i=\boldsymbol\Delta_i^*\mathsf H\boldsymbol\Delta_i .
\]
Composing (CJA16a) with \(M\), and using (CJA9), proves
\begin{equation}
 K_i:=S\beta_i-\beta_iA_H
 =\boldsymbol\Delta_iR+
   \boldsymbol\Delta_i\boldsymbol C_iM+\boldsymbol N_iM.
 \tag{CJA17}
\end{equation}
For clarity, the three terms here are maps from \(H\) to the original
proper-source value space.
Writing
\(U_i=\boldsymbol\Delta_iR\),
\(V_i=\boldsymbol\Delta_i\boldsymbol C_iM\), and
\(W_i=\boldsymbol N_iM\), its original squared residual form is
\begin{equation}
 \begin{split}
 K_i^*\mathsf HK_i={}&
 U_i^*\mathsf HU_i+V_i^*\mathsf HV_i+W_i^*\mathsf HW_i\\
 &+U_i^*\mathsf HV_i+V_i^*\mathsf HU_i\\
 &+U_i^*\mathsf HW_i+W_i^*\mathsf HU_i\\
 &+V_i^*\mathsf HW_i+W_i^*\mathsf HV_i .
 \end{split}
 \tag{CJA18}
\end{equation}
All nine directed terms have now been displayed.  The last four
cross terms are exactly zero by (CJA16a), with its proved band
orthogonality.  The remaining original form is
\begin{equation}
 \begin{split}
 K_i^*\mathsf HK_i={}&R^*\mathsf M_iR
  +R^*\mathsf M_i\boldsymbol C_iM
  +M^*\boldsymbol C_i^*\mathsf M_iR\\
 &+M^*\boldsymbol C_i^*\mathsf M_i\boldsymbol C_iM
  +M^*\boldsymbol N_i^*\mathsf H\boldsymbol N_iM .
 \end{split}
 \tag{CJA18a}
\end{equation}
In particular the two Gamma-metric cross terms involving \(R\)
remain.  The normal energies are exactly
\begin{equation}
 \begin{split}
 \boldsymbol N_0^*\mathsf H\boldsymbol N_0
 &=I_4\otimes\mathfrak a_q^2c_{q-1}^*c_{q-1},\\
 \boldsymbol N_1^*\mathsf H\boldsymbol N_1
 &=I_4\otimes
  (\mathfrak a_1^2c_1^*c_1+
                  \mathfrak a_{q+1}^2c_q^*c_q).
 \end{split}
 \tag{CJA18b}
\end{equation}
These follow from orthonormality of the different \(e_r\);
neither original boundary edge is discarded.
Similarly the full proper-source centered form is
\begin{equation}
 \begin{split}
 (S\beta_i)^*\mathsf H\beta_i+\beta_i^*\mathsf H(S\beta_i)-a\Gamma_i
 ={}&A_H^*\Gamma_i+\Gamma_iA_H-a\Gamma_i\\
 &+R^*\mathsf M_iM+M^*\mathsf M_iR\\
 &+M^*\boldsymbol C_i^*\mathsf M_iM
       +M^*\mathsf M_i\boldsymbol C_iM\\
 &+M^*\boldsymbol N_i^*\mathsf H\boldsymbol\Delta_iM
       +M^*\boldsymbol\Delta_i^*\mathsf H\boldsymbol N_iM .
 \end{split}
 \tag{CJA19}
\end{equation}
This follows by first substituting
\(S\beta_i=\beta_iA_H+K_i\), then inserting (CJA17).
On the actual line \(\overline S+S=a\), so the left side is
exactly zero, pointwise before integration.  The last two terms
are zero by (CJA16a).  Put \(B_i=R+\boldsymbol C_iM\).
The complete original weight identity is consequently
\begin{equation}
 \begin{split}
 A_H^*\Gamma_i+\Gamma_iA_H-a\Gamma_i
   &=-M^*\mathsf M_iB_i-B_i^*\mathsf M_iM,\\
 i[Q_i,\mathbf J]
   &=-M_0^*\mathsf M_i(B_i\mathcal T^{-1})
       -(B_i\mathcal T^{-1})^*\mathsf M_iM_0 .
 \end{split}
 \tag{CJA19a}
\end{equation}
The centered scalar is the same retained \(a\) throughout.
The surviving normal leakage is measured by (CJA18a)--(CJA18b),
while the weight identity retains both the arithmetic compression
residual and the boundary tangential correction.

\subsection{The actual cyclic enlargement and its positive trace}

The one-copy COB proper-source packet is
\(E_2=\mathbb C[S]/(\chi^{a+1})\).
For its original nonzero boundary vector, with the originally selected
\(L_I\), COB7--9 give the literal
form
\begin{equation}
 p=\Delta_iL_Iz=\chi R_p\pmod{\chi^{a+1}},
 \qquad 0\ne p,\quad \deg R_p\leq q .
 \tag{CJA20}
\end{equation}
The subscript prevents confusion with the action residual \(R\).
The monic greatest common divisor
\(g_p=\gcd(\chi^a,R_p)\) gives the exact annihilator and exact
cyclic-module map
\begin{equation}
 f_p=\frac{\chi^a}{g_p},\qquad
 \mathbb C[S]/(f_p)\ \longrightarrow\ \mathbb C[S]p\subset E_2,
 \qquad [F]\longmapsto Fp .
 \tag{CJA21}
\end{equation}
This map is an isomorphism of \(\mathbb C[S]\)-modules.
Indeed, \(Fp=0\) is equivalent in the original quotient to
\(\chi^{a+1}\mid F\chi R_p\), and cancellation of the nonzero
polynomial \(\chi\) gives \(\chi^a\mid FR_p\).
At each of the distinct roots \(\beta\) of \(\chi\) this requires
\(\operatorname{ord}_\beta F\geq
a-\min(a,\operatorname{ord}_\beta R_p)\), and requires no other
factor.  These are exactly the factors of \(f_p\).
The map is onto by the definition of the cyclic module, and the
calculated kernel proves injectivity.
Writing \(d_\beta=\min(a,\operatorname{ord}_\beta R_p)\), we obtain
\begin{equation}
 \dim\mathbb C[S]p
   =aq-\sum_\beta d_\beta\geq(a-1)q,\qquad
 \sum_\beta d_\beta=\deg g_p\leq\deg R_p\leq q .
 \tag{CJA22}
\end{equation}
Thus every invariant subspace containing this actual nonzero vector
contains its explicitly given cyclic module.
The original full-unit proper-source map in COB.8 is
\[
 \Psi_2:E_2\longrightarrow Q_h^{\otimes a},\qquad
 [P]\longmapsto
 [P(D_1+\cdots+D_a)F_h^{\otimes a}],
 \quad Q_h=\mathscr B/\Theta h(D)W_* .
 \]
It is injective and intertwines \(S\) with the sum of the original
Euler actions.  Here is the algebra behind that retained assertion.
The original identities are \(h(D)F_h=\Theta\phi_*\),
\(\mathcal MF_h=v_h\), \(g=2\xi=hv_h\), and
\(W_*=\mathbb C[D]\phi_*\).
If \(P(D)F_h=\Theta h(D)T(D)\phi_*\), its original Mellin
identity is \(Pv_h=ghT=h^2v_hT\), so
\(P=h^2T\), since \(v_h\) is not the zero entire function.
The converse follows by applying \(h(D)T(D)\) to
\(h(D)F_h=\Theta\phi_*\), using the same original Euler
intertwining.  The single-factor cyclic annihilator is exactly
\(h^2\).
Tensoring the resulting injections over \(\mathbb C\) preserves
injectivity.  In every ordered root tuple the local nilpotents
\(\epsilon_1,\ldots,\epsilon_a\) satisfy \(\epsilon_j^2=0\);
their sum has \(a\)-th power
\(a!\epsilon_1\cdots\epsilon_a\ne0\) and \((a+1)\)-st power zero.
Every original grid sum occurs, so the cyclic sum has precisely
the annihilator \(\chi^{a+1}\).
The actual Taylor observation multiplies by the full unit
\([v_h]_{h^2}^{\otimes a}\), which is invertible at every original
simple root and is retained.
These facts prove the stated map and injectivity with the original
source labels.
Now \(F\Psi_2(p)=\Psi_2(Fp)\); injectivity proves equality of the
two annihilators, without changing or discarding that unit.

Here is a direct metric proof of the needed positive-trace estimate.
Let \(B_p\) denote multiplication by \(S\) on the cyclic module
(CJA21), let \(H_p>0\) be any specified positive metric on that
module, and define
\[
 D_p=B_p^*H_p+H_pB_p-aH_p,\qquad
 K_p=H_p^{-1/2}D_pH_p^{-1/2}.
\]
The polynomial factorization of \(f_p\) gives its primary subspaces
with dimensions \(a-d_\beta\).
In the local basis \(1,S-\beta,\ldots,(S-\beta)^{a-d_\beta-1}\),
multiplication is \(\beta I\) plus a nilpotent shift.
This proves its eigenvalue and trace, retaining all multiplicities.
Let \(I_+\) be any full basis of the direct sum of the primary
subspaces with \(\operatorname{Re}\beta>c\), and write
\(B_pI_+=I_+B_+\), \(H_+=I_+^*H_pI_+\).
The orthogonal projector in Euclidean coordinates onto
\(H_p^{1/2}\operatorname{im}I_+\) is
\[
 P_+=H_p^{1/2}I_+H_+^{-1}I_+^*H_p^{1/2}.
 \]
Cyclicity of matrix trace and the invariant-subspace identity give
\[
 \operatorname{tr}(P_+K_p)
  =\operatorname{tr}\{H_+^{-1}I_+^*D_pI_+\}
  =2\operatorname{Re}\operatorname{tr}B_+
                  -a\dim\operatorname{im}I_+ .
\]
For every Hermitian \(K\) and orthogonal projector \(P\),
\(\operatorname{tr}(PK)\leq\operatorname{tr}K_+\).
To verify this last elementary inequality, diagonalize \(K\) in an
orthonormal basis: the diagonal entries of \(P\) lie in \([0,1]\),
so each positive eigenvalue contributes at most itself, and each
negative eigenvalue contributes a nonpositive number.
We have therefore proved, for the original metric,
\begin{equation}
 \operatorname{tr}(K_p)_+
 \geq\sum_{\beta}(a-d_\beta)
               \max(2\operatorname{Re}\beta-a,0).
 \tag{CJA23}
\end{equation}
The notation on the left means the trace of the positive spectral
part of \(K_p\), rather than the positive part of its scalar trace.

For odd \(a\), the full simple-packet positive sum is, explicitly,
\[
 L_{h,a}
 =2\delta(a+1)\sum_{r=(a+1)/2}^{a}(2r-a)
 =\frac{\delta(a+1)^3}{2}
 =\frac{\delta q(a+1)}2 .
\]
Since the largest positive displacement is \(2a\delta\), (CJA22)
and (CJA23) yield the required finite bound
\begin{equation}
 \operatorname{tr}(K_p)_+
 \geq aL_{h,a}-2a\delta q
 =\frac{\delta aq(a-3)}2 .
 \tag{CJA24}
\end{equation}
The exact preceding spectral sum remains available with the actual
\(d_\beta\); the inequality does not assign them generic values.

\subsection{The consecutive invariant hull and exact Sylvester inverse}

There is no nonzero \(\mathbf A\)-invariant subspace contained in
\(\operatorname{im}M=(\mathcal P_{m-1})^4\).
Indeed, in a nonzero such subspace choose a vector whose largest
component degree \(d\) is maximal among all its vectors.
Then \(d\leq m-1\).
Multiplication by \(S\) raises that degree to \(d+1\leq m<q\);
no reduction modulo \(\chi\) is available to change the leading
coefficient.
The resulting vector has larger degree, contradicting invariance.
This proves that the invariant core is exactly zero.
On the other hand, the invariant enlargement is exact:
\begin{equation}
 \sum_{r=0}^{d}\mathbf A^r\operatorname{im}M
   =\bigl(\mathcal P_{\min(q-1,m+d-1)}\bigr)^4
 \quad(d\geq0).
 \tag{CJA25}
\end{equation}
For \(m+d-1<q\), ordinary polynomial multiplication proves both
containments, since every power up to \(m+d-1\) occurs.
At \(d=q-m\) this is all of \(\mathbf E\), and it stays all of
\(\mathbf E\) for every larger \(d\).

Let \(\mathcal V=I_4\otimes V\) be the original root-evaluation
map, \(V_{\beta n}=\beta^n\), \(0\leq n<q\).
Distinctness of the complete grid makes \(V\) invertible: a
degree-\((q-1)\) polynomial vanishing at all \(q\) roots is zero.
Let \(O\) be a real orthogonal eigenvector matrix for \(J_m\),
with \(O^*J_mO=\operatorname{diag}(t_1,\ldots,t_m)\), and put
\[
 \mathcal O=I_4\otimes O,\qquad
 \mathcal W=\mathcal O^*\mathcal T,\qquad
 K=\mathcal V^{-*}\mathbf G\mathcal V^{-1}.
 \]
Thus \(\mathcal W^*\mathcal W=G_H\).
The original target metric in root-value coordinates is the full
positive matrix \(K\), not the identity.

For a map \(Y:H\to\mathbf E\) the equation
\(\mathbf AX-XA_H=Y\) has the exact solution
\begin{equation}
 \begin{gathered}
 Z=\mathcal VY\mathcal W^{-1},\qquad
 \widehat X_{(b,\beta),(e,j)}
   =\frac{Z_{(b,\beta),(e,j)}}{\beta-c-it_j},\\
 X=\mathcal V^{-1}\widehat X\mathcal W .
 \end{gathered}
 \tag{CJA26}
\end{equation}
The indices \(b,e\) range independently over the four copies;
this formula includes all sixteen copy-to-copy blocks.
Indeed, conjugation by these exact maps turns the Sylvester equation
into those scalar equations, because the target action is the
diagonal matrix of all original \(\beta\), and the source action is
the diagonal matrix of all \(c+it_j\).
The real part of each denominator is \((2r-a)\delta\).
Since \(a\) is odd, it is nonzero and has absolute value at least
\(\delta\).
Thus every entry is uniquely determined, proving invertibility on
the full original space \(\operatorname{Hom}(H,\mathbf E)\).

The entrywise bound does not by itself give a
\(\delta^{-1}\) bound in the original metrics.
The exact metric formula can be stated without that inference.
Use the original Hilbert--Schmidt norm
\[
 \|X\|_{\mathrm{HS}(G_H,\mathbf G)}^2
 =\operatorname{tr}(G_H^{-1}X^*\mathbf GX).
 \]
Since \(\mathcal W\) is the proved source isometry, transport gives
\begin{equation}
 \|X\|_{\mathrm{HS}(G_H,\mathbf G)}^2
      =\operatorname{tr}(\widehat X^*K\widehat X).
 \tag{CJA27}
\end{equation}
For source index \(j\), let
\(D_j=\operatorname{diag}_{(b,\beta)}(\beta-c-it_j)\).
Different source columns in (CJA27) contribute separately.
For each column the exact norm of multiplication by \(D_j^{-1}\)
in the target \(K\) metric is
\(\|K^{1/2}D_j^{-1}K^{-1/2}\|_2\).
Taking the largest column norm, and choosing a maximizing vector
in that column for the reverse inequality, proves
\begin{equation}
 \boxed{\|\mathscr S^{-1}\|_{\mathrm{HS}(G_H,\mathbf G)}
 =\max_j\|K^{1/2}D_j^{-1}K^{-1/2}\|_2
 \leq\frac{\sqrt{\lambda_{\max}(K)/\lambda_{\min}(K)}}{\delta}.}
 \tag{CJA28}
\end{equation}
Here \(\mathscr S(X)=\mathbf AX-XA_H\).
Both target coordinate factors remain.
No bound on their condition ratio has been assumed.

For completeness, a different specified positive source form
\(\Gamma\), including the boundary pullback for the indicated map, is retained
by setting
\(Q_\Gamma=\mathcal W^{-*}\Gamma\mathcal W^{-1}\).
If columns are vectorized in their actual order, put
\[
 \mathcal H_\Gamma=\overline{Q_\Gamma^{-1}}\otimes K,\qquad
 \mathcal D=\operatorname{diag}_{((e,j),(b,\beta))}
                         (\beta-c-it_j).
 \]
The identity
\(\operatorname{tr}(Q_\Gamma^{-1}\widehat X^*K\widehat X)
=\operatorname{vec}(\widehat X)^*
 \mathcal H_\Gamma\operatorname{vec}(\widehat X)\)
is obtained by expanding all row and column indices.
Thus the corresponding exact original-metric inverse norm is
\begin{equation}
 \|\mathscr S^{-1}\|_{\mathrm{HS}(\Gamma,\mathbf G)}
 =\|\mathcal H_\Gamma^{1/2}\mathcal D^{-1}
             \mathcal H_\Gamma^{-1/2}\|_2 .
 \tag{CJA29}
\end{equation}
It retains both the source and target Gram changes, including
their cross entries.  No native metric is replaced by \(G_H\).

Finally the computed residual satisfies \(\mathscr S(M)=R\).
The complete correction equation
\(\mathbf A(M+X)=(M+X)A_H\) is therefore
\(\mathscr S(X)=-R\).
Invertibility just proved gives its unique solution
\begin{equation}
 \boxed{X=-M.}
 \tag{CJA30}
\end{equation}
This computes the correction for the consecutive source action \(A_H\),
the original target action \(\mathbf A\), and their full map space.
Its relation to enlargement is already explicit in (CJA25) and
(CJA20)--(CJA24): the small compressed image has zero invariant core,
while the corresponding invariant closure retains the original action
and grows by the stated exact module maps.

\subsection{The original selected minor: its core and complete action receiver}

We now retain the originally selected row set \(I\), its complement
\(J_I\), and its actual injection \(L_I\), without a consecutive-row
assumption.  Put
\[
 \mathcal K_I=\{v+j:j\in J_I\}\subset\{0,\ldots,q-1\},\qquad
 \mathfrak p_n(S)=\sum_{r=n+1}^q c_rS^{r-n-1}.
 \]
The residue calculation COI6--8 proves
\(\operatorname{im}L_I=\operatorname{span}\{\mathfrak p_n:
n\in\mathcal K_I\}\); the full family
\(\mathfrak p_0,\ldots,\mathfrak p_{q-1}\) is a basis because its
degrees are respectively \(q-1,\ldots,0\) with leading coefficient
one.
For each original root \(\beta\), the original Lagrange class is
\[
 v_\beta=\frac{\chi(S)}{(S-\beta)\chi'(\beta)}
       =\sum_{n=0}^{q-1}
               \frac{\beta^n}{\chi'(\beta)}\mathfrak p_n(S).
 \]
To prove the last identity, expand
\(\chi(S)/(S-\beta)=\sum_{n=0}^{q-1}\beta^n\mathfrak p_n(S)\);
the coefficient of \(S^r\) on the right is
\(\sum_{n=0}^{q-r-1}c_{n+r+1}\beta^n\), exactly the quotient
coefficient in division by \(S-\beta\).
All displayed coordinates of \(v_\beta\) are nonzero:
\(\chi'(\beta)\ne0\) by the simple grid, and
\(\operatorname{Re}\beta\geq a(1/2-\delta)>0\), so \(\beta\ne0\).
Since \(|\mathcal K_I|=m<q\), at least one nonzero coordinate is
omitted.  Therefore
\begin{equation}
 v_\beta\notin\operatorname{im}L_I
       \quad\text{for every original root and every retained minor.}
 \tag{CJA31}
\end{equation}

Let \(\mathcal U\) be a nonzero invariant subspace of
\(\mathbf E\) contained in \(\operatorname{im}M_I\), where
\(M_I=I_4\otimes L_I\).
For any nonzero \(u\in\mathcal U\), at least one original Lagrange
spectral projector \(P_\beta=v_\beta(\mathbf A)\) has
\(P_\beta u\ne0\), since their sum is the identity.
Here \(v_\beta(\mathbf A)\) means evaluation of its literal
Lagrange polynomial at \(\mathbf A\), on each of the four copies.
Each \(P_\beta\) is a polynomial in \(\mathbf A\), so this vector
still belongs to \(\mathcal U\).
It has the form
\((z_1v_\beta,z_2v_\beta,z_3v_\beta,z_4v_\beta)\), with some
\(z_b\ne0\).
Its membership in the four-copy image would imply
\(v_\beta\in\operatorname{im}L_I\), contradicting (CJA31).
This proves the actual all-minor strengthening
\begin{equation}
 \boxed{\operatorname{core}_{\mathbf A}(\operatorname{im}M_I)
             =\{0\}\qquad(\text{every retained }I).}
 \tag{CJA32}
\end{equation}
Its exact invariant hull is also given by original evaluations:
put
\(\mathcal B_I=\{\beta:\mathfrak p_n(\beta)\ne0
\text{ for some }n\in\mathcal K_I\}\).
The same spectral projectors prove
\[
 \mathbb C[\mathbf A]\operatorname{im}M_I
   =\left(\bigoplus_{\beta\in\mathcal B_I}
                       \mathbb Cv_\beta\right)^4 .
 \]
For inclusion from left to right, every missing root evaluation
stays zero under multiplication by \(\beta\).
For the reverse inclusion apply \(P_\beta\) to a column with nonzero
evaluation.  Thus this hull formula retains the actual selected
image rather than assigning it the consecutive hull without a map.

The exact original-to-consecutive transition COI9--11 is
\begin{equation}
 \begin{gathered}
 L_I=F_CD_{0I}+L_0M_{0I},\qquad M_I=Y+Z,\\
 Y=(I_4\otimes L_0)(I_4\otimes M_{0I}),\qquad
 Z=(I_4\otimes F_C)(I_4\otimes D_{0I}).
 \end{gathered}
 \tag{CJA33}
\end{equation}
Its matrices are evaluated using the full original conductor and
the selected invertible minor in COI9; \(M_{0I}\) is invertible
there by the full block transition.
In particular \(Z\) is a retained actual arithmetic vector, not
a discarded relation.
For every full Hermitian form \(Q\) on \(\mathbf E\), direct
multiplication gives the exact receiver
\begin{equation}
 M_I^*QM_I=Y^*QY+Z^*QZ+Y^*QZ+Z^*QY .
 \tag{CJA34}
\end{equation}
This applies separately to \(\mathbf G\), to each \(\mathsf M_i\),
and to the full centered action
\(\mathbf D=\mathbf A^*\mathbf G+\mathbf G\mathbf A-a\mathbf G\).
Define the actual quantities
\[
 G_I=M_I^*\mathbf GM_I,\quad
 A_I=G_I^{-1}M_I^*\mathbf G\mathbf AM_I,\quad
 R_I=\mathbf AM_I-M_IA_I,\quad
 \Gamma_{i,I}=M_I^*\mathsf M_iM_I.
 \]
They obey
\begin{equation}
 \begin{split}
 M_I^*\mathbf GR_I&=0,\\
 A_I^*G_I+G_IA_I-aG_I&=M_I^*\mathbf DM_I .
 \end{split}
 \tag{CJA35}
\end{equation}
The right-hand side retains all four terms of (CJA34), with
\(Q=\mathbf D\).
The proper-source identities (CJA17)--(CJA19) and the first line
of (CJA19a), including both
Gamma-metric cross terms and both normal edges, hold verbatim
with \(M,A_H,R,\Gamma_i\) replaced by
\(M_I,A_I,R_I,\Gamma_{i,I}\): their proofs used the actual
compression equation and the proved COB band orthogonality,
not consecutive degrees.
In particular
\[
 \begin{split}
 A_I^*\Gamma_{i,I}+\Gamma_{i,I}A_I-a\Gamma_{i,I}
 ={}&-M_I^*\mathsf M_i(R_I+\boldsymbol C_iM_I)\\
    &-(R_I+\boldsymbol C_iM_I)^*\mathsf M_iM_I .
 \end{split}
 \]

There is a further exact original-chart domain where the
critical-line and Sylvester conclusions follow without a
consecutive selector:
\begin{equation}
 \mathcal K_I\subset\{1,\ldots,q-1\}.
 \tag{CJA36}
\end{equation}
This holds for every original minor when \(v\geq1\); for \(v=0\)
it holds exactly when \(0\notin J_I\).
On this explicitly identified domain the degree of every
\(\mathfrak p_n\) in the image is at most \(q-2\).
Multiplication by \(S\) on that image is therefore literal
through degree \(q-1\), without polynomial reduction.
The physical identity \(\overline S+S=a\), applied to each pair
of original image polynomials, proves
\(M_I^*\mathbf DM_I=0\).
By (CJA35), take \(T_I=G_I^{1/2}\) and put
\begin{equation}
 J_I^{\rm phys}=-i(T_IA_IT_I^{-1}-cI).
 \quad\text{Then}\quad
 (J_I^{\rm phys})^*=J_I^{\rm phys}.
 \tag{CJA37}
\end{equation}
The proof is to conjugate the zero identity in (CJA35) by
\(T_I^{-1}\); the resulting centered action is skew-Hermitian.
This Hermitian compression need not be the tridiagonal \(J_m\);
no rank-four assertion is transferred to \(R_I\).
Diagonalize this actual finite Hermitian matrix by a unitary
\(O_I\).  Replacing \(\mathcal W\) in (CJA26) by \(O_I^*T_I\),
and \(t_j\) by its actual real eigenvalues, gives the exact
Sylvester inverse for
\(\mathbf AX-XA_I\).
Every denominator still has real part \((2r-a)\delta\ne0\).
The original-metric norm is exactly (CJA28) with these new
eigenvalues and the same retained target Gram \(K\);
the unique action-compatible correction is \(X=-M_I\).
For each \(\Gamma_{i,I}\), its compressed defect is the exact
commutator
\[
 T_I^{-*}(A_I^*\Gamma_{i,I}+\Gamma_{i,I}A_I-a\Gamma_{i,I})T_I^{-1}
  =i[T_I^{-*}\Gamma_{i,I}T_I^{-1},J_I^{\rm phys}],
 \]
and its own original-metric Sylvester norm is transported by
(CJA29) with \(O_I^*T_I\).

On the remaining possible original chart \(v=0,\ 0\in J_I\),
the degree-\((q-1)\) direction is present.
The actual action and all its centered terms are (CJA33)--(CJA35).
The physical critical-line argument above has an explicit
degree-\(q\) polynomial-reduction term there, so it does not
assign that chart a critical-line compression or the
consecutive Jacobi spectrum.
The all-minor zero-core theorem (CJA32), the actual cyclic
enlargement (CJA20)--(CJA24), and all complete proper-source
cross terms remain established on that chart as well.


% Complete additive derivation for arrival03 Note09.
% Original coefficient section retained. No document preamble or bibliography.
\section{The full native--arithmetic relative return on the original coefficients}

\subsection{Original forms, source orders, and orientations}
Retain the actual simple quartet, its period, branch, full conductor
$E_A$, and all original root guards. Write
\begin{equation}\tag{NMR1}\label{a03cm:nmr:domain}
\begin{gathered}
a=k+4\equiv1\pmod4,\quad q=(a+1)^2,\quad Q=(a-7)^2,\quad
\Delta=q-Q=16a-48,\\ v=\operatorname{ord}_0E_A,\quad
0\leq v\leq80,\quad m=\Delta-v>0,\\
s\in\{1,a\},\quad c_0=a/2,\quad c_1=a/2-4,\\
(N_0,N_1,N_2,N_3)=(q-1,q,2q-1,2q),\quad D_i=N_i-v,\\
\mathcal R f=-f(N_0)-f(N_1)+f(N_2)+f(N_3).
\end{gathered}
\end{equation}
The lower grid keeps source order $s$, not order $a-8$.
The original Gamma measure, mass, and polynomial norms are
\begin{equation}\tag{NMR2}\label{a03cm:nmr:gamma}
dm_s(y)=\frac{(2\pi)^{s/2}2^{s/2}}{4\pi\Gamma(s/2)}
 |\Gamma(s/4+iy/2)|^2\,dy,\qquad
h_{n,s}=(2\pi)^{s/2}n!(s/2)_n .
\end{equation}
Use its original orthonormal polynomials
$\phi_{n,s,c}(S)=p_{n,s}((S-c)/i)/\sqrt{h_{n,s}}$,
where $p_{0,s}=1$, $p_{1,s}=y$, and
$p_{n+1,s}=yp_{n,s}-n(n-1+s/2)p_{n-1,s}$.
These are the mass-retaining Meixner--Pollaczek formulas with
$p_{n,s}(y)=n!P_n^{(s/4)}(y/2;\pi/2)$
\cite{a03cm:human:dlmf18}.

Let $Z=\mathbb V^{-1}P_vE_{\mathcal J}$ be the actual DNE numerator
section selected by its actual admitted lower-conductor minor, and let
$L=\mathcal LZ:\mathbb C^m\to E_a=\mathbb C[S]/(\chi_a)$ be its
residue map, retaining $\chi_a'(\beta)$ in $\mathcal L$.
The arithmetic form $G_N^{\rm ar}$ is the attained polynomial
remainder metric for the complete source
$d\mu_a(y)=w_h^{*a}(y)\,dy$ on the physical line $S=c_0+iy$.
Set
\begin{equation}\tag{NMR3}\label{a03cm:nmr:forms}
\begin{gathered}
M=I_4\otimes L,\qquad
\mathcal G_N=M^*(G_N^{\rm ar})^{\oplus4}M,\qquad
\mathcal D_N=(D_{a,N-v}^{(s)})^{\oplus4},\\
T_N=\mathcal G_N^{-1/2}\mathcal D_N\mathcal G_N^{-1/2},\qquad
F_{\rm ar}=-\mathcal R\log\det\mathcal G_N,\qquad
\mathcal T_{a,s}=\mathcal R\log\det D_{a,N-v}^{(s)} .
\end{gathered}
\end{equation}
All are literal original coefficient forms. The positive square root
only represents the identity map between the two original Hilbert
metrics; it does not rescale either source or discard a coordinate.

\subsection{The native metric and its monotonicity before any determinant}
The native basis germs are
$F_b(z)=\sum_\beta Z_{\beta b}e^{\beta z}/E_A(z)$, analytic at zero.
At every $D\geq q-1-v$, let
\[
H_D[n,\beta']=\phi_{n,s,c_1}(\beta'),\qquad
J_D[n,b]=\Lambda_{F_b}(\phi_{n,s,c_1}),
\qquad0\leq n\leq D.
\]
Every derivative in this expression is retained. Explicitly,
with $n_{jb}=\sum_\beta Z_{\beta b}\beta^j$ and $f_{rb}=F_b^{(r)}(0)$,
\[
f_{rb}=
\frac{n_{v+r,b}-\sum_{j=v+1}^{v+r}
             \binom{v+r}{j}\mu_j f_{v+r-j,b}}
     {\binom{v+r}{v}\mu_v}.
\]
It follows directly by differentiating the original identity
$E_AF_b=\sum_\beta Z_{\beta b}e^{\beta z}$.
The exact original metric is
\begin{equation}\tag{NMR4}\label{a03cm:nmr:native-min}
x^*D_{a,D}^{(s)}x
 =\min_{\eta\in\mathbb C^Q}\|J_Dx+H_D\eta\|_2^2,\qquad
D_{a,D}^{(s)}
 =J_D^*J_D-J_D^*H_D(H_D^*H_D)^{-1}H_D^*J_D .
\end{equation}
The lower evaluation columns are independent. The combined columns
$[H_D,J_D]$ are independent: multiplying a germ by $E_A$ and taking
its first $q$ numerator moments identifies them with the independent
columns $[C_a^{\mathsf T},Z]$ through the upper Vandermonde.
Thus the minimum is positive for $x\ne0$ and its unique minimizer is
$-(H_D^*H_D)^{-1}H_D^*J_Dx$.

At $D+1$ the objective in \eqref{a03cm:nmr:native-min}, for each fixed
$x,\eta$, equals its old objective plus the squared modulus of the
new observation row. Taking the minimum over the identical complete
space $\mathbb C^Q$ therefore proves
\begin{equation}\tag{NMR5}\label{a03cm:nmr:monotonicity}
D_{a,D+1}^{(s)}\succeq D_{a,D}^{(s)}>0,\qquad
G_{N+1}^{\rm ar}\preceq G_N^{\rm ar},\qquad
\mathcal G_{N+1}\preceq\mathcal G_N .
\end{equation}
For the arithmetic inequalities, the full affine fibre of a remainder
$p$ grows from $p+\chi_a\mathbb C[S]_{\leq N-q}$ to
$p+\chi_a\mathbb C[S]_{\leq N+1-q}$, in the same original integral
norm. Minimizing over the larger fibre can only decrease its squared
norm. Pullback along the fixed $M$ preserves this inequality.
In particular $\mathcal T_{a,s}\geq0$ and $F_{\rm ar}\geq0$.

\subsection{Reconstruction of the native total and its signs}
Here the Gamma numerator metric is denoted $G_N^{(s)}$, to retain
its distinction from the arithmetic source. Let
$\mathsf E_{+,N}[n,\beta]=\phi_{n,s,c_0}(\beta)$ and
$K_N=\mathsf E_{+,N}^*\mathsf E_{+,N}$.
An exponential coefficient $\nu$ acts on a remainder coefficient
vector $x$ by $(\mathbb V\nu)^{\mathsf T}x$.
Its squared dual norm is
$(\mathbb V\nu)^*\overline{(G_N^{(s)})^{-1}}(\mathbb V\nu)$,
by solving for its Hermitian representing vector.
The same functional in the orthonormal polynomial coordinates has
squared norm $\nu^*K_N\nu$. Equality for every $\nu$ proves
\begin{equation}\tag{NMR6}\label{a03cm:nmr:dual}
K_N=\mathbb V^*\overline{(G_N^{(s)})^{-1}}\mathbb V,\qquad
\det K_N=\frac{|\det\mathbb V|^2}{\det G_N^{(s)}} .
\end{equation}
The conjugation and ordinary transpose are both retained.

Let $Z_v=\mathbb V^{-1}P_{<v}$ and
$\mathsf F=[C_a^{\mathsf T},Z,Z_v]$. This frame is invertible:
the first $v$ moment coordinates remove a putative $Z_v$ coefficient,
and the selected lower-conductor minor then proves independence of
the other columns. Define the three successive factors of
$\mathsf F^*K_N\mathsf F$:
\[
K_{C,N}=\overline C_aK_NC_a^{\mathsf T},\qquad
\widehat D_N=Z^*\{K_N-K_NC_a^{\mathsf T}K_{C,N}^{-1}
                                      \overline C_aK_N\}Z,
\]
and let $Q_{v,N}$ be the attained last-block Gram after minimizing
over every column of $[C_a^{\mathsf T},Z]$.
Subtracting the first-block minimizer and its conjugate row
operation has determinant one and splits off $K_{C,N}$.
Applying the same elimination to the second block splits off
$\widehat D_N$. The remaining minimum is exactly $Q_{v,N}$.
Consequently
\begin{equation}\tag{NMR7}\label{a03cm:nmr:factorization}
\det K_{C,N}\det\widehat D_N\det Q_{v,N}
=|\det\mathsf F|^2\det K_N
=\frac{|\det\mathsf F|^2|\det\mathbb V|^2}{\det G_N^{(s)}} .
\end{equation}
The two fixed frame factors are not set equal to one.

The lower original covariance is $K_{-,D}=H_D^*H_D$.
The literal polynomial conductor matrix $A_{D,s}$ in the original
Gamma frames has ordinary dual $B_{D,s}=A_{D,s}^{\mathsf T}$.
The exact product-rule square is
\[
B_{D,s}H_D=\mathsf E_{+,N}^{\rm red}C_a^{\mathsf T},
\qquad
B_{D,s}J_D=\mathsf E_{+,N}^{\rm red}Z,
\]
where the omitted first $v$ upper rows vanish on both numerator
subspaces. Thus $\widehat D_N$ is the minimum of
$\|B_{D,s}(J_Dx+H_D\eta)\|^2$ over that same $\eta$, with no
scalar $v!/\mu_v$ inserted.
Put
\begin{equation}\tag{NMR8}\label{a03cm:nmr:corrections}
\begin{gathered}
e_N^C=\log\det K_{C,N}-\log\det K_{-,N-v},\qquad
\epsilon_N=\log\det\widehat D_N-\log\det D_{a,N-v}^{(s)},\\
\mathcal B_a^{(s)}=-\mathcal R\log\det G_N^{(s)},\qquad
\mathcal C^-_{a,s}=\mathcal R\log\det K_{-,N-v}.
\end{gathered}
\end{equation}
The last quotient factor is
$Q_{v,N}=\overline{(P_{<v}^*G_N^{(s)}P_{<v})^{-1}}$.
Indeed its value map on dual coordinates is
$P_{<v}^{\mathsf T}\mathbb V$; its kernel is the full
vanishing-moment space, and the inverse of its covariance, using
\eqref{a03cm:nmr:dual}, is precisely the displayed conjugate inverse.
Hence
$\mathcal R\log\det Q_{v,N}=F_{L_v}^{(s)}$, the original
low-polynomial return, with signs $(+,+,-,-)$ on its primal Gram.
Applying $\mathcal R$ to \eqref{a03cm:nmr:factorization} cancels its fixed
frame factors and proves the literal native identity
\begin{equation}\tag{NMR9}\label{a03cm:nmr:total}
\mathcal T_{a,s}
=\mathcal B_a^{(s)}-\mathcal C^-_{a,s}
 -\mathcal R e^C-F_{L_v}^{(s)}-\mathcal R\epsilon .
\end{equation}
This independently verifies the signs and the one-copy meaning of
the LRC24 total.

For clarity its error control is on the same finite maps. The complete
LRC exterior bounds give
$-2I_{D,s,Q,1}\leq e_N^C\leq2F_{D,s,Q,1}$.
For the native quotient the induced map and inverse are
$P_{B Y\cap(BK)^\perp}B$ and $P_{Y\cap K^\perp}B^{-1}$,
where $Y=\operatorname{im}[H_D,J_D]$, $K=\operatorname{im}H_D$.
Each projection removes exactly the relevant denominator. Taking the
top exterior power in dimension $m$ gives
$-2I_{D,s,m,1}\leq\epsilon_N\leq2F_{D,s,m,1}$.
These are the original all-rank bounds, each $O_{\rm actual}(q)$;
no additional denominator-rank factor is charged.

The low return is also $O_{\rm actual}(q)$ with its fixed dimension
$v$. Explicitly retain the original full-remainder estimate
$K_{q-1}\preceq K_N\preceq Z_{a,s}K_{q-1}$, where
\begin{equation}\tag{NMR10}\label{a03cm:nmr:Z}
\begin{gathered}
t_q=\pi/2-1/q,\quad A_q=2\log3/t_q,\quad
S_q=\sum_{j=0}^{q-1}A_q^j,\\
B_{s,q}=[(9/25)\sin(1/q)]^{-s/2}
        \exp\{2a\sqrt{\delta^2+\gamma^2}(\log3+t_q/2)\},\\
Z_{a,s}=1+B_{s,q}S_q^2
           \sum_{n=q}^{2q}\frac{n!}{(s/2)_n}(25/16)^n .
\end{gathered}
\end{equation}
This is the original AKS whole-remainder bound, as proved in the
primitive-endpoint integration. Inverting \eqref{a03cm:nmr:dual},
restricting to the fixed low subspace, and using its actual dimension
$v$ gives $0\leq F_{L_v}^{(s)}\leq2v\log Z_{a,s}$.
Here $\log Z_{a,s}=O_{\rm actual}(q)$: one has
$n!/(s/2)_n\leq2n+1$, $\log S_q=O(q)$ and
$\log B_{s,q}=O_{\rm actual}(a\log q)=o(q)$.
Thus every error on the right of \eqref{a03cm:nmr:total} is $O(q)$.
At $v=0$ the low factor is the empty determinant one.

To check the evaluated coefficients rather than merely repeat LRC24,
write $\ell_s=(s-1)/4$ and use the proved literal covariance
expansions in their exact LRS21--22 source orders:
\begin{equation}\tag{NMR11}\label{a03cm:nmr:covariances}
\begin{aligned}
\mathcal B_a^{(s)}
 &=C_Bq^2+\ell_sq(2C_B-4a_1)+O_{\rm actual}(q),\\
\mathcal C^-_{a,1}
 &=C_Bq^2-16C_\partial aq+128q\log a+O_{\rm actual}(q),\\
\mathcal C^-_{a,a}
 &=C_Bq^2+(-16C_\partial+C_B/2-a_1)aq
                  +128q\log a+O_{\rm actual}(q).
\end{aligned}
\end{equation}
These use the lower interval starting at $m$ with width $q$ and
cutoff $N-v$, not a new natural lower window. The asymptotic
statements retain the original finite domains
$a\geq257$,
$a\sqrt{\delta^2+\gamma^2}\leq2^{-33}q$, and, for $b=a,a-8$,
$n_b=(b+1)^2+\ell_s\geq100$ and
$\max\{b\sqrt{\delta^2+\gamma^2},2\ell_s+1/2\}\leq2^{-17}n_b$.
All hold eventually for the fixed actual quartet.

For $s=1$ the upper $\ell_s$ term is zero. For $s=a$ it is exactly
$(a-1)q(C_B/2-a_1)$, whose $aq$ coefficient cancels the
$(C_B/2-a_1)aq$ term in the lower covariance; the remaining
$-(C_B/2-a_1)q$ is within the stated $O(q)$.
Subtraction in \eqref{a03cm:nmr:total} therefore proves
\begin{equation}\tag{NMR12}\label{a03cm:nmr:native-value}
\mathcal T_{a,s}
=16C_\partial aq-128q\log a+O_{\rm actual}(q),
\qquad s=1,a .
\end{equation}
The sign $-128$ is forced by subtracting the actual lower covariance.
The positive coefficient is the original
$C_\partial=2\int_0^1\log((1+\rho(t))/(1-\rho(t)))\,dt>0$:
the proved profile--covariance dictionary has
$0<\rho(t)<1$ and an integrable logarithmic endpoint at zero.
Thus no new outer coefficient has been substituted.

\subsection{Exact arithmetic correction for the actual section}
This step applies to the actual $L$, without assigning its image a
lower polynomial degree than its original minor proves.
Let $g_0=L^*G_{q-1}^{\rm ar}L$ in one copy.
Represent its columns by their literal remainders $l_1,\ldots,l_m$.
At cutoff $N$ write the complete relation frame
$r_j=\chi_aS^j$, $0\leq j\leq N-q$.
For $N=q-1$ it is empty. Define full integral matrices
\[
(H_N^{\rm ar})_{ij}=\langle r_i,r_j\rangle_{\mu_a},\qquad
(C_N^{\rm ar})_{ib}=\langle r_i,l_b\rangle_{\mu_a},
\]
with conjugate-linearity in the first argument. Completing the square
over every relation coefficient gives
\begin{equation}\tag{NMR13}\label{a03cm:nmr:arithmetic-min}
g_N=L^*G_N^{\rm ar}L
 =g_0-(C_N^{\rm ar})^*(H_N^{\rm ar})^{-1}C_N^{\rm ar}.
\end{equation}
No relation cross block is deleted. Independence of the original
remainders from the relation ideal implies $g_N>0$.
Define the finite, actual projection matrices in four copies:
\begin{equation}\tag{NMR14}\label{a03cm:nmr:projection}
S_i=\left[
g_0^{-1/2}(C_{N_i}^{\rm ar})^*
(H_{N_i}^{\rm ar})^{-1}C_{N_i}^{\rm ar}g_0^{-1/2}
\right]^{\oplus4}.
\end{equation}
The empty case is $S_0=0$. The nested relation projections and
\eqref{a03cm:nmr:arithmetic-min} give
$0=S_0\preceq S_1\preceq S_2\preceq S_3<I_{4m}$.
Taking determinants, the fixed $\det(g_0^{\oplus4})$ cancels and
proves
\begin{equation}\tag{NMR15}\label{a03cm:nmr:Far}
\begin{gathered}
F_{\rm ar}
=\log\det(I-S_1)-\log\det(I-S_2)-\log\det(I-S_3),\\
0\leq F_{\rm ar}\leq-2\log\det(I-S_3)
              \leq-8m\log(1-\beta_{\rm orig}),\qquad
\beta_{\rm orig}=\|S_3\|<1 .
\end{gathered}
\end{equation}
For the lower bound, $\log\det(I-S_1)\geq\log\det(I-S_2)$
and $-\log\det(I-S_3)\geq0$.
For the first upper bound, discard the nonpositive first term and
use $\log\det(I-S_2)\geq\log\det(I-S_3)$.
For the last upper bound, all $4m$ eigenvalues of $S_3$ are at most
$\beta_{\rm orig}$.
This is a finite original-matrix value and bound; it asserts no
unproved asymptotic rate for $\beta_{\rm orig}$.

An explicit coarse bound also applies uniformly to every original
section. The unchanged simple-quartet source estimates are
\[
b_a(1+y^2)^{-25}\,d\sigma(y)\leq d\mu_a(y)
\leq B_a\,d\sigma(y),\qquad d\sigma=dm_1,
\]
where the complete original constants are
\begin{equation}\tag{NMR16}\label{a03cm:nmr:source-constants}
\begin{gathered}
b_a=\frac{c_h\vartheta_h^{a-3}e^{-(\pi/2)(a-3)}}
 {C_\Gamma\,2^{25}(a-2)^{50}},\qquad
B_a=\frac{C_h^{\rm tilt}}{c_\Gamma}
                   a^{9/2}M_{\pi/2}^{a-1},\\
\vartheta_h=\int_{-1}^{1}w_h(y)\,dy,\qquad
M_{\pi/2}=\int_{\mathbb R}e^{(\pi/2)y}w_h(y)\,dy,\qquad
c_\Gamma=\sqrt2e^{-7/6},\quad
C_\Gamma=\sqrt{2\sqrt5}e^{2/3}.
\end{gathered}
\end{equation}
The constants $c_h,C_h^{\rm tilt}$ retain their original
convolution lower-envelope and tilted-envelope definitions.
For a polynomial of degree at most $D$, the Gamma recurrence gives
$\|yP\|_\sigma\leq2(D+1)\|P\|_\sigma$.
For nonzero $P$ put $I_0=\int|P|^2d\sigma$,
$I_2=\int y^2|P|^2d\sigma$ and $u=I_2/I_0$.
The convex function $\varphi(x)=(1+x)^{-25}$ satisfies its exact
tangent inequality
$\varphi(y^2)\geq\varphi(u)+\varphi'(u)(y^2-u)$.
Multiply by $|P|^2$ and integrate in the unchanged $d\sigma$.
The linear term is $\varphi'(u)(I_2-uI_0)=0$, so the result is at
least $\varphi(u)I_0\geq[1+4(D+1)^2]^{-25}I_0$.
The zero polynomial obeys the same inequality directly. Consequently
\[
b_a[1+4(D+1)^2]^{-25}\|P\|_\sigma^2
\leq\|P\|_{\mu_a}^2\leq B_a\|P\|_\sigma^2 .
\]
The physical translate $P(c_0+iy)$ remains its actual polynomial
in $y$; no Euclidean coefficient norm or source mass is changed.
Put
$K_{a,2q}=(B_a/b_a)[1+4(2q+1)^2]^{25}$.
Taking these inequalities on every original affine fibre and then
using \eqref{a03cm:nmr:Z} at order $s=1$ gives
\begin{equation}\tag{NMR17}\label{a03cm:nmr:coarse}
\begin{gathered}
\Omega_a=K_{a,2q}Z_{a,1},\qquad
\Omega_a^{-1}\mathcal G_{q-1}\preceq
                 \mathcal G_N\preceq\mathcal G_{q-1},\\
0\leq F_{\rm ar}\leq8m\log\Omega_a=O_{\rm actual,h}(aq).
\end{gathered}
\end{equation}
Indeed $G_N^{\rm ar}\succeq
b_a[1+4(2q+1)^2]^{-25}G_N^{(1)}
\succeq b_a[1+4(2q+1)^2]^{-25}Z_{a,1}^{-1}G_{q-1}^{(1)}$,
whereas $G_{q-1}^{\rm ar}\preceq B_aG_{q-1}^{(1)}$.
Pulling back proves the first line. The second uses the same two
high-endpoint determinant bound as \eqref{a03cm:nmr:Far}.
Here $\log K_{a,2q}=O_h(a+\log q)$ and $\log Z_{a,1}=O_{\rm actual}(q)$.

\subsection{The exact relative determinant and its mandatory spectral cost}
Taking determinants in \eqref{a03cm:nmr:forms}, with dimension exactly
$4m$, yields at each endpoint
$\log\det T_N=4\log\det D_{a,N-v}^{(s)}-\log\det\mathcal G_N$.
Therefore
\begin{equation}\tag{NMR18}\label{a03cm:nmr:relative}
\mathcal R\log\det T_N=4\mathcal T_{a,s}+F_{\rm ar}.
\end{equation}
Combining the exact finite bounds above gives
\begin{equation}\tag{NMR19}\label{a03cm:nmr:relative-interval}
4\mathcal T_{a,s}\leq\mathcal R\log\det T_N
\leq4\mathcal T_{a,s}
+\min\{-2\log\det(I-S_3),\,8m\log\Omega_a\}.
\end{equation}
On the original LRC finite domains, its correct evaluated expression
for the actual section is
\begin{equation}\tag{NMR20}\label{a03cm:nmr:relative-asymptotic}
\mathcal R\log\det T_N
=64C_\partial aq-512q\log a+F_{\rm ar}+O_{\rm actual}(q).
\end{equation}
The explicit arithmetic term cannot be absorbed into $O(q)$ merely
from the coarse bound \eqref{a03cm:nmr:coarse}.

Let $\mathscr L_{a,s}$ be the larger maximum generalized eigenvalue
on the two original native links $D_2$ relative to $D_0$ and $D_3$
relative to $D_1$, with $D_i=D_{a,N_i-v}^{(s)}$.
Each determinant ratio is at most $\mathscr L_{a,s}^{m}$,
so
\begin{equation}\tag{NMR21}\label{a03cm:nmr:link-cost}
\log\mathscr L_{a,s}\geq\frac{\mathcal T_{a,s}}{2m},\qquad
\liminf_{a\to\infty}\frac{\log\mathscr L_{a,s}}q
\geq\frac{C_\partial}{2}>0 .
\end{equation}
The determinant assertion follows by taking the eigenvalues of
$D_{\rm low}^{-1/2}D_{\rm high}D_{\rm low}^{-1/2}$.
For the limit, substitute \eqref{a03cm:nmr:native-value} and
$m=16a-48-v$, keeping $s=1$ or $s=a$ on its original sequence.

There is a stronger finite result for every uniform comparison of the
two actual metrics at all four endpoints. Such a comparison has
$\alpha_a\mathcal G_N\preceq\mathcal D_N
\preceq\beta'_a\mathcal G_N$ with $0<\alpha_a\leq\beta'_a$.
Every eigenvalue of $T_N$ then lies in
$[\alpha_a,\beta'_a]$, so its four-endpoint determinant return is at
most $8m\log(\beta'_a/\alpha_a)$. By \eqref{a03cm:nmr:relative},
\begin{equation}\tag{NMR22}\label{a03cm:nmr:uniform-cost}
\log\frac{\beta'_a}{\alpha_a}
\geq\frac{4\mathcal T_{a,s}+F_{\rm ar}}{8m}
=\frac{\mathcal T_{a,s}}{2m}+\frac{F_{\rm ar}}{8m},\qquad
\liminf_{a\to\infty}\frac1q\log\frac{\beta'_a}{\alpha_a}
\geq\frac{C_\partial}{2}.
\end{equation}
This applies in particular to the actual optimal constants
$\alpha_a=\min_i\lambda_{\min}(T_{N_i})>0$ and
$\beta'_a=\max_i\lambda_{\max}(T_{N_i})<\infty$.
It thus measures a finite, existing spectral comparison, not an
assumed additional estimate. It strengthens the submitted N9 by
retaining the nonnegative exact arithmetic correction.
Alternatively \eqref{a03cm:nmr:monotonicity} gives on either link
$\mathcal D_{\rm high}\preceq\beta'_a\mathcal G_{\rm high}
\preceq\beta'_a\mathcal G_{\rm low}
\preceq(\beta'_a/\alpha_a)\mathcal D_{\rm low}$,
recovering its original weaker bound \eqref{a03cm:nmr:link-cost}.
Any single fixed invertible coefficient change applied to both
sequences preserves these inequalities and determinant returns by
congruence. No claim is made that every individual eigenvalue has
this growth.

For each original fixed positive boundary form $\Gamma_i$, $i=0,1$,
the exact identity is independently
\begin{equation}\tag{NMR23}\label{a03cm:nmr:fixed-boundary}
\mathcal R\log\det
  (\Gamma_i^{-1/2}\mathcal D_N\Gamma_i^{-1/2})
=4\mathcal T_{a,s}.
\end{equation}
The fixed $\log\det\Gamma_i$ has coefficient
$-1-1+1+1=0$; its individual endpoint value is retained.

\subsection{The consecutive-chart value and the retained original-section correction}
Let $I_0=\{0,\ldots,Q-1\}$ be the proved nonzero consecutive
lower-conductor minor. Its section $L_0$ has the proved image
$\mathbb C[S]_{\leq m-1}$. Denote its arithmetic pullback by
$\mathcal G_N^0$ and its arithmetic return by $F_{\rm ar}^0$.
Use $T_N^0=(\mathcal G_N^0)^{-1/2}
 (D_{0,N}^{(s)})^{\oplus4}(\mathcal G_N^0)^{-1/2}$ for its
native form in that same chart.
The full low-degree localization proved in CAL1--18 gives on its
stated eventual guards
\[
(1-\varepsilon_a)\mathcal G_{q-1}^0
\preceq\mathcal G_N^0\preceq\mathcal G_{q-1}^0,\qquad
\varepsilon_a=(\sqrt{\eta_a}+\sqrt{\zeta_a^*})^2<1 .
\]
Here $\eta_a$ is CAL7 and CAL9 retains the sharper literal bound
$\zeta_a^*=K_{a,m}[2^d(m+d)!/(m!Y^d)]^2\leq K_{a,m}4^{-2d}$,
where $Y=2^{-16}q$, $d=\lfloor Y/16\rfloor$.
These are the proved CAL constants, not an estimate for an
unspecified original minor.
It follows by precisely the determinant argument of
\eqref{a03cm:nmr:Far} that
\begin{equation}\tag{NMR24}\label{a03cm:nmr:consecutive-return}
0\leq F_{\rm ar}^0\leq-8m\log(1-\varepsilon_a),\qquad
\mathcal R\log\det T_N^0
=64C_\partial aq-512q\log a+O_{\rm actual,h}(q).
\end{equation}
Here the exponential estimate for $\varepsilon_a$ is a proved
low-degree estimate on this actual section.

For the original admitted minor $I$, keep the exact chart matrices
$U=M_{I_0\leftarrow I}$ and $D=D_{I_0\leftarrow I}$ from NCB16--28.
They satisfy
$L_I=L_0U+L_CD$ with $L_C=\mathcal LC_a^{\mathsf T}$ and $U$
invertible. Division of the numerator by $E_A$ turns $L_CD$ into
the actual lower exponential evaluation germ. Its whole contribution
is already inside the unchanged native minimization
\eqref{a03cm:nmr:native-min}. Therefore, at every cutoff,
\begin{equation}\tag{NMR25}\label{a03cm:nmr:native-transition}
D_{I,N}=U^*D_{0,N}U,\qquad
L_IU^{-1}=L_0+L_CW,\qquad W=DU^{-1}.
\end{equation}
The native identity follows by the bijective shift of the minimizing
lower coefficient $\eta$ by $Dx$, so it does not assume an isometry
of the arithmetic section.

Since $\det U$ is fixed, the native four-endpoint determinant return
is the same in the original and consecutive charts. The arithmetic
section instead retains the full $L_CW$ term in
\eqref{a03cm:nmr:native-transition}. Its exact arithmetic graph metric,
with both mixed conductor--native blocks, is proved in the companion
original-coefficient image calculation. The direct original-section
formula \eqref{a03cm:nmr:arithmetic-min} already retains those full
polynomial representatives and provides its finite correction
$F_{\rm ar}$ without any chart identification.

Thus the exact identities, native monotonicity, and exponential lower
uniform comparison cost in the submitted Note09 are retained.
Its exponentially small arithmetic correction and the expression
with that correction absorbed into $O(q)$ are established for the
proved consecutive section. For the literal original section the
correct complete formula is \eqref{a03cm:nmr:relative-asymptotic}, with
its explicit finite arithmetic matrices. None of these determinant identities allocates
a scalar to an individual native flag, replaces the independent
degree-$k$ proper-source metric, or adds a term to the canonical action.



% END PRESERVED INPUT 03_COEFFICIENT_METRIC_COMPLETE.tex
\endgroup

\clearpage
\begingroup


\setcounter{equation}{0}
\renewcommand{\theequation}{A03RP.\arabic{equation}}
\section*{Negative reflected trace, both boundary energies and uniform period recovery}
\addcontentsline{toc}{section}{Negative reflected trace, both boundary energies and uniform period recovery}
% BEGIN PRESERVED INPUT 04_RESIDUE_PERIOD_COMPLETE.tex SHA256 cf600e129b35ee865c58173fa0408218fd8c524b16a8558fa1c23aae5e5c8205
% Independent mathematical derivation of incoming 02_RESIDUE_BOUNDARY_VALUE.md.
% Inclusion fragment. The fixed-minor receiving map is proved for every chart.
\begingroup
\section{An explicit reflected trace and both original boundary metrics}
\label{a03rp:sec:RBV}

Keep the original fixed hypothetical simple quartet, its full arithmetic
unit, actual period, logarithm branch and conductor. Write
\[
\begin{gathered}
a\equiv1\pmod4,\quad a\ge17,\quad c=a/2,\\
q=(a+1)^2,\quad q'=(a-7)^2,\quad m=16a-48-v,\quad0\le v\le80,\\
\lambda_{u,b}=c+(2u-a)\delta+i(2b-a)\gamma,\qquad
0<\delta<\tfrac12,\quad\gamma>2,\quad0\le u,b\le a,\\
\chi(S)=\prod_{u,b=0}^a(S-\lambda_{u,b}),\quad
E=\mathbb C[S]/(\chi),\quad
\mathcal P_{q-1}\xrightarrow{\ \sim\ }E .
\end{gathered}\tag{RBV1}
\]
The last map sends an actual polynomial to its residue class, with
the unique degree-below-$q$ remainder as inverse. Roots are distinct
by equality of real and imaginary parts. No packet factors or coordinate
phases are discarded.

\subsection{The explicit polynomial and its exact trace}
Put
\[
t_0=\frac\pi{2\gamma},\qquad r_0=t_0\delta,\qquad
c_0=8\log4-10>0,\qquad
P_a(S)=\sum_{j=0}^{4a}\frac{t_0^{2j}(S-c)^{2j}}{(2j)!}.
\tag{RBV2}
\]
Its degree is $8a\le m-1$ since $8a\ge129\ge49+v$.
Also $8a<q$. Define the original finite reflected form
\[
\mathcal W(f,g)=\sum_{\lambda}\overline{f(a-\overline\lambda)}g(\lambda).
\tag{RBV3}
\]
The reflection permutes the complete grid and is an involution.
Changing its index in the conjugate of $\mathcal W(g,f)$ proves
that this form is Hermitian. The polynomial $P_a$ has real even
coefficients in $S-c$, so
$\overline{P_a(a-\overline\lambda)}=P_a(\lambda)$.
Hence $\mathcal W(P_a,P_a)=\sum_\lambda P_a(\lambda)^2$, a real
number because conjugation also permutes the complete grid.

For $z=\lambda_{u,b}-c$, the actual comparison value is
\[
\cosh(t_0z)=i\sin\!\left(\frac{\pi(2b-a)}2\right)
                   \sinh(r_0(2u-a)).\tag{RBV4}
\]
The sine is $+1$ or $-1$, and $2u-a$ is a nonzero odd integer.
Moreover
$t_0|z|\le at_0\sqrt{\delta^2+\gamma^2}<2a$:
$\delta/\gamma<1/4$ and $\pi<22/7$ give the final strict bound.
For $x\ge0$ and an integer $n\ge1$, the exponential series gives
$\sum_{j\ge n}x^j/j!\le x^ne^x/n!$, because
$(n+j)!\ge n!j!$. The elementary integral comparison
$\log n!=\sum_{j=1}^n\log j\ge\int_1^n\log x\,dx
=n\log n-n+1$ gives $n!\ge(n/e)^n$. Therefore
\[
\begin{split}
|P_a(c+z)-\cosh(t_0z)|
&\le\frac{e^{2a}(2a)^{8a+1}}{(8a+1)!}\\
&\le e^{2a}(e/4)^{8a+1}\le e^{-c_0a}=:\epsilon_a.
\end{split}\tag{RBV5}
\]
This bounds the full omitted even tail by a larger exponential tail.
Impose the explicit eventual guard $\epsilon_a\le r_0/4$.
Since $|\sinh(r_0(2u-a))|\ge r_0$, write the actual value as
$is+e$ with $s$ real and $|e|\le|s|/4$. Direct expansion gives
\[
\Re(is+e)^2\le-s^2+2|s||e|+|e|^2\le-7s^2/16.
\tag{RBV6}
\]
The elementary finite geometric sum gives exactly
\[
\begin{gathered}
\sum_{u=0}^a\cosh(2r_0(2u-a))
=\frac{\sinh(2r_0(a+1))}{\sinh(2r_0)},\\
\sum_{\lambda}\cosh(t_0(\lambda-c))^2
=-\frac{a+1}{2}
 \left[\frac{\sinh(2r_0(a+1))}{\sinh(2r_0)}-(a+1)\right].
\end{gathered}\tag{RBV7}
\]
Consequently the incoming finite sign and its constants are correct:
\[
\boxed{\mathcal W(P_a,P_a)\le-N_a<0,\qquad
N_a=\frac{7(a+1)}{32}
\left[\frac{\sinh(2r_0(a+1))}{\sinh(2r_0)}-(a+1)\right].}
\tag{RBV8}
\]
The bracket is $2\sum_u\sinh^2(r_0(2u-a))>0$.
There is a stronger explicit finite lower bound. For real $s$ and
$0\le\epsilon\le|s|/2$, maximizing $\Re(is+x+iy)^2$ over
$x^2+y^2\le\epsilon^2$ gives $-(|s|-\epsilon)^2$.
For $s>0$, substitute $x^2\le\epsilon^2-y^2$: the remaining
concave quadratic $-s^2-2sy+\epsilon^2-2y^2$ takes its maximum
on $[-\epsilon,\epsilon]$ at $y=-\epsilon$ because its vertex
is $-s/2\le-\epsilon$. The case $s<0$ follows by reflection.
Consequently RBV8 also holds with
\[
\begin{split}
N_a^{\rm sharp}
&=(a+1)\sum_{u=0}^a
 (|\sinh(r_0(2u-a))|-\epsilon_a)^2\\
&=(a+1)\left\{\frac12
 \left[\frac{\sinh(2r_0(a+1))}{\sinh(2r_0)}-(a+1)\right]
-2\epsilon_a\frac{\cosh(r_0(a+1))-1}{\sinh r_0}
 +(a+1)\epsilon_a^2\right\}.
\end{split}\tag{RBV8a}
\]
The additional finite sum follows by summing the positive and
negative odd-index hyperbolic sines separately. In particular the
coefficient $7/32$ may be strengthened to $9/32$ on the same guard.
The original smaller $N_a$ is retained in the downstream displayed
bounds, all of which also hold with $N_a^{\rm sharp}$.
The exact error from RBV7 is at most
\[
2\epsilon_a(a+1)\sum_{u=0}^a|\sinh(r_0(2u-a))|
                         +q\epsilon_a^2.\tag{RBV9}
\]
Since the sum is bounded by $(a+1)e^{r_0a}$, this error divided by
the magnitude in RBV7 is
$O_{\delta,\gamma}(a e^{-(c_0+r_0)a})$.
Keeping also the lower exponential and $(a+1)$ term in RBV7 proves
\[
\mathcal W(P_a,P_a)
=-\frac{(a+1)e^{2r_0(a+1)}}{4\sinh(2r_0)}[1+o(1)].
\tag{RBV10}
\]

\subsection{Original arithmetic density and finite full-form constants}
The source remains
\[
\begin{gathered}
v_h(s)=2\xi(s)/h(s),\quad
w_h(y)=|v_h(1/2+iy)|^2/(2\pi),\quad
d\mu_a(y)=w_h^{*a}(y)\,dy,\\
M_h(\theta)=\int_{\mathbb R}e^{\theta y}w_h(y)\,dy,
\quad\alpha=\pi/2,\quad
d\sigma(y)=|\Gamma(1/4+iy/2)|^2dy/(2\pi),\quad
\sigma(\mathbb R)=\sqrt{2\pi}.
\end{gathered}\tag{RBV11}
\]
RTM1--3/GTM1 proves evenness, positivity almost everywhere and the
finite transforms at $|\theta|\le\alpha$, using the complete unit.
For the simple quartet, the original global constants in GTM1 are
\[
\begin{gathered}
B=50,\quad p=7/2,\quad
\vartheta_h=\int_{-1}^1w_h(y)\,dy,\quad\mathfrak M=M_h(\alpha),\\
b_a=\frac{c_h\vartheta_h^{a-3}e^{-\alpha(a-3)}(a-2)^{-50}}
 {C_\Gamma2^{25}},\qquad
A_a=\frac{C_h^{\rm tilt}}{c_\Gamma}a^{p+1}\mathfrak M^{a-1},\\
b_a(1+y^2)^{-25}d\sigma\le d\mu_a\le A_a d\sigma,\qquad
c_\Gamma=\sqrt2e^{-7/6},\quad C_\Gamma=\sqrt{2\sqrt5}e^{2/3}.
\end{gathered}\tag{RBV12}
\]
The constants $c_h,C_h^{\rm tilt}$ are exactly the original
TW7/BT bounds retained in RTM2, including both half-lines. This proof
does not posit a positive lower bound for the one-factor zeta numerator.

Here is the finite polynomial comparison needed for the whole
boundary bands. In the original orthonormal Gamma frame,
multiplication by $y$ has adjacent entries
$\sqrt{(n+1)(n+1/2)}\le n+1$ (GTM7).
On polynomials of degree at most $D$, its two shifts each have norm
at most $D+1$. Thus
$\|yF\|_\sigma\le2(D+1)\|F\|_\sigma$.
For $F\ne0$, divide the positive integral $|F|^2d\sigma$ by its
own total only to apply the following scalar inequality; the source
measure itself is never changed. Convexity of $(1+x)^{-25}$ on
$x\ge0$, proved by its positive second derivative, gives its tangent
lower bound at the mean of $y^2$. Integration therefore gives
\[
\|F\|_{\mu_a}^2\ge
b_a[1+4(D+1)^2]^{-25}\|F\|_\sigma^2.
\tag{RBV13}
\]
For all degrees through $2q$ we may use the exact constants
\[
L_a^-=b_a[1+4(2q+1)^2]^{-25},\qquad L_a^+=A_a,
\qquad\log(L_a^+/L_a^-)=O_h(a+\log q).\tag{RBV14}
\]
The original Gamma recurrence and mass dictionary are the
Meixner--Pollaczek formulas in DLMF (18.19.8)--(18.19.9), (18.22.8)
\cite{a03rp:human:dlmf18}, with the exact substitution and mass verified
in LT2 and NTG3. This is a full coefficient-form comparison, so it
applies before either original boundary minimization.

\subsection{Concentration of the complete relation bands}
Set $Y=2^{-16}q$ and $R=a\sqrt{\delta^2+\gamma^2}$, and retain the
explicit eventual guard $R\le Y/4$.
For every polynomial $P$ of degree at most $q$ and
$F(S)=\chi(S)P(S)$, the original roots give
\[
|\chi(c+iy)|\le(Y+R)^q\quad(|y|\le Y),\qquad
|\chi(c+iy)|\ge(q-R)^q\quad(q\le y\le2q).
\tag{RBV15}
\]
Use the shifted Legendre basis
$\phi_d(y)=\sqrt{(2d+1)/q}\,L_d(2y/q-3)$ on $[q,2q]$,
$0\le d\le q$. Its Rodrigues formula and squared norm are the
Legendre entries of DLMF (18.5.5), Tables 18.5.1 and 18.3.1
\cite{a03rp:human:dlmf18}. The full finite identity
\[
L_d(z)=2^{-d}\sum_{j=0}^d\binom dj^2(z-1)^{d-j}(z+1)^j
\]
gives $|L_d(z)|\le10^d$ on $|z|\le4$, using
$\sum_j\binom dj^2=\binom{2d}{d}\le4^d$.
For $|y|\le Y$, expansion of the original polynomial
$P(c+iy)$ in this basis and Cauchy--Schwarz prove
\[
|P(c+iy)|^2\le\frac{(q+1)^2}{q}100^q
                         \int_q^{2q}|P(c+it)|^2dt.
\tag{RBV16}
\]
Indeed the basis-square sum is bounded by
$q^{-1}100^q\sum_{d=0}^q(2d+1)=q^{-1}100^q(q+1)^2$.
The lower Gamma density on $[q,2q]$ is
$c_\Gamma e^{-\pi q}(1+2q)^{-1/2}$. Multiplying RBV15--16,
integrating the full central interval using mass $\sqrt{2\pi}$,
and using RBV14 in the denominator proves
\[
\begin{gathered}
\int_{|y|\le Y}|F(c+iy)|^2d\mu_a
 \le\bar\eta_a\|F\|_{\mu_a}^2,\qquad
\bar\eta_a=\min(1,\eta_a),\\
\eta_a=\frac{L_a^+}{L_a^-}\frac{\sqrt{2\pi}}{c_\Gamma}
 \frac{(q+1)^2\sqrt{1+2q}}q e^{\pi q}100^q
 \left(\frac{Y+R}{q-R}\right)^{2q},\\
\log\eta_a\le-c_1q+O_h(a+\log q),\qquad
c_1=30\log2-\pi-\log100>0.
\end{gathered}\tag{RBV17}
\]
The last ratio is at most $2^{-15}$ on the stated guard. This
bound applies to every vector of either complete relation band,
including arbitrary linear combinations and all original cross terms.

\subsection{Both boundary energies, including strict positivity}
Let $e_j=\chi\pi_j/\sqrt{\nu_j}$, where $\pi_j$ is the actual
monic polynomial orthogonal for the full weight
$|\chi(c+iy)|^2d\mu_a$ in the original coordinate. Define
\[
\mathcal B_0=\operatorname{span}(e_0,\ldots,e_{q-1}),\qquad
\mathcal B_1=\operatorname{span}(e_1,\ldots,e_q),\qquad
\Delta_iP=-\operatorname{Proj}_{\mathcal B_i}P.
\tag{RBV18}
\]
These are the original COB two windows. Each band consists of
$\chi$ times polynomials of degree at most $q$, so RBV17 applies
to every unit vector in it.
Their exact link to the original minimum forms follows without a
choice of diagonal trial metric. Let $s_N:E\to\mathcal P_N$ send
a class to its unique minimum-norm representative in $d\mu_a$.
For the original representative $P\in\mathcal P_{q-1}$, its entire
affine fibre is $P+\chi\mathcal P_{N-q}$. Completing its orthogonal
square gives
$s_NP=P-\operatorname{Proj}_{\chi\mathcal P_{N-q}}P$.
Consequently
$\Delta_0=s_{2q-1}-s_{q-1}$ and
$\Delta_1=s_{2q}-s_q$, exactly as in RBV18.
Writing $G_N=s_N^*s_N$ in the unchanged class frame, orthogonality
of these two nested projection differences gives
$G_{q-1}-G_{2q-1}=\Delta_0^*\Delta_0$ and
$G_q-G_{2q}=\Delta_1^*\Delta_1$.
Thus each coefficient matrix is exactly
$\Gamma_{i,I}=L_I^*\Delta_i^*\Delta_iL_I$ on its stated original
residue injection. This proves the map into both boundary energies,
including all cross terms of the actual minimum.
These are precisely COB.5--6 and COB.13 in
\path{COUPLED_ORIGINAL_CONE_BOUNDARY.tex}. The general positive
coefficient-square operation is the Schur minimum in Boyd and
Vandenberghe \cite[Appendix A.5.5, p.~650]{human:boyd-vandenberghe};
the preceding formulas prove its full original-fibre specialization.

In fact both $\Delta_i$ are injective on the entire original
$\mathcal P_{q-1}$. Pair the roots with indices $u$ and $a-u$.
Since $(a+1)^2/2$ is even, their original factors give exactly
\[
\chi(c+iy)=\prod_{b=0}^a\prod_{0\le u<a/2}
\left[(y-(2b-a)\gamma)^2+(2u-a)^2\delta^2\right]>0.
\tag{RBV19}
\]
If $P\perp\mathcal B_0$, test its orthogonality against
$\chi P\in\mathcal B_0$ to obtain
$\int\chi(c+iy)|P(c+iy)|^2d\mu_a=0$, hence $P=0$.
If $P\perp\mathcal B_1$, choose the unique $b$ such that
$F=P-b\chi$ is orthogonal to $e_0$. Since $\chi$ is a multiple
of $e_0$, $F$ remains orthogonal to $e_1,\ldots,e_q$.
Its degree is at most $q$, so $\chi F$ belongs to
$\operatorname{span}(e_0,\ldots,e_q)$. Testing against it and using
RBV19 gives $F=0$. Then $P=b\chi$ and $\deg P<q$ force $P=b=0$.
This proves strict positivity without an assumed missing metric bound.

On the original source line,
$|P_a(c+iy)|\le\cosh(t_0|y|)\le e^{t_0|y|}$.
The exact convolution transform, with all masses retained, is
\[
\int e^{\theta y}d\mu_a=M_h(\theta)^a.
\]
Evenness and $2t_0<\alpha$ therefore give, with
$\theta_1=(\alpha-2t_0)/2>0$, $\theta_2=2t_0+\theta_1<\alpha$,
\[
\|P_a\|_{\mu_a}^2\le2M_h(2t_0)^a,\qquad
\|1_{|y|>Y}P_a\|_{\mu_a}^2
 \le2e^{-\theta_1Y}M_h(\theta_2)^a.
\tag{RBV20}
\]
For a unit $f\in\mathcal B_i$, split its inner product with $P_a$
over the two original integration regions. Cauchy--Schwarz and RBV17
give
$|\langle f,P_a\rangle|\le\sqrt{\bar\eta_a}\|P_a\|+
\|1_{|y|>Y}P_a\|$. Taking the supremum over the full unit sphere
of the band, then $(x+y)^2\le2x^2+2y^2$, proves
\[
\boxed{0<\|\Delta_iP_a\|_{\mu_a}^2\le B_a,
\quad B_a=4\bar\eta_aM_h(2t_0)^a+
4e^{-\theta_1 2^{-16}q}M_h(\theta_2)^a,\quad i=0,1.}
\tag{RBV21}
\]
No band-dimension multiplier enters this estimate.

\subsection{Relative eigenvalue and evaluated discrepancy}
Let $G_i^E$ be the Gram of $\Delta_i$ on the actual monomial frame
of $\mathcal P_{q-1}$; it is strictly positive by RBV19.
Let $W^E$ be the matrix of RBV3 in that same frame. Then
\[
T_i^E=(G_i^E)^{-1/2}W^E(G_i^E)^{-1/2}
\]
is Hermitian. The coefficient column of $P_a$ gives the actual
Rayleigh quotient, so
\[
\lambda_{\min}(T_i^E)\le-N_a/B_a<0,\qquad
\|I-T_i^E\|\ge1+N_a/B_a.\tag{RBV22}
\]
This follows by diagonalizing the Hermitian matrix and expressing
the Rayleigh quotient as a positive weighted average of its real
eigenvalues. With the original two matrices, the discrepancy is
\[
\begin{split}
\|\Delta_iP_a\|_{\mu_a}^2-\mathcal W(P_a,P_a)
&=\frac{(a+1)e^{2r_0(a+1)}}{4\sinh(2r_0)}[1+o_h(1)],\\
\lim_{a\to\infty}\frac{\|\Delta_iP_a\|^2-\mathcal W(P_a,P_a)}
 {a e^{\pi\delta a/\gamma}}
&=\frac{e^{\pi\delta/\gamma}}{4\sinh(\pi\delta/\gamma)}.
\end{split}\tag{RBV23}
\]
Indeed RBV21 has logarithm at most
$-c_*q+O_h(a+\log q)$, where
\[
c_*=\min\{30\log2-\pi-\log100,
                   2^{-16}(\pi/2-\pi/\gamma)/2\}>0.
\tag{RBV24}
\]
This is negligible compared with RBV10. Also
$\log N_a=2r_0a+\log a+O_h(1)$, so RBV22 proves
\[
\liminf_{a\to\infty}q^{-1}\log[-\lambda_{\min}(T_i^E)]\ge c_*.
\tag{RBV25}
\]
All limits run over the original $a\equiv1\pmod4$ family at fixed
actual quartet/unit and eventually satisfy the explicitly stated
guards. No actual-period numerical eigenvalue is invented.

\subsection{Exact return through every original conductor chart}
This paragraph records the precise map needed to receive these
results on the supplied COB coefficient image, without identifying
two different sections of the same quotient. Write
$\chi(S)=\sum_{j=0}^qc_jS^j$, $c_q=1$, and define
\[
\mathfrak p_n(S)=\sum_{j=n+1}^qc_jS^{j-n-1},\qquad
\mathbb V_{n,\lambda}=\lambda^n,\qquad
\mathcal L\nu=\sum_\lambda\nu_\lambda\frac{\chi(S)}{S-\lambda},
\quad F=\mathcal L\mathbb V^{-1}.
\tag{RBV26}
\]
NCB5--7 proves $Fe_n=\mathfrak p_n$ by the coefficient of $S^{-1}$
in $\mathfrak p_nS^j/\chi$. Thus $F$ is invertible and its inverse is
the exact Laurent map
$(F^{-1}P)_n=[S^{-n-1}]P(S)/\chi(S)$, $0\le n<q$.
Let $C^{\mathsf T}$ be the original conductor multiplication in
exponential coefficients, and write
$\mathbb VC^{\mathsf T}=P_vJ$ with the original $v$ leading zeros.
For every retained invertible $q'$-row minor $J_I$, define
\[
\begin{gathered}
H_I=J_I^{-1}E_I^{\mathsf T},\quad
Q_I=E_{I^c}^{\mathsf T}-J_{I^c}H_I,\quad
L_I=FP_vE_{I^c},\quad C_E=\mathcal LC^{\mathsf T},\\
\eta_I=H_IP_v^{\mathsf T}F^{-1}P_a,\quad
z_I=Q_IP_v^{\mathsf T}F^{-1}P_a,\qquad
\boxed{P_a=C_E\eta_I+L_Iz_I.}
\end{gathered}\tag{RBV27}
\]
There is no low component because $\deg P_a<m$ makes all Laurent
coordinates below $q-m\ge v$ zero. To prove the displayed identity,
check the $I$ and $I^c$ rows to get
$JH_I+E_{I^c}Q_I=I$. Multiply by $FP_v$ and use
$FP_vJ=C_E$. These are the exact original residue and conductor maps;
no coefficient section is inserted in place of another.

In particular the consecutive minor
$I_0=\{0,\ldots,q'-1\}$ has determinant
\[
\det J_{I_0}=\mu_v^{q'}
 \prod_{r=0}^{q'-1}\binom{v+r}{v}\det\mathbb V_{\rm lower}\ne0.
\tag{RBV28}
\]
Indeed its row $r$ is the lower-root evaluation of a degree-$r$
polynomial with leading coefficient $\binom{v+r}{v}\mu_v$;
triangular coefficient elimination gives exactly this determinant.
Here $\eta_{I_0}=0$ and
$\operatorname{im}L_{I_0}=\mathcal P_{m-1}$ by the distinct monic
degrees $m-1,m-2,\ldots,0$ of its columns. The literal inverse is
\[
(z_{I_0})_j=[S^{-(q-m+j+1)}]P_a(S)/\chi(S).
\tag{RBV29}
\]
Therefore RBV8--25 apply to the Gram
$\Gamma_{i,I_0}=L_{I_0}^*G_i^EL_{I_0}$ and reflected matrix
$L_{I_0}^*W^EL_{I_0}$ with precisely the same bounds, because the
Rayleigh vector is the explicit column RBV29. For any actual other
chart $I$, RBV27 supplies the full original coefficient vector
$(\eta_I,z_I)$ instead. With $M_I=[C_E,L_I]$, the same calculation
holds on its full image with the two exact matrices
\[
M_I^*G_i^EM_I=\begin{pmatrix}
C_E^*G_i^EC_E&C_E^*G_i^EL_I\\
L_I^*G_i^EC_E&L_I^*G_i^EL_I
\end{pmatrix},\qquad M_I^*W^EM_I.
\tag{RBV30}
\]
The columns of $M_I$ form the original direct conductor/native
decomposition in the high-moment space, hence this Gram is positive
definite. All cross terms in RBV30 are essential.
This gives a proved negative class and exponential relative bound
through every complete original chart. Selecting only the native
block in a nonconsecutive chart is a further, different restriction:
RBV27 evaluates its omitted conductor component exactly instead of
declaring it zero. The companion original-image derivation identifies
which selector is the specified current COB map; no root ordering
alone is used here as a proof of an ordering of row subsets.
For the original four-copy maps use $I_4\otimes L_{I_0}$ or
$I_4\otimes M_I$ and the coefficient column with the displayed
one-copy preimage in its first block and zero in the other three.
Its reflected value and each boundary energy are exactly the
one-copy values. Hence RBV22--25 have the same constants on these
original four-copy images, with no additional rank factor.

\subsection{The complete tensor residue, its multiplier, and the residual}
We now expand the full FTR/FRT comparison on the present original
simple-quartet domain. For the four original roots
$z_{\epsilon\eta}=1/2+(2\epsilon-1)\delta+i(2\eta-1)\gamma$,
retain every original unit value in
\[
w_{\epsilon\eta}=\frac{v_h(z_{\epsilon\eta})}{h'(z_{\epsilon\eta})}\ne0,
\quad p(X,Y)=\sum_{\epsilon,\eta=0}^1w_{\epsilon\eta}X^\epsilon Y^\eta,
\quad c_{u,b}^{\rm ten}=[X^uY^b]p(X,Y)^a.
\tag{RBV31}
\]
There are no omitted jets: the one-factor multiplicity is exactly
one here. The original full-unit residue formula FTR.8, by taking
the simple residues in every ordered tensor factor, gives
\[
\mathcal B(f,g)=\sum_{u,b=0}^a
c_{u,b}^{\rm ten} f^{\sharp_a}(\lambda_{u,b})g(\lambda_{u,b}).
\tag{RBV32}
\]
Indeed each ordered root tuple contributes the product of its
$a$ unchanged $w_{\epsilon\eta}$ factors; the exponents of $X,Y$
record the two original plus counts. Expansion of $p^a$ counts
every ordered tuple once. Thus the tensor pullback multiplier in
the cyclic residue form is the actual class with values
$b_{\rm ten}(\lambda_{u,b})=\chi'(\lambda_{u,b})c_{u,b}^{\rm ten}$.
This is the complete simple-root specialization of FTR.6--12.

Put $\mathcal D=\{(u,b):c_{u,b}^{\rm ten}\ne0\}$ and
$\mathcal Z=\{0,\ldots,a\}^2\setminus\mathcal D$, and retain the
original idempotents
$e_{u,b}=\chi(S)/((S-\lambda_{u,b})\chi'(\lambda_{u,b}))$.
The trace multiplier and complementary form are exactly
\[
\begin{gathered}
a_{\rm tr}=\sum_{(u,b)\in\mathcal D}
              (c_{u,b}^{\rm ten})^{-1}e_{u,b},\qquad
\mathcal W_{\mathcal Z}(f,g)=\sum_{(u,b)\in\mathcal Z}
              f^{\sharp_a}(\lambda_{u,b})g(\lambda_{u,b}),\\
b_{\rm ten}a_{\rm tr}=e_{\mathcal D}\chi'\quad\text{in }E,
\qquad
\boxed{\mathcal W(f,g)=\mathcal B(f,a_{\rm tr}g)
                              +\mathcal W_{\mathcal Z}(f,g).}
\end{gathered}\tag{RBV33}
\]
All equations follow by evaluation at every simple root; the
evaluation map is an isomorphism, so they prove the typed algebra
map, not just matching dimensions. They are exactly FRT.2--5.
The trace equality can also be verified by taking the residue of
$\chi'f^{\sharp_a}g/\chi$ at each root. The classical multiplication
trace construction is recorded by the Stacks Project Authors,
Section 49.3, tags 0BVH and 0BJF \cite{a03rp:human:stacks-trace}; the
displayed root computations prove the full specialization here.

In particular the retained residual has the exact finite value
\[
\mathcal W_{\mathcal Z}(P_a,P_a)
 =\sum_{(u,b)\in\mathcal Z}P_a(\lambda_{u,b})^2.
\tag{RBV34}
\]
Conjugation of the full unit and $h$ gives
$w_{\epsilon,1-\eta}=\overline{w_{\epsilon\eta}}$;
their original reflection gives
$w_{1-\epsilon,\eta}=-\overline{w_{\epsilon\eta}}$.
Since $a$ is odd, the coefficient sets $\mathcal D,\mathcal Z$
are closed under both grid symmetries. Thus both partial sums of
$P_a(\lambda)^2$ are real. RBV6 proves that RBV34 is nonpositive,
and strictly negative whenever $\mathcal Z$ is nonempty. This
does not determine that set by a rank assumption: its exact definition
retains all cancellations in the actual coefficients RBV31.

There is also a negative pairing directly in the retained tensor
residue without needing to determine its interior zero coefficients.
The four edge coefficients of $p^a$ are binomial products of
nonzero original $w$'s, so every perimeter point belongs to
$\mathcal D$. In particular its two complete edges $u=0,a$
contain $2(a+1)$ points with absolute comparison value
$\sinh(r_0a)$. Applying RBV8a at these points and RBV6 at all
remaining points of $\mathcal D$ proves
\[
\boxed{\mathcal B(P_a,a_{\rm tr}P_a)
 \le-2(a+1)(\sinh(r_0a)-\epsilon_a)^2<0.}\tag{RBV35}
\]
The second argument is the explicitly computed full polynomial
class $a_{\rm tr}P_a\bmod\chi$; it is not identified with a native
coefficient vector unless its Laurent coordinates satisfy that
specific section's membership test.

Finally the original arithmetic injection is
$\mathcal I:E\hookrightarrow Q^{\otimes a}$ with the complete
$[v_h]_h^{\otimes a}$ factor. Its inverse on its named image is
the original full Mellin-jet observation, division by that invertible
Taylor factor and inverse sum substitution. Thus
$\mathcal I M_{a_{\rm tr}}\mathcal I^{-1}$ is the exact arithmetic
map realizing RBV33--35 on $\mathcal I(e_{\mathcal D}E)$.
In RBV35 the first argument may be written $e_{\mathcal D}P_a$,
since the complementary first-input block pairs to zero under
$\mathcal B$ by RBV32; the separate value RBV34 remains in
the full reflected trace.
It commutes with the original sum action because both act by
multiplication before this injection. The proper boundary image
uses instead the full COB.8 map
$\Psi_2:[P]\mapsto[P(D_1+\cdots+D_a)F_h^{\otimes a}]$
on $\mathbb C[S]/\chi^{a+1}$, whose one-factor Taylor unit is
$[v_h]_{h^2}$; these two different original target spaces are not
silently identified. The ordinary class is recovered on the named
boundary image by $\mathcal I\Delta_i^{-1}\Psi_2^{-1}$ exactly
as in COB.12. This retains the full tensor-residue residual and
every original unit. The finite packet trace is not substituted
for the full infinite Weil distribution, and none of these values
is assigned to the native minimum metric.

The tensor calculation also has a relative-metric receiver, with
the residual retained at matrix level. On any of the certified
images above let $L$ denote its displayed injection and set
$\Gamma_i=L^*G_i^EL$. With $W_{\mathcal D}^E,W_{\mathcal Z}^E$
the two explicitly computed reflected matrices, define
\[
\begin{gathered}
T_{i,\mathcal D}=\Gamma_i^{-1/2}L^*W_{\mathcal D}^EL\Gamma_i^{-1/2},
\qquad T_{i,\mathcal Z}=\Gamma_i^{-1/2}L^*W_{\mathcal Z}^EL\Gamma_i^{-1/2},\\
T_i=T_{i,\mathcal D}+T_{i,\mathcal Z},\qquad
\boxed{\lambda_{\min}(T_{i,\mathcal D})
\le-\frac{2(a+1)(\sinh(r_0a)-\epsilon_a)^2}{B_a}<0.}
\end{gathered}\tag{RBV36}
\]
Both partial matrices are Hermitian by their proved symmetry sets.
Use the very same preimage of $P_a$ as in RBV22 and the strict
boundary bound RBV21 to prove the displayed inequality. Its
negative magnitude again has logarithmic lower rate $c_*$ after
division by $q$. This is a bound for the actual tensor-residue
comparison $\mathcal B(f,a_{\rm tr}g)$; the other matrix remains
in the exact displayed sum. No equality of the smallest eigenvalues
of those three matrices is asserted.
\endgroup


% Complete proof fragment. Bibliography entries are in REFERENCES.tex.
\begingroup
\section{Uniform recovery on the original fixed-period observation image}
\label{a03rp:sec:UPR}

This calculation derives U1--U8 of the supplied
\emph{A uniform inverse bound for the original period observation},
19 September 2026.  It corrects one finite-module degree statement and
makes the linear division operator and its degree-uniform bound explicit.
The actual observation, Fourier frame and inherited Gram are AKS7--10
in the existing cumulative source \cite{a03rp:programme:aks-current}.
Every period constant below is at the same fixed original nonzero
admitted period and its original logarithm branch.  No numerical
certificate for that period is claimed.

\subsection{The unchanged coefficient-dual map and its exact kernel}
Retain a simple stipulated quartet
\[
\tfrac12\pm\delta\pm i\gamma,\qquad
0<\delta<\tfrac12,\quad\gamma>2,
\]
the original arithmetic unit \(U\), primary idempotent matrix
\(\mathcal E_{\rm prim}\), and period matrix \(\Pi\).  The precise
one-factor matrix and Fourier data are
\[
F_{jt}=\zeta^{jt}-1,\quad
(F^{-1})_{tj}=\zeta^{-tj}/5,\quad
\zeta=e^{2\pi i/5},\quad
H=F^{-1}\Pi U\mathcal E_{\rm prim},\quad
D=\operatorname{diag}(\zeta,\zeta^2,\zeta^3,\zeta^4).
\tag{UPR1}
\]
The original provider establishes invertibility of \(H\).  In
\(\mathcal A=\mathbb C[x_1,x_2,x_3,x_4]\), put
\[
Q(x)=Q_0(H^{-1}x),\qquad Q_0(t)=t_1t_4-t_2t_3.
\]
For degree \(k\ge0\) let
\(\mathcal B_k=\operatorname{span}\{Z^aW^b:0\le a,b\le k\}\).
The original coefficient-dual map is
\[
\Theta_k:\mathcal A_k\longrightarrow\mathcal B_k,\qquad
f(x)\longmapsto f\bigl(H(ZW,Z,W,1)^{\mathsf T}\bigr).
\tag{UPR2}
\]
The arithmetic receiver below uses the original subsequence
\(k\equiv1\pmod4,\ k\ge5\); these finite polynomial constructions
also define the intermediate graded maps at every degree.

Here is an exact proof of its kernel and an explicit full lift.
For each pair \(0\le a,b\le k\), put \(n=\max(0,a+b-k)\), and define
\[
\ell_{a,b}(t)=
t_1^nt_2^{a-n}t_3^{b-n}t_4^{k-a-b+n}.
\tag{UPR3}
\]
All exponents are nonnegative, their sum is \(k\), and substitution
\(t=(ZW,Z,W,1)\) gives \(Z^aW^b\).  Every homogeneous monomial in
the \(t_i\) can be reduced by \(t_1t_4=t_2t_3\) until either its
first or its fourth exponent is zero.  The two exponents
\(a=n_1+n_2\), \(b=n_1+n_3\), and its total degree \(k\) then
determine all four exponents uniquely, yielding (UPR3).
Thus the remaining monomials have linearly independent images, and
\[
\ker\Theta_k=Q\mathcal A_{k-2},\qquad
\Theta_k\mathcal A_k=\mathcal B_k,
\tag{UPR4}
\]
where a negative homogeneous degree means zero.  This argument
also proves the complete relation ideal; it is not a dimension-only
identification.

Define the linear lift on the displayed biform frame by
\(L_k(Z^aW^b)=\ell_{a,b}(H^{-1}x)\).
For the ordinary coefficient \(\ell^1\) norm put
\[
h_-=\max\left(1,\max_i\sum_j|(H^{-1})_{ij}|\right).
\]
The sum of absolute coefficients of a product is at most the
product of those sums.  Apply this to the \(k\) linear factors of
(UPR3), and then to a linear combination of biform monomials:
\[
\Theta_kL_k=I_{\mathcal B_k},\qquad
\|L_kf\|_1\le h_-^k\|f\|_1.
\tag{UPR5}
\]
The full arithmetic unit and period remain in \(H\) and \(h_-\).

\subsection{A fixed homogeneous module with the corrected degree bound}
Put
\[
\mathcal S=\mathbb C[X_1,X_2,X_3,X_4],\quad X_i=x_i^5,\qquad
\mathcal R=\{0,1,2,3,4\}^4.
\]
The exact decomposition
\[
\mathcal A=\bigoplus_{r\in\mathcal R}\mathcal S x^r
\tag{UPR6}
\]
follows by Euclidean division of each monomial exponent by five.
The invariant submodule for the original \(D\)-action is
\[
\mathcal A^D=\bigoplus_{r\in\mathcal R_0}\mathcal S x^r,\qquad
\mathcal R_0=\{r:\textstyle\sum_{i=1}^4 i r_i=0\bmod5\}.
\tag{UPR7}
\]
There are \(625\) residue vectors and \(125\) invariant residue
vectors: the fourth residue is uniquely determined by the first
three in the displayed congruence.

Use the original grading \(\deg X_i=5\) and \(\deg x^r=|r|\).
For the second source summand below the residue basis has the
additional degree-two shift:
\[
\mathcal M:
\left(\bigoplus_{r\in\mathcal R_0}\mathcal Sx^r\right)
\oplus
\left(\bigoplus_{r\in\mathcal R}\mathcal Sx^r[2]\right)
\longrightarrow
\bigoplus_{r\in\mathcal R}\mathcal Sx^r,\qquad
(f,g)\longmapsto f+Qg.
\tag{UPR8}
\]
It is an exact degree-preserving map, and
\[
\operatorname{im}\mathcal M=\mathcal A^D+Q\mathcal A.
\tag{UPR9}
\]

The entries of the multiplication-by-\(Q\) block have total
\(X\)-degree at most \emph{two}, rather than the at-most-one assertion
in the incoming U2.  Indeed multiplication of \(x^r\) by a quadratic
monomial can cause at most two carries in division of exponents
by five.  For distinct \(i,j\), the term \(x_ix_j\) and a residue
with \(r_i=r_j=4\) produce the entry \(X_iX_j\).  The actual
grading and the fact that this is one fixed finite polynomial
matrix are unchanged by this correction.

\subsection{Finite basis, homogeneous lifts and a specified linear division}
We use Buchberger's polynomial-basis method
\cite{a03rp:human:buchberger2006}, with its finiteness supplied by
Dickson's lemma \cite{a03rp:human:dickson1913}.  The needed module argument
and norm estimates are supplied here.

Enumerate the \(625\) target residues once in lexicographic order,
with integer position \(b_r\in\{0,\ldots,624\}\).  Order target
monomials \(X^\alpha x^r\) lexicographically by
\((\alpha_1,\alpha_2,\alpha_3,\alpha_4,b_r)\).
This is a well-order compatible with multiplication by every
\(X^\beta\).

Start with the nonzero columns of (UPR8), retaining their corresponding
source coordinate vectors.  Make each column monic by division by
its actual nonzero leading coefficient.  For pairs whose leading
target position agrees, form the usual least-common-monomial
cancellation, reduce it using the current monic list, and append
each nonzero remainder with its simultaneously updated source
lift.  Order the finite pair lists lexicographically to specify
the choices.  A reduction of a homogeneous vector by a homogeneous
column retains its full shifted degree, because the cancelled
monomials occupy the same target position.  Every resulting column
\(g_j\) and retained lift \(h_j\) are therefore homogeneous and
\[
\mathcal M h_j=g_j.
\tag{UPR10}
\]

For completeness, an infinite sequence in \(\mathbb N^d\) has an
earlier term componentwise at most a later one.  For \(d=1\), either
some value occurs infinitely often, or an infinite strictly
increasing subsequence can be selected.  In either case there is
an infinite nondecreasing subsequence.  Apply this to the first
coordinate, and apply the induction hypothesis to the other
coordinates of that subsequence.  This proves the assertion for
every finite \(d\), including \(d=4\).
An infinite sequence of appended leading monomials would contain
infinitely many in one target position and then an earlier one
dividing a later one.  That contradicts the complete remainder
reduction before appending.  Hence only finitely many columns are
appended, and the pair algorithm terminates.

The output is a module Gr\"obner basis.  To verify this criterion
directly, take a nonzero element of \(\operatorname{im}\mathcal M\)
and a representation by monomial multiples of the output columns
whose largest occurring leading monomial is least possible; among
these, minimize the number of terms at that largest monomial.
If it exceeds the leading monomial of the represented vector,
at least two terms at the same target position cancel.
Their least-common-monomial combination is a multiple of one of
the processed pairs.  That pair has a representation with all
occurring leading monomials strictly smaller than its cancelled
monomial, because its remainder is zero.  Substitute this
representation and combine the remaining terms at the largest
monomial.  This decreases their number, or the largest monomial
itself, a contradiction.  Thus some leading basis monomial divides
the leading monomial of the given element.  Division consequently
has zero remainder exactly on the image module.
A reduced basis may be obtained by discarding divisible leading
monomials and reducing the tails, updating their lifts throughout;
reducedness is not needed in the subsequent estimate.

We now specify positive integer scores that strictly decrease in
every rewrite.  Let \(D_*\) be the maximum of all absolute exponent
differences, coordinate by coordinate, between a leading term and
a tail term in the finite list \(g_j\); set \(D_*=0\) when all tails
are empty.  Set
\[
T_*=2(D_*+625)+2,\quad
w_i=T_*^{\,5-i}\quad(1\le i\le4),\quad
W_*=\max_i w_i,\quad B_*=624.
\tag{UPR11}
\]
Give \(X^\alpha x^r\) score
\(\operatorname{sc}(\alpha,r)=\sum_iw_i\alpha_i+b_r\).
If a leading and tail term first differ at variable index \(i\),
their score difference is at least
\[
T_*^{5-i}
-D_*\sum_{j>i}T_*^{5-j}-624>0.
\]
The last inequality follows from
\(D_*/(T_*-1)<1/2\) and \(T_*/2>624\).
If all four exponents agree, the leading position has strictly
larger \(b_r\).  Thus every such difference is an integer at least
one.  Multiplication by a monomial shifts both scores equally.
All scores are nonnegative.

Set the following fixed constants from the finite basis and its
actual lifts:
\[
C_*=\max\left(2,1+\max_j\|\operatorname{tail}(g_j)\|_1\right),
\qquad H_*=\max\left(1,\max_j\|h_j\|_1\right).
\tag{UPR12}
\]
Here and below coefficient norms in the free module use its
displayed residue basis; under (UPR6) these are precisely the
ordinary polynomial coefficient norms.

Define a linear division operator by a recursion on individual
monomials, rather than by choices depending on a whole input
polynomial.  If no leading \(g_j\) divides a target monomial \(t\),
set \(\mathcal R(t)=0\), \(\operatorname{NF}(t)=t\).
Otherwise use the first such \(g_j\), put \(m=t/\operatorname{LM}g_j\),
and define
\[
\begin{split}
\mathcal R(t)&=m h_j-\mathcal R(m\,\operatorname{tail}g_j),\\
\operatorname{NF}(t)&=-\operatorname{NF}
                         (m\,\operatorname{tail}g_j).
\end{split}
\tag{UPR13}
\]
Extend each right side linearly over its finitely many monomials.
All recursive arguments have lower score, so this is a finite
definition.  Finally extend to arbitrary polynomials by their
monomial expansion.  The resulting maps are linear, and induction
in the recursion proves
\[
f=\mathcal M\mathcal R(f)+\operatorname{NF}(f).
\tag{UPR14}
\]
The normal form has only irreducible monomials; if \(f\) is in
the image module, (UPR14) and the just-proved leading-monomial
criterion imply that it is zero.  No choice of a degree-dependent
minor is involved.

If \(t\) has score \(s\), another induction proves
\[
\|\mathcal R(t)\|_1\le H_*C_*^{s+1}.
\tag{UPR15}
\]
Indeed the first lift has norm at most \(H_*\) and the total
absolute coefficient mass of its tails is at most \(C_*-1\).
Their scores are at most \(s-1\), so the induction gives the bound
\[
H_*+(C_*-1)H_*C_*^s\le H_*C_*^{s+1}.
\]
At score zero a reducible term has no tail, and the same inequality
holds directly.  Monomial multiplication does not change coefficient
norms.  Homogeneity of (UPR10) shows that this recursion preserves
the original shifted degree throughout.

On degree \(k\), \(5|\alpha|+|r|=k\), so the score is at most
\(W_*k/5+B_*\).  By linearity, for every degree-\(k\) element of
the actual image module,
\[
\mathcal M\mathcal R(f)=f,\qquad
\|\mathcal R(f)\|_1
\le H_*C_*^{W_*k/5+B_*+1}\|f\|_1.
\tag{UPR16}
\]
These constants are defined by one finite calculation over the
fixed field of actual coefficients.  The proof establishes their
finiteness and the uniform bound.  It does not claim that this
basis has been numerically computed for an unspecified hypothetical
quartet or that exact zero tests on arbitrary supplied complex
numbers are available to a numerical implementation.

\subsection{A linear right inverse on the original invariant image}
Let \(f\in\Theta_k(\mathcal A_k^D)\).  Its fixed lift \(L_kf\)
differs from an invariant lift by a multiple of \(Q\), by (UPR4).
Thus \(L_kf\in\operatorname{im}\mathcal M\).
Write \(\mathcal R(L_kf)=(p,g)\) in (UPR8).  Then
\[
p\in\mathcal A_k^D,\qquad L_kf=p+Qg,\qquad \Theta_kp=f.
\]
The map \(f\mapsto p\) is linear by (UPR5) and (UPR13).
Since \(\dim\mathcal B_k=(k+1)^2\),
\[
\|p\|_2\le\|p\|_1
\le\mathscr R_k\|f\|_2,\qquad
\mathscr R_k=
H_*C_*^{W_*k/5+B_*+1}h_-^k(k+1).
\tag{UPR17}
\]
In particular \(\log\mathscr R_k=O_{H}(k+\log(k+1))\).
Here \(k\) is the original tensor degree; the coefficient module
and all its constants are fixed.

To compare with the actual matrix, expand
\(\prod_t(\sum_\alpha H_{t\alpha}t_\alpha)^{\nu_t}\)
for each invariant \(\nu\).
The coefficient with plus-real count \(a\), plus-imaginary count
\(b\), is
\[
B_{\nu;(a,b)}
=\sum_{\substack{\sum_\alpha n_{t\alpha}=\nu_t\\
                  \text{original plus counts }a,b}}
\left(\prod_t\frac{\nu_t!}{\prod_\alpha n_{t\alpha}!}\right)
\prod_{t,\alpha}H_{t\alpha}^{\,n_{t\alpha}},
\tag{UPR18}
\]
exactly AKS8.  Thus the matrix of the invariant restriction of
(UPR2) is \(B_k^{\mathsf T}\) in these unchanged coefficient frames.
For any map \(T\), its least-norm preimage has norm at most the
preimage supplied by a right inverse.  Equation (UPR17) therefore
bounds the Moore--Penrose inverse of \(B_k^{\mathsf T}\) by
\(\mathscr R_k\).  A singular-value factorization
\(B_k=U\Sigma V^*\) gives
\(B_k^{\mathsf T}=\bar V\Sigma(\bar U)^*\), with exactly the
same nonzero singular values.  Consequently
\[
\|(B_k)^+\|\le\mathscr R_k.
\tag{UPR19}
\]
Zero singular values are kept as the actual observation kernel.
If its whole image is zero, the pseudoinverse is zero and the
same assertion holds.

For comparison the forward norm is also degree-exponential.
With \(h_+=\max(1,\max_i\sum_j|H_{ij}|)\), expansion of the same
linear factors has coefficient \(\ell^1\) norm at most \(h_+^k\).
Substitution \(t=(ZW,Z,W,1)\) only combines coefficients and cannot
increase that norm.  It follows that
\[
\|B_k\|=\|B_k^{\mathsf T}\|
\le h_+^k\sqrt{\binom{k+3}{3}}.
\tag{UPR20}
\]
Hence the condition number on its nonzero coefficient singular
values is \(\exp(O_H(k+\log(k+1)))\).

\subsection{The exact original word Gram, with each multiplicity once}
Direct summation of fifth roots of unity gives
\[
(F^*F)_{tu}
=\sum_{j=1}^4(\zeta^{-jt}-1)(\zeta^{ju}-1)
=5(\delta_{tu}+1).
\tag{UPR21}
\]
Thus \(F^*F=5(I+\mathbf1\mathbf1^*)\), whose eigenvalues are
five with multiplicity three and twenty-five with multiplicity one.

For \(|\nu|=k\), let \(e_\nu\) be the sum of all distinct ordered
standard tensor words with multiplicity vector \(\nu\).
These are the unscaled symmetric-word vectors of the original
provider.  Different \(e_\nu\) are orthogonal and
\[
\|e_\nu\|^2=k!/\nu!,\qquad
1\le k!/\nu!\le4^k.
\tag{UPR22}
\]
The upper bound follows because the sum of all these multinomial
integers is \(4^k\).
If \(\mathcal F_0\) has columns \(F^{\otimes k}e_\nu\) for the
original invariant index set, then its inherited Gram entries are
\[
(\mathcal F_0^*\mathcal F_0)_{\nu\mu}
=\sum_{\substack{\sum_u n_{tu}=\nu_t\\
                  \sum_t n_{tu}=\mu_u}}
\frac{k!}{\prod_{t,u}n_{tu}!}
\prod_{t,u}[5(\delta_{tu}+1)]^{n_{tu}}.
\tag{UPR23}
\]
Indeed expansion of
\((\bar x^{\mathsf T}F^*Fy)^k\) gives the displayed coefficients;
equivalently expand both ordered tensor-word sums.
This is the exact AKS10 formula, including every row cross term.

The squared singular values of \(F^{\otimes k}\) lie between
\(5^k\) and \(25^k\), since tensor eigenvalues multiply.
For an arbitrary coefficient column \(a=(a_\nu)\), (UPR22) gives
\[
\sum_\nu|a_\nu|^2
\le\left\|\sum_\nu a_\nu e_\nu\right\|^2
\le4^k\sum_\nu|a_\nu|^2.
\]
Apply the tensor singular-value bounds to this same vector,
restricted to the actual invariant indices:
\[
\boxed{5^kI\le\mathcal G_0:=\mathcal F_0^*\mathcal F_0
\le100^kI.}
\tag{UPR24}
\]
The row multiplicities in (UPR18) and the word Gram (UPR23) have
different exact roles.  No additional multinomial is inserted
into either of them.

The original AKS9 map in root-value coordinates is
\[
\jmath\Lambda_k v=\mathcal F_0B_kv.
\tag{UPR25}
\]
The map \(\mathcal F_0\) is injective, with
\(\|\mathcal F_0^+\|\le5^{-k/2}\) by (UPR24).
For \(y\) in the actual image of (UPR25), define
\[
T_k y=B_k^+\mathcal F_0^+y.
\]
Then \(B_kT_ky=\mathcal F_0^+y\) and
\(\mathcal F_0B_kT_ky=y\), because the latter vector lies in
\(\operatorname{im}B_k\).  Furthermore
\[
\|T_ky\|_2\le5^{-k/2}\mathscr R_k\|y\|_{\rm original\ output}.
\tag{UPR26}
\]
This is a specified linear right inverse on the image.  Its values
may be passed to the quotient by \(\ker B_k\); it does not recover
components in that kernel.

\subsection{Original arithmetic quotient metric at every cutoff}
Now keep the original grid
\[
q=(k+1)^2,\quad c=k/2,\quad
\lambda_{a,b}=c+(2a-k)\delta+i(2b-k)\gamma,\quad
\chi(S)=\prod_{a,b=0}^k(S-\lambda_{a,b}),
\]
and let \(R=k\sqrt{\delta^2+\gamma^2}\),
\(d_*=2\min(\delta,\gamma)>0\).
Distinct roots have separation at least \(d_*\), because at least
one of their two grid-coordinate differences is a nonzero integer.
For an actual root-value column \(v\), the degree-\((q-1)\)
Lagrange representative is
\[
P_v(S)=\sum_\lambda
v_\lambda\frac{\chi(S)}{(S-\lambda)\chi'(\lambda)}.
\tag{UPR27}
\]
It has exactly those root values.  For \(S=c+iy\),
each numerator factor has modulus at most \(|y|+R\), whereas
\(|\chi'(\lambda)|\ge d_*^{q-1}\).
Cauchy--Schwarz in the actual \(q\)-term sum proves
\[
|P_v(c+iy)|^2
\le q\,d_*^{-2q+2}(|y|+R)^{2q-2}\|v\|_2^2.
\tag{UPR28}
\]

Retain the arithmetic source
\[
w_h(y)=|2\xi(1/2+iy)/h(1/2+iy)|^2/(2\pi),\qquad
m_k=w_h^{*k},
\quad
M_h(\eta)=\int_{\mathbb R}e^{\eta y}w_h(y)\,dy,\quad
\eta=\pi/4.
\]
The inherited exponential envelope gives a finite positive
\(M_h(\eta)\); no probability rescaling is made.
The nonnegative exponential series proves
\[
t^n\le n!\eta^{-n}e^{\eta t}\qquad(t\ge0,\ n\in\mathbb N).
\]
Tonelli's theorem in the original convolution coordinates gives
\(\int e^{\eta y}m_k(y)\,dy=M_h(\eta)^k\).
Evenness of the unchanged density yields
\[
\int_{\mathbb R}e^{\eta|y|}m_k(y)\,dy
=2\int_0^\infty e^{\eta y}m_k(y)\,dy
\le2M_h(\eta)^k.
\]
Apply these facts to (UPR28), including the factor \(e^{\eta R}\):
\[
\|P_v\|_{m_k}^2
\le\mathscr B_{h,k}\|v\|_2^2,\qquad
\mathscr B_{h,k}=
2q(2q-2)!(\eta d_*)^{-2q+2}
e^{\eta R}M_h(\eta)^k.
\tag{UPR29}
\]
Both half-lines, every root and every original source mass remain.

For each original cutoff \(N\ge q-1\), the minimum metric \(G_N\)
on the class of \(P_v\) is bounded above by the norm of this
particular admitted representative.  Composing (UPR26) with
(UPR27), and then with the actual class map, gives a linear
right inverse \(\mathscr T_{k,N}\) on the original observation image:
\[
\boxed{\displaystyle
\|\mathscr T_{k,N}\|_{{\rm original\ output}\to G_N}
\le \sqrt{\mathscr B_{h,k}}\,5^{-k/2}\mathscr R_k.}
\tag{UPR30}
\]
Passing its value further to the original quotient by
\(\ker\Lambda_k\) can only decrease the quotient norm, so the same
bound applies to that receiver.
Since \(\log(2q-2)!\le(2q-2)\log(2q-2)\), \(R=O_h(k)\),
and (UPR17) has logarithm \(O_{h,u}(k+\log(k+1))\),
\[
\log\|\mathscr T_{k,N}\|
\le O_{h,u}(q\log(q+2)),
\tag{UPR31}
\]
uniformly over every such \(N\).
All constants refer to the original fixed quartet, unit, period
and branch.  The construction is a large, finite recovery allowance
on the actual image.  It supplies no small native angular estimate,
does not remove the observation kernel, and does not identify the
arithmetic quotient metric with the independent proper-source metric.
\endgroup



% END PRESERVED INPUT 04_RESIDUE_PERIOD_COMPLETE.tex
\endgroup

\clearpage
\begingroup
\def\R{\mathbb R}
\def\C{\mathbb C}
\def\PF{\operatorname{PF}}
\def\dd{\,d}

\setcounter{equation}{0}
\renewcommand{\theequation}{A03MA.\arabic{equation}}
\section*{The original Montgomery finite-part identity}
\addcontentsline{toc}{section}{The original Montgomery finite-part identity}
% BEGIN PRESERVED INPUT 05_MONTGOMERY_FINITE_PART_COMPLETION.tex SHA256 ead0f65531011c7202aa84ecad750fcaff412c05c53680cb8ce1fac652249792

\noindent The original Six Readers files are preserved unchanged. This note
completes the scalar identity left open after
\texttt{eq:poisson-pole-complete} in
\nolinkurl{proofs/B_guinand_weil_montgomery_interfaces.tex}, and in the
Montgomery paragraph of
\texttt{satellites/10\_literature\_foundations.tex}.
It independently derives the new incoming note
\texttt{03\_MONTGOMERY\_ARCHIMEDEAN\_IDENTITY.md} from the literal test and
finite-part convention in those complete sources. The distinct false local
Guinand--Weil residual equality is corrected in the companion derivation;
it is not used anywhere in this proof.

\section{The original parameters, transform and domain}
Throughout retain
\begin{equation}\tag{MA1}\begin{gathered}
 1<\sigma_0<2,\quad t\in\R,\quad x\ge1,\quad U=\log x,\quad
 a=\sigma_0-\tfrac12,\\
 F(v)=e^{-a|v-U|}e^{-it(v-U)},\quad
 C=F(0)=x^{1/2-\sigma_0+it},\\
 z=1-\sigma_0+it,\quad w=\sigma_0-it=1-z,\quad
 p=\sigma_0+it,\quad A=(2\pi)^{-1}.
\end{gathered}\end{equation}
For every positive base, a complex power means $y^s=\exp(s\log y)$ with
the real natural logarithm. The digamma function $\psi=\Gamma'/\Gamma$
is meromorphic, not a chosen branch of $\log\Gamma$.
Weil's real-place kernel and the retained finite part \cite{a03ma:Weil} are
\begin{equation}\tag{MA2}\begin{gathered}
 K(v)=\frac{e^{|v|/2}}{|e^v-e^{-v}|},\\
 \PF(FK)=\lim_{\Lambda\to\infty}
 \left\{\int_\R(1-e^{-\Lambda|v|})F(v)K(v)\dd v-C\log\Lambda\right\}.
\end{gathered}\end{equation}
Indeed $|v|K(v)\to1/2$ and $F(v)\to C$ as $v\to0$, so the original
subtraction $2\beta(0)\log\Lambda$ has $\beta(0)=C/2$. In particular its
coefficient is $C$, not $2C$.

The test $F$ is continuous. On the two intervals separated by $U$ its
derivative is $(a-it)F$ and $(-a-it)F$, respectively. The average derivative
value at the single jump is $-it$, since $F(U)=1$. These are exactly the
first-kind jump data of Weil's test convention. For every
$0<b<\sigma_0-1$,
\[
 |F(v)|+|F'(v)|\le(1+a+|t|)e^{aU}e^{-a|v|}
 =O(e^{-(1/2+b)|v|}).
\]
Thus both literal conditions (A) and (B) of the recorded Weil source hold,
including when $U=0$.
For $1-\sigma_0<\Re s<\sigma_0$, absolute convergence permits $v=U+u$:
\begin{equation}\tag{MA3}\begin{aligned}
 \Phi(s)&:=\int_\R F(v)e^{(s-1/2)v}\dd v\\
 &=x^{s-1/2}\left[
 \int_0^\infty e^{-[a-(s-1/2-it)]u}\dd u
 +\int_0^\infty e^{-[a+(s-1/2-it)]u}\dd u\right]\\
 &=\frac{2a x^{s-1/2}}{a^2-(s-1/2-it)^2}.
\end{aligned}\end{equation}
The strip contains $s=0,1$ and the entire closed critical strip. Factoring
the two original denominators yields
\begin{equation}\tag{MA4}
 \Phi(0)=-\frac{2a x^{-1/2}}{zp},\qquad
 \Phi(1)=\frac{2a x^{1/2}}{(\sigma_0-1+it)(\sigma_0-it)}.
\end{equation}
There is no zero denominator on (MA1): $-1<\Re z<0$, $1<\Re w,\Re p<2$.

\section{The finite part as an ordinary, convergent regularized integral}
On $v>0$ write
\[
 J(v)=[F(v)+F(-v)]K(v).
\]
The exact kernel is $K(v)=e^{-v/2}/(1-e^{-2v})$.
It is $1/(2v)+O(1)$ near zero, while
$F(v)+F(-v)=2C+O(v)$ by local Lipschitz continuity, including at $U=0$.
Consequently $J(v)=C/v+O(1)$. For $v\ge\max(1,U)$, the test has decay
$O_{\sigma_0,t,x}(e^{-av})$ and $K(v)=O(e^{-v/2})$, giving
$J(v)=O_{\sigma_0,t,x}(e^{-\sigma_0v})$.
Therefore the function
\begin{equation}\tag{MA5}
 H(v)=J(v)-C\frac{e^{-v}}v
\end{equation}
is absolutely integrable on $(0,\infty)$.
For finite $\Lambda>0$, elementary differentiation in $\Lambda$ and the
value zero at $\Lambda=0$ prove
\[
 \int_0^\infty\frac{e^{-v}-e^{-(1+\Lambda)v}}v\dd v=\log(1+\Lambda).
\]
It follows exactly that
\begin{equation}\tag{MA6}
 \int_\R(1-e^{-\Lambda|v|})FK\dd v-C\log\Lambda
 =\int_0^\infty(1-e^{-\Lambda v})H(v)\dd v
  +C\log(1+\Lambda^{-1}).
\end{equation}
Dominated convergence with the proved $L^1$ function $H$ now gives
\begin{equation}\tag{MA7}
 \boxed{\PF(FK)=\int_0^\infty
 \left\{[F(v)+F(-v)]K(v)-C\frac{e^{-v}}v\right\}\dd v.}
\end{equation}
This is an exact linear relation between the original finite-part functional
and an ordinary regularized integral. The subtracted function and its
constant have both been retained. The same estimates are uniform when
$(\sigma_0,t,x)$ varies over compact subsets of the domain (MA1), so this
functional is continuous there, including the relative endpoint $x=1$.

\section{Direct integration on the two original intervals}
The positive kernel has the geometric expansion
\begin{equation}\tag{MA8}
 K(v)=\sum_{m=0}^\infty e^{-(2m+1/2)v},\qquad v>0.
\end{equation}
For a fixed $m\ge0$ put $\lambda_m=2m+1/2$. The actual test on the two
intervals is
\begin{equation}\tag{MA9}\begin{aligned}
 F(v)&=Ce^{(a-it)v}&& (0\le v\le U),\\
 F(v)&=Cx^{2a}e^{-(a+it)v}&& (v\ge U),\\
 F(-v)&=Ce^{-(a-it)v}&& (v\ge0).
\end{aligned}\end{equation}
In particular the first interval is empty when $x=1$, and remains finite
when $m=0$ and $\Re z<0$. Direct integration gives
\begin{equation}\tag{MA10}\begin{aligned}
 I_m&:=\int_0^\infty[F(v)+F(-v)]e^{-\lambda_m v}\dd v\\
 &=C\frac{1-e^{-(2m+z)U}}{2m+z}
   +Cx^{2a}\frac{e^{-(2m+p)U}}{2m+p}
   +\frac C{2m+w}\\
 &=C\left(\frac1{2m+w}+\frac1{2m+z}\right)
       -\frac{2a x^{-2m-1/2}}{(2m+z)(2m+p)}.
\end{aligned}\end{equation}
Here $Cx^{-(2m+z)}=Cx^{2a-(2m+p)}=x^{-2m-1/2}$ and $p-z=2a$.
No integral of $e^{-(2m+z)v}$ over the infinite half-line was used.

For finite $\Lambda>2$ all integrations and (MA8) may be exchanged in the
regulated integral. To prove this for complex $F$, apply Tonelli to the
absolute values: the dominating function is
$(1-e^{-\Lambda v})(|F(v)|+|F(-v)|)K(v)$, bounded near zero and
exponentially integrable at infinity by the estimates preceding (MA5).
The same integration as (MA10), with every $2m$ replaced by $2m+\Lambda$,
gives
\begin{equation}\tag{MA11}\begin{aligned}
 \int_\R(1-e^{-\Lambda|v|})FK\dd v
 ={}&C\sum_{m\ge0}\left[
 \frac1{2m+w}-\frac1{2m+w+\Lambda}
 +\frac1{2m+z}-\frac1{2m+z+\Lambda}\right]\\
 &-2a x^{-1/2}\sum_{m\ge0}
 \frac{x^{-2m}}{(2m+z)(2m+p)}\\
 &+2a x^{-1/2-\Lambda}\sum_{m\ge0}
 \frac{x^{-2m}}{(2m+z+\Lambda)(2m+p+\Lambda)}.
\end{aligned}\end{equation}
Each bracketed harmonic difference is $O_{\Lambda,\sigma_0,t}(m^{-2})$;
the two rational series are separately absolutely convergent. Thus (MA11)
splits only convergent series, not divergent harmonic sums.

\section{The exact digamma difference and its complex limit}
Euler's locally uniformly convergent product for $1/\Gamma$, as recorded in
DLMF 5.8.2 \cite{a03ma:Gamma}, gives by local logarithmic differentiation
\begin{equation}\tag{MA12}
 \psi(b)-\psi(a_0)
 =\sum_{m=0}^\infty\left(\frac1{m+a_0}-\frac1{m+b}\right),
\end{equation}
for $a_0,b$ away from the nonpositive integers. The differentiated differences
converge locally uniformly since they are $O(m^{-2})$. One can first
differentiate on the positive real half-plane and then continue meromorphically;
this yields the same expression without imposing a logarithm branch at
$z/2$, whose real part lies in $(-1/2,0)$.
Applying (MA12) to (MA11) evaluates its first line exactly as
\begin{equation}\tag{MA13}
 \frac C2\left[
 \psi((w+\Lambda)/2)-\psi(w/2)
 +\psi((z+\Lambda)/2)-\psi(z/2)\right].
\end{equation}

For clarity, the limiting formula needed here follows directly from the same
product. For real $r>0$ it gives
$\psi(r)=\lim_{N\to\infty}[\log N-\sum_{m=0}^N(m+r)^{-1}]$.
The decreasing-integrand sum comparison gives
\[
 0\le\sum_{m=0}^N\frac1{m+r}
      -\int_0^{N+1}\frac{du}{u+r}\le\frac1r,
 \qquad |\psi(r)-\log r|\le\frac1r.
\]
For $r\ge\max(1,2|c_0|)$, (MA12) also gives
\[
 |\psi(r+c_0)-\psi(r)|
 \le2|c_0|\sum_{m\ge0}(m+r)^{-2}
 \le\frac{4|c_0|}{r}.
\]
Taking $r=\Lambda/2$, $c_0=s/2$ proves the complex estimate
\begin{equation}\tag{MA14}
 \left|\psi((\Lambda+s)/2)-\log\Lambda+\log2\right|
 \le\frac{2+4|s|}{\Lambda},\qquad
 \Lambda\ge\max(2,2|s|).
\end{equation}
Thus no passage to real parts of the complex identity is involved.
The shifted rational series in (MA11) satisfies, for $\Lambda>2$,
\begin{equation}\tag{MA15}
 \left|\sum_{m\ge0}
 \frac{x^{-2m}}{(2m+z+\Lambda)(2m+p+\Lambda)}\right|
 \le\frac1{(\Lambda-1)^2}+\frac1{2(\Lambda-1)}.
\end{equation}
Indeed the real parts of the two denominator factors exceed
$2m+\Lambda-1$, and the last bound is the sum-plus-integral bound for
$(2u+\Lambda-1)^{-2}$. Multiplication by $x^{-\Lambda}$ makes the error
smaller when $x>1$; (MA15) also proves its convergence to zero directly
at $x=1$.

Subtracting $C\log\Lambda$ in (MA11), then using (MA13)--(MA15), proves
\begin{equation}\tag{MA16}\boxed{\begin{aligned}
 \PF(FK)={}&-C\log2
 -\frac C2[\psi(w/2)+\psi(z/2)]\\
 &-2a x^{-1/2}\sum_{m=0}^\infty
       \frac{x^{-2m}}{(2m+z)(2m+p)}.
\end{aligned}}\end{equation}
More quantitatively, for
$\Lambda\ge\max(3,2|w|,2|z|)$ the difference between the regulated finite
part in (MA2) and (MA16) is bounded by
\begin{equation}\tag{MA17}
 \frac{|C|[2+2(|w|+|z|)]}{\Lambda}
 +2a x^{-1/2-\Lambda}
 \left[\frac1{(\Lambda-1)^2}+\frac1{2(\Lambda-1)}\right].
\end{equation}
This controls the actual regulator with its original subtraction.
For $M\ge1$, the tail of the unshifted rational series after $m=M$ satisfies
\begin{equation}\tag{MA18}
 \left|\sum_{m>M}\frac{x^{-2m}}{(2m+z)(2m+p)}\right|
 \le\frac1{2M}.
\end{equation}
For $x>1$ it also satisfies the stronger bound
$x^{-2M-2}/[2(M+1)^2(1-x^{-2})]$.
To verify this, use $|2m+z|>2m-1\ge m$ and $|2m+p|>2m$ for $m\ge1$,
then the integral or geometric bound. The separate $m=0$ term is finite
since $z,p\ne0$.

\section{Both prime orientations and the exact reflection identity}
Let $\Lambda(n)$ be the von Mangoldt function, with $\Lambda(1)=0$.
The two orientations are, directly from (MA9),
\begin{equation}\tag{MA19}\begin{aligned}
 n^{-1/2}F(\log n)
 &=\begin{cases}
 Cn^{\sigma_0-1-it},& n\le x,\\
 x^{\sigma_0-1/2+it}n^{-\sigma_0-it},& n>x,
 \end{cases}\\
 n^{-1/2}F(-\log n)&=Cn^{-w}.
\end{aligned}\end{equation}
At $n=x$ the two displayed positive-orientation expressions both equal
$n^{-1/2}$. Thus the split does not create a jump or a half-weight at
a prime power. At $x=1$ the first sum contains only $n=1$ and is zero.
Absolute convergence of the infinite prime sums follows from
$0\le\Lambda(n)\le\log n$ and $\sigma_0>1$.

For completeness, the Euler product for $\zeta(s)$ on $\Re s>1$
\cite{a03ma:Zeta} is nonzero and has the locally absolutely convergent logarithm
$\sum_{\ell\ \mathrm{prime}}\sum_{r\ge1}\ell^{-rs}/r$.
Absolute convergence follows from
$\sum_\ell\sum_{r\ge1}\ell^{-r\sigma_0}/r<\infty$;
local termwise differentiation is permitted by the corresponding convergent
$\log n\,n^{-\sigma_0}$ bound. Unique prime-power factorization therefore
gives the exact signs
\begin{equation}\tag{MA20}
 P_-:=\sum_{n\ge1}\Lambda(n)n^{-1/2}F(-\log n)
 =C\sum_{n\ge1}\Lambda(n)n^{-w}
 =-C\frac{\zeta'(w)}{\zeta(w)},\qquad
 -P_-=+C\frac{\zeta'(w)}{\zeta(w)}.
\end{equation}
For example, for any integer $N$ beyond $e^{1/\sigma_0}$, the absolute tail
of the Dirichlet sum after $N$ is at most
\[
 \int_N^\infty(\log u)u^{-\sigma_0}\dd u
 =N^{1-\sigma_0}\left[\frac{\log N}{\sigma_0-1}
                       +\frac1{(\sigma_0-1)^2}\right].
\]
Its outside factor $|C|$ is retained for $P_-$.

Put $\mathscr L(s)=\pi^{-s/2}\Gamma(s/2)\zeta(s)$.
The exact zeta functional equation is $\mathscr L(s)=\mathscr L(1-s)$,
equivalently the $\xi$ reflection with its $s(s-1)/2$ factor
\cite{a03ma:Zeta}. At $w$ the Euler product and the nonvanishing Gamma factor
give $\mathscr L(w)\ne0$. At $z=1-w$, which has $-1<\Re z<0$,
$\Gamma(z/2)$ is finite and nonzero, so reflection also proves
$\zeta(z)\ne0$. Thus every logarithmic derivative below exists.
Differentiating the equality as a meromorphic equality and dividing by its
nonzero value gives
\begin{equation}\tag{MA21}\begin{aligned}
 -\tfrac12\log\pi+\tfrac12\psi(z/2)+\frac{\zeta'(z)}{\zeta(z)}
 &=-\left[-\tfrac12\log\pi+\tfrac12\psi(w/2)
                           +\frac{\zeta'(w)}{\zeta(w)}\right],\\
 \frac{\zeta'(w)}{\zeta(w)}-\log\pi
 +\tfrac12[\psi(w/2)+\psi(z/2)]
 &=-\frac{\zeta'(z)}{\zeta(z)}.
\end{aligned}\end{equation}
This derivation uses ratios, not an unproved global logarithm branch.

\section{Pole cancellation and the completed identity}
Substitute (MA16) and (MA20) into the literal residual expression. Its
constant is $\log A=-\log(2\pi)$, so $-\log(2\pi)+\log2=-\log\pi$.
Using (MA21) only after that cancellation yields
\begin{equation}\tag{MA22}\begin{aligned}
 &\Phi(0)+C\log A-\PF(FK)-P_-\\
 &=\Phi(0)+C\left\{\frac{\zeta'(w)}{\zeta(w)}-\log\pi
              +\frac{\psi(w/2)+\psi(z/2)}2\right\}
       +2a x^{-1/2}\sum_{m\ge0}\frac{x^{-2m}}{(2m+z)(2m+p)}\\
 &=-C\frac{\zeta'(z)}{\zeta(z)}
   +\underbrace{\Phi(0)+\frac{2a x^{-1/2}}{zp}}_{=0\text{ by (MA4)}}
   +2a x^{-1/2}\sum_{m\ge1}\frac{x^{-2m}}{(2m+z)(2m+p)}.
\end{aligned}\end{equation}
The two denominator presentations are exactly related by
\[
 2m+z=-(\sigma_0-1-it-2m),\qquad 2m+p=\sigma_0+it+2m.
\]
Therefore (MA22) is precisely the required complete scalar identity
\begin{theorem}
For every original parameter tuple (MA1), with the literal finite part
(MA2), all the terms below exist and
\begin{equation}\tag{MA23}\begin{aligned}
 &\Phi(0)+C\log((2\pi)^{-1})-\PF(FK)
   -\sum_{n\ge1}\Lambda(n)n^{-1/2}F(-\log n)\\
 &\quad=-C\frac{\zeta'(1-\sigma_0+it)}{\zeta(1-\sigma_0+it)}
 -x^{-1/2}\sum_{m\ge1}
 \frac{(2\sigma_0-1)x^{-2m}}
      {(\sigma_0-1-it-2m)(\sigma_0+it+2m)}.
\end{aligned}\end{equation}
\end{theorem}
The proof has not used a zero sum, a zero-location hypothesis, or the
incorrect printed Guinand--Weil residual equality.

\section{The endpoint and the receiving Montgomery formula}
At $x=1$, the actual test is $F(v)=e^{-a|v|}e^{-itv}$ and $C=1$.
The direct two-interval integration was already valid with $U=0$.
Independently, (MA12) and $p-z=2a$ give
\[
 2a\sum_{m\ge0}\frac1{(2m+z)(2m+p)}
 =\sum_{m\ge0}\left(\frac1{2m+z}-\frac1{2m+p}\right)
 =\tfrac12[\psi(p/2)-\psi(z/2)].
\]
Substitution in (MA16) proves the endpoint value
\begin{equation}\tag{MA24}
 \boxed{\PF(FK)=-\log2-\tfrac12
 [\psi((\sigma_0-it)/2)+\psi((\sigma_0+it)/2)]\quad(x=1).}
\end{equation}
All series converge directly at this endpoint by (MA18); no omitted
limiting or prime-power convention is required.

We now receive the established identity into the recorded original Weil
formula, with its symmetric zero cutoff. Its rational trivial-character
specialization is \cite[formula (11)]{a03ma:Weil}
\begin{equation}\tag{MA25}
 \lim_{T\to\infty}\sum_{|\Im\rho|<T}\Phi(\rho)
 =\Phi(0)+\Phi(1)+C\log A-P_+-P_--\PF(FK),
\end{equation}
where $P_+$ is the positive-orientation sum from (MA19) and all nontrivial
zeros are counted with multiplicity. We use this classical explicit formula
at its original test domain, whose hypotheses were verified after (MA2).
The full source provides its literal normalization and the two prime orientations.
Substitution of the independently proved (MA23) gives the zero-side formula
without changing any zero coordinate.

On the real-ordinate domain used in Montgomery's 1973 paper, namely RH,
$\rho=1/2+i\gamma$ with $\gamma\in\R$ and
\[
 \Phi(\rho)=\frac{(2\sigma_0-1)x^{i\gamma}}
 { (\sigma_0-1/2)^2+(t-\gamma)^2}.
\]
Equations (MA4), (MA19), (MA23) and (MA25) therefore produce exactly
\begin{equation}\tag{MA26}\begin{aligned}
 &(2\sigma_0-1)\sum_\gamma
 \frac{x^{i\gamma}}{(\sigma_0-1/2)^2+(t-\gamma)^2}\\
 &=-x^{-1/2}\left[
 \sum_{n\le x}\Lambda(n)(x/n)^{1-\sigma_0+it}
 +\sum_{n>x}\Lambda(n)(x/n)^{\sigma_0+it}\right]\\
 &\quad-C\frac{\zeta'(1-\sigma_0+it)}{\zeta(1-\sigma_0+it)}
 +\frac{x^{1/2}(2\sigma_0-1)}{(\sigma_0-1+it)(\sigma_0-it)}\\
 &\quad-x^{-1/2}\sum_{m\ge1}
 \frac{(2\sigma_0-1)x^{-2m}}
 {(\sigma_0-1-it-2m)(\sigma_0+it+2m)}.
\end{aligned}\end{equation}
This is Montgomery's complete retained equation \cite[pp.~185--187]{a03ma:Montgomery}.
The sum means at least the original symmetric cutoff in (MA25); its
convergence is supplied by that cited theorem. No later real-ordinate estimate
is declared valid for a complex ordinate by this receiving calculation.
The scalar (MA23) itself is unconditional on (MA1); the RH qualification in
(MA26) belongs only to the original choice of real zero coordinates.

\section{Exact backward replacement}
In the complete Original Six Readers provider, the proved transform,
prime-range and pole formulas remain as printed. Replace the status sentence
``The remaining gamma/trivial-zero commutator is not claimed proved'' and
the concluding paragraph beginning ``Still to be proved is'' in
\nolinkurl{proofs/B_guinand_weil_montgomery_interfaces.tex} by (MA23),
with the proof (MA5)--(MA24). In
\texttt{satellites/10\_literature\_foundations.tex}, replace the paragraph
beginning ``What is still open is one explicit commutator'' by the same
completed identity and its exact parameter domain. These are successor
replacement instructions; no archived source bytes are changed here.

The separate local labels \texttt{eq:guinand-weil-arch-map} and
\nolinkurl{eq:guinand-weil-residual-complete} need the nonzero correction proved
in the companion note, not (MA23) with a discarded term. On its explicit
common test subspace the exact additional functional is
\[
 \mathcal D(f)=\int_0^\infty f(v)
       [2e^{-v/2}-\operatorname{sech}(v/2)]\dd v.
\]
The original right side must receive $+\mathcal D(f)$. In particular this
change must be propagated to the asserted local subtraction proof and its
status records; it is not a correction to Guinand's published theorem.
The companion proves the exact test-space inclusion with decay exponent
strictly greater than $1/2$, all jump terms, and a nonzero admissible example.



% END PRESERVED INPUT 05_MONTGOMERY_FINITE_PART_COMPLETION.tex
\endgroup

\clearpage
\begingroup
\def\R{\mathbb R}
\def\C{\mathbb C}
\def\PF{\operatorname{PF}}
\def\sech{\operatorname{sech}}

\setcounter{equation}{0}
\renewcommand{\theequation}{A03GC.\arabic{equation}}
\section*{Exact correction of the archived Guinand--Weil identity}
\addcontentsline{toc}{section}{Exact correction of the archived Guinand--Weil identity}
% BEGIN PRESERVED INPUT 06_GUINAND_ARCHIMEDEAN_EXACT_CORRECTION.tex SHA256 d58f7cfd78571071edb19d97c9fd1dadc71f62324410ba595e69ffd6ce432f40



This note concerns the equation labelled
\texttt{eq:guinand-weil-arch-map} and its expanded version
\texttt{eq:guinand-weil-residual-complete} in the programme's
\emph{Original Six Readers}.  The local equations, cosine convention and
finite-part subtraction are retained literally.  The result is a direct
calculation of their exact difference, with no use of RH or of a zero sum.
The incoming note \emph{Exact correction to the archived Guinand--Weil
comparison}, 19 September 2026, identified this difference.  Here every
interchange and endpoint is proved, the test-domain assertion is corrected,
and the entire exponential family is evaluated.  The result corrects the
local transcription/comparison; it makes no assertion against Guinand's
published theorem \cite{a03gc:guinand1948} or Weil's published formula
\cite{a03gc:weil1952}.

\section{The original functions and the exact domain inclusion}

For $c>0$, define $\mathcal D_c$ to consist of complex continuous functions
$f:[0,\infty)\to\C$ with a finite partition
$0=a_0<a_1<\cdots<a_m<\infty$ such that $f$ is $C^2$ on each open piece,
$f''$ is absolutely integrable on every piece, and $f'$ has finite
one-sided values at the finite endpoints.  On the final unbounded piece
require
\begin{equation}
 |f(x)|+|f'(x)|+|f''(x)|\le C_f e^{-cx}
 \quad (x\ge X_f).
 \label{a03gc:gc:decay}
\end{equation}
At a derivative jump use its two one-sided values, and use their average
as its assigned representative.  This representative does not change an
integral.  Put
\begin{equation}
 \mathcal D_0^{\exp}=\bigcup_{c>0}\mathcal D_c,
 \qquad \mathcal D_W^{\exp}=\bigcup_{c>1/2}\mathcal D_c.
 \label{a03gc:gc:domains}
\end{equation}
These are vector spaces: for addition take the common refinement of the
two finite partitions and the smaller of the two decay exponents.

The original definitions are
\begin{align}
 g_f(t)&=\sqrt{\frac2\pi}\int_0^\infty f(x)\cos(tx)\,dx,
 &&t\ge0,\label{a03gc:gc:cosine}\\
 2H_f&=\int_0^\infty f(x)
       \left(\frac1x-\frac{e^{-x/2}}{\sinh x}\right)dx,
 \label{a03gc:gc:H}\\
 K(x)&=\frac{e^{|x|/2}}{|e^x-e^{-x}|},\qquad x\ne0,
 \label{a03gc:gc:K}\\
 P_f&=\lim_{\Lambda\to\infty}
 \left\{\int_\R(1-e^{-\Lambda|x|})f(|x|)K(x)\,dx
              -f(0)\log\Lambda\right\},
 \label{a03gc:gc:PF}\\
 \mathscr L_f&=\sqrt{\frac2\pi}\int_0^\infty
             g_f(t)\log\frac{t}{2\pi}\,dt,
 \label{a03gc:gc:L}\\
 W_f&=2\int_0^\infty f(x)e^{-x/2}\,dx-P_f-f(0)\log(2\pi).
 \label{a03gc:gc:W}
\end{align}
In \eqref{a03gc:gc:PF}, $|x|f(|x|)K(x)\to f(0)/2$, so the original
Weil subtraction $2\beta(0)\log\Lambda$ is exactly
$f(0)\log\Lambda$; this factor is not adjusted.
The archived equality was $W_f=2H_f+\mathscr L_f$.

\begin{proposition}[Exact inclusion in the recorded common domain]
For every $f\in\mathcal D_W^{\exp}$ and $1<p\le2$, the function $f$
satisfies the Guinand test conditions recorded in the source:
it is an indefinite integral of its locally integrable derivative,
$f(x)\to0$, and $xf'(x)\in L^p(0,\infty)$.
The map
\begin{equation}
 E:\mathcal D_W^{\exp}\longrightarrow
       \mathcal W_{\rm even}^{G,p},\qquad
 E(f)(x)=\frac{f(|x|)}{\sqrt{2\pi}},
 \qquad f(t)=\sqrt{2\pi}\,E(f)(t)\ (t\ge0)
 \label{a03gc:gc:E}
\end{equation}
is an injective linear map into the recorded Guinand--Weil common test
space, and the displayed restriction is its inverse on its image.
Moreover
\begin{equation}
 \int_\R E(f)(x)e^{itx}\,dx=g_f(t),\qquad
 T^{1/2}f(T)\longrightarrow0.
 \label{a03gc:gc:transform-map}
\end{equation}
\end{proposition}
\begin{proof}
The integrability of $f''$ and finite derivative traces make $f'$ bounded
on each bounded piece.  Continuity of $f$ cancels the boundary values in
the fundamental theorem on adjacent pieces.  Thus
$f(x)=f(0)+\int_0^x f'(u)du$ on every bounded interval.
The exponential bound gives the limit at infinity and the $L^p$ assertion;
near zero and each finite break the derivative is bounded.
The even extension is continuous and piecewise continuously
differentiable.  Its derivative at $0$ has one-sided values
$-f'(0+)/\sqrt{2\pi}$ and $f'(0+)/\sqrt{2\pi}$; at $\pm a_j$ it
has precisely the corresponding reflected derivative jumps.  Thus it
satisfies the recorded Weil condition (A), including the averaged
representatives.  Choose a decay exponent $c>1/2$ in \eqref{a03gc:gc:decay} and
put $b=c-1/2>0$.  Then both $E(f)$ and its piecewise derivative are
$O(e^{-(1/2+b)|x|})$, which is condition (B).
The two expressions in \eqref{a03gc:gc:E} compose to the identity.  Splitting
the integral at zero gives
$2/\sqrt{2\pi}=\sqrt{2/\pi}$ times the cosine integral, proving
\eqref{a03gc:gc:transform-map}.  Its last limit follows directly from
\eqref{a03gc:gc:decay}.  These are the common-domain conditions recorded from
Guinand's Theorems 3 and $3'$ and Weil's (A), (B)
\cite[pp.~112--115]{a03gc:guinand1948}\cite[Collected Papers II, pp.~52--58]{a03gc:weil1952}.
\end{proof}

The same expression $E(f)$ is a defined even Fourier test for
$f\in\mathcal D_0^{\exp}$ and satisfies the same transform identity,
but the inclusion into Weil's (B) requires the faster exponent.
For example $f(x)=e^{-x/4}$ belongs to $\mathcal D_0^{\exp}$, whereas
\begin{equation}
 e^{(1/2+b)x}|E(f)(x)|
   =\frac{e^{(1/4+b)x}}{\sqrt{2\pi}}\longrightarrow\infty
   \quad (b>0).
 \label{a03gc:gc:domain-counterexample}
\end{equation}
Consequently the incoming note's phrase ``exponential decay'' alone does
not prove its claimed inclusion in the original common domain.  Equations
\eqref{a03gc:gc:domains}--\eqref{a03gc:gc:transform-map} give the exact inclusion and
inverse.  All functional identities below hold on
$\mathcal D_0^{\exp}$, and their application as a correction within the
recorded common domain is restricted to $\mathcal D_W^{\exp}$.

\section{Cosine decay, inversion at zero, and the logarithmic functional}

\begin{lemma}[Jump-retaining decay and the endpoint inversion]
For $f\in\mathcal D_0^{\exp}$ set
$\Delta_j=f'(a_j+)-f'(a_j-)$ and
$B_f=|f'(0+)|+\sum_{j=1}^m|\Delta_j|+\int_0^\infty|f''(x)|dx$,
where $f''$ is the piecewise derivative.  For $t>0$,
\begin{equation}
 g_f(t)=-\frac{\sqrt{2/\pi}}{t^2}
 \left(f'(0+)+\sum_{j=1}^m\Delta_j\cos(ta_j)
                   +\int_0^\infty f''(x)\cos(tx)dx\right).
 \label{a03gc:gc:IBP}
\end{equation}
In particular $g_f$ and $g_f(t)\log(t/(2\pi))$ are absolutely integrable,
and
\begin{equation}
 \sqrt{\frac2\pi}\int_0^\infty g_f(t)dt=f(0).
 \label{a03gc:gc:inversion}
\end{equation}
\end{lemma}
\begin{proof}
First integrate $f(x)\cos(tx)$ by parts on all pieces up to a finite
endpoint $R$ in the last piece.  Continuity makes the internal $f$
boundary terms cancel.  The term at zero contains $\sin0=0$, and the
term at $R$ tends to zero.  This gives
$\int f\cos=-t^{-1}\int f'\sin$.
On the second integration, $\sin(tx)=-t^{-1}(\cos(tx))'$.
The endpoint at zero contributes $f'(0+)/t$; the two values at $a_j$
give $\Delta_j\cos(ta_j)/t$; the endpoint at infinity vanishes.
Multiplication by the first factor $-1/t$ proves \eqref{a03gc:gc:IBP}.
Thus
$|g_f(t)|\le\sqrt{2/\pi}\min(\|f\|_1,B_f/t^2)$.
This proves integrability of the transform and of its logarithmic weight
both at zero and infinity.

For $r>0$, absolute convergence permits the exact Laplace--cosine
calculation
\begin{equation}
 \sqrt{\frac2\pi}\int_0^\infty e^{-rt}g_f(t)dt
  =\frac2\pi\int_0^\infty f(x)\frac{r}{r^2+x^2}dx
  =\frac2\pi\int_0^\infty\frac{f(ru)}{1+u^2}du.
 \label{a03gc:gc:Poisson}
\end{equation}
Indeed $\int_0^\infty e^{-rt}\cos(tx)dt$
is the real part of $(r-ix)^{-1}$ and equals $r/(r^2+x^2)$;
the double integral is bounded by $\|f\|_1/r$.
The function $f$ is bounded.  Dominated convergence in the right-hand
integral as $r\downarrow0$ gives $f(0)$, since
$\int_0^\infty(1+u^2)^{-1}du=\pi/2$.
Dominated convergence by $|g_f|$ in the left-hand integral proves
\eqref{a03gc:gc:inversion} without an unproved endpoint inversion assertion.
\end{proof}

\begin{theorem}[The exact cosine-log functional]
For $f\in\mathcal D_0^{\exp}$,
\begin{equation}
 \boxed{\displaystyle
 \mathscr L_f=-\int_0^\infty\frac{f(x)-f(0)e^{-x}}{x}dx
                      -f(0)\log(2\pi).}
 \label{a03gc:gc:log-formula}
\end{equation}
\end{theorem}
\begin{proof}
Set $h(x)=f(x)-f(0)e^{-x}$.  It belongs to $\mathcal D_0^{\exp}$,
$h(0)=0$, and $h(x)=O(x)$ near zero because its right derivative is
bounded there.  Its exponential decay gives
\begin{equation}
 \int_0^\infty\frac{|h(x)|}{x}dx<\infty,
 \qquad \int_0^\infty g_h(t)dt=0.
 \label{a03gc:gc:h}
\end{equation}
For $t>0$ the Frullani identity used here has an elementary proof:
\begin{equation}
 \int_0^\infty\frac{e^{-r}-e^{-tr}}{r}dr
 =\int_1^t\left(\int_0^\infty e^{-ur}dr\right)du
 =\int_1^t\frac{du}{u}=\log t.
 \label{a03gc:gc:Frullani}
\end{equation}
The interval $[1,t]$ is oriented if $t<1$; absolute integrability on the
bounded $u$ interval justifies the exchange.  For fixed $t$ the integrand
in $r$ has constant sign.  Therefore every truncated integral
$\ell_{\epsilon,R}(t)=\int_\epsilon^R(e^{-r}-e^{-tr})dr/r$
satisfies $|\ell_{\epsilon,R}(t)|\le|\log t|$.

At $0<\epsilon<R<\infty$, Fubini and \eqref{a03gc:gc:h} give
\begin{align}
 \sqrt{\frac2\pi}\int_0^\infty g_h(t)\ell_{\epsilon,R}(t)dt
 &=-\sqrt{\frac2\pi}\int_\epsilon^R\frac1r
                        \int_0^\infty e^{-rt}g_h(t)dt\,dr\notag\\
 &=-\frac2\pi\int_\epsilon^R\int_0^\infty
                           \frac{h(x)}{r^2+x^2}dx\,dr.
 \label{a03gc:gc:truncated-log}
\end{align}
For the first exchange, $2\|g_h\|_1\log(R/\epsilon)$ is a finite
majorant.  For inserting the cosine transform in the second line use
$\int_\epsilon^R r^{-1}\int_0^\infty e^{-rt}dt\,dr\,\|h\|_1<\infty$.
For the untruncated double integral, Tonelli applied to the absolute value
gives exactly
\begin{equation}
 \int_0^\infty\int_0^\infty\frac{|h(x)|}{r^2+x^2}dx\,dr
   =\frac\pi2\int_0^\infty\frac{|h(x)|}{x}dx<\infty.
 \label{a03gc:gc:absolute-double}
\end{equation}
The left side of \eqref{a03gc:gc:truncated-log} tends to the logarithmic
integral by the majorant $|g_h(t)\log t|$ from the lemma.  The right side
tends by \eqref{a03gc:gc:absolute-double} to $-\int h(x)/x\,dx$.

For the retained reference $e^{-x}$ its transform is
$\sqrt{2/\pi}/(1+t^2)$.  The substitution $t=1/u$ in the absolutely
convergent integral yields
$\int_0^\infty\log t\,(1+t^2)^{-1}dt=0$.
Hence its $\mathscr L$ value is $-\log(2\pi)$.
Adding $f(0)$ times this reference to the result for $h$ proves
\eqref{a03gc:gc:log-formula}, including its constant.
\end{proof}

\section{The original finite part and the defect map}

\begin{theorem}[Evaluation with the original regulator]
For $f\in\mathcal D_0^{\exp}$ the limit in \eqref{a03gc:gc:PF} exists and
\begin{equation}
 \boxed{\displaystyle
 P_f=\int_0^\infty\left(
       \frac{f(x)e^{x/2}}{\sinh x}-\frac{f(0)e^{-x}}{x}\right)dx.}
 \label{a03gc:gc:PF-formula}
\end{equation}
The displayed combined integrand is absolutely integrable.
\end{theorem}
\begin{proof}
Write $A_f(x)=f(x)e^{x/2}/\sinh x-f(0)e^{-x}/x$ for $x>0$.
The expansions
$e^{x/2}/\sinh x=x^{-1}+1/2+O(x)$ and
$e^{-x}/x=x^{-1}-1+O(x)$, together with
$f(x)=f(0)+f'(0+)x+o(x)$, show
$A_f(x)\to f'(0+)+3f(0)/2$ at zero.
At infinity its two terms decay exponentially.  Thus $A_f\in L^1$.
Evenness changes the full-line regulated integral into
$\int_0^\infty(1-e^{-\Lambda x})f(x)e^{x/2}/\sinh x\,dx$,
with no residual factor of two.  Subtract and add the reference to obtain
the exact finite-$\Lambda$ identity
\begin{align}
 &\int_\R(1-e^{-\Lambda|x|})f(|x|)K(x)dx-f(0)\log\Lambda\notag\\
 &\qquad=\int_0^\infty(1-e^{-\Lambda x})A_f(x)dx
              +f(0)\log(1+\Lambda^{-1}).
 \label{a03gc:gc:finite-regulator}
\end{align}
Here the reference integral is
$\int_0^\infty(e^{-x}-e^{-(1+\Lambda)x})dx/x=\log(1+\Lambda)$
by \eqref{a03gc:gc:Frullani}.  Dominated convergence by $|A_f|$, and
$\log(1+\Lambda^{-1})\to0$, prove \eqref{a03gc:gc:PF-formula}.
\end{proof}

\begin{theorem}[Exact comparison and receiving map]
Define
\begin{equation}
 d(x)=2e^{-x/2}-\sech(x/2)
     =\frac{2e^{-3x/2}}{1+e^{-x}},\qquad
 D:\mathcal D_0^{\exp}\to\C,\quad D(f)=\int_0^\infty f(x)d(x)dx.
 \label{a03gc:gc:D}
\end{equation}
Then the original expressions satisfy the exact identity
\begin{equation}
 \boxed{W_f-(2H_f+\mathscr L_f)=D(f).}
 \label{a03gc:gc:defect}
\end{equation}
In particular the map
\begin{equation}
 T:\mathcal D_0^{\exp}\to\C^3,\quad T(f)=(W_f,H_f,\mathscr L_f),
 \qquad (1,-2,-1)\,T=D
 \label{a03gc:gc:typed}
\end{equation}
is the exact typed relation between the two presentations.  Its restriction
along the inclusion $\mathcal D_W^{\exp}\hookrightarrow\mathcal D_0^{\exp}$
is the corrected common-domain comparison.
The defect is strictly positive on every nonzero real nonnegative test.
\end{theorem}
\begin{proof}
The kernel in \eqref{a03gc:gc:H} tends to $1/2$ at zero and is $O(1/x)$ at
infinity, proving absolute convergence of $H_f$.  All other terms were
proved convergent above.  Substitute \eqref{a03gc:gc:log-formula} and
\eqref{a03gc:gc:PF-formula}; to avoid separating divergent terms, first combine
their integrands on $[\eta,R]$ and then pass to the limits.  The
$f(0)\log(2\pi)$ constants and the reference subtractions cancel exactly,
as do the $f(x)/x$ terms.  The remaining kernel is
\begin{equation}
 2e^{-x/2}+\frac{e^{-x/2}-e^{x/2}}{\sinh x}
 =2e^{-x/2}-\frac1{\cosh(x/2)}=d(x).
 \label{a03gc:gc:kernel-algebra}
\end{equation}
It is bounded at zero and bounded by $2e^{-3x/2}$ at infinity.
Thus the limit is $D(f)$, proving \eqref{a03gc:gc:defect} and \eqref{a03gc:gc:typed}.
The second expression for $d$ in \eqref{a03gc:gc:D} is strictly positive.
A continuous nonnegative nonzero $f$ is positive on an interval, so its
integral against $d$ is strictly positive.
\end{proof}

For an explicit receiving functional retaining all old terms, set
\begin{align}
 H_f^{\rm rec}&=H_f+\tfrac12D(f),\label{a03gc:gc:receiving-H}\\
 2H_f^{\rm rec}
 &=\int_0^\infty f(x)
  \left(\frac1x-\frac{e^{x/2}}{\sinh x}+2e^{-x/2}\right)dx,
 \qquad W_f=2H_f^{\rm rec}+\mathscr L_f.
 \label{a03gc:gc:receiving-kernel}
\end{align}
The kernel in this integral tends to $3/2$ at zero.  Equations
\eqref{a03gc:gc:receiving-H}--\eqref{a03gc:gc:receiving-kernel} follow by the already
proved identity and \eqref{a03gc:gc:kernel-algebra}; they explicitly record the
map from the archived $H_f$.  They do not rename the archived $H_f$ or
identify $H_f^{\rm rec}$ with an uninspected historical transcription.

\section{The full exponential family and the original counterexample}

For $a>0$ let $f_a(x)=e^{-ax}$.  Its cosine transform, obtained by taking
the real part of $(a-it)^{-1}$, is
\begin{equation}
 g_{f_a}(t)=\sqrt{\frac2\pi}\frac{a}{a^2+t^2}.
 \label{a03gc:gc:exponential-transform}
\end{equation}
In particular $f_a\in\mathcal D_0^{\exp}$ for all $a>0$ and
$f_a\in\mathcal D_W^{\exp}$ for $a>1/2$.
Put $z=a/2+1/4$ and $w=a/2+3/4$.  We obtain
\begin{align}
 P_{f_a}&=-\log2-\psi(z),\label{a03gc:gc:exp-PF}\\
 2H_{f_a}&=\psi(w)+\log(2/a),\label{a03gc:gc:exp-H}\\
 \mathscr L_{f_a}&=\log(a/(2\pi)),\label{a03gc:gc:exp-L}\\
 W_{f_a}&=\frac{2}{a+1/2}+\psi(z)-\log\pi,
 \label{a03gc:gc:exp-W}\\
 2H_{f_a}+\mathscr L_{f_a}&=\psi(w)-\log\pi,
 \label{a03gc:gc:exp-right}\\
 D(f_a)&=\frac{2}{a+1/2}+\psi(z)-\psi(w)
       =\psi(z+1)-\psi(w)>0.
 \label{a03gc:gc:exp-D}
\end{align}
Here $\psi=\Gamma'/\Gamma$ has its usual unscaled argument.

To verify every constant in these identities, use the primary special
function formula \cite[5.9.12]{a03gc:dlmf}
\begin{equation}
 \psi(v)=\int_0^\infty
 \left(\frac{e^{-t}}t-\frac{e^{-vt}}{1-e^{-t}}\right)dt,
 \qquad v>0.
 \label{a03gc:gc:digamma}
\end{equation}
Its integrand is $O(1)$ at zero and exponentially integrable at infinity.
In \eqref{a03gc:gc:PF-formula}, substitute $t=2x$ to get the combined
integrand $e^{-zt}/(1-e^{-t})-e^{-t/2}/t$.
Adding and subtracting $e^{-t}/t$ inside the combined integral gives
$-\psi(z)$ and
$\int_0^\infty(e^{-t}-e^{-t/2})dt/t=-\log2$ by
\eqref{a03gc:gc:Frullani}.  This proves \eqref{a03gc:gc:exp-PF}.
The same substitution in \eqref{a03gc:gc:H} gives
$e^{-at/2}/t-e^{-wt}/(1-e^{-t})$; the difference from
\eqref{a03gc:gc:digamma} integrates to $\log(2/a)$, proving
\eqref{a03gc:gc:exp-H}.
In the transform integral put $t=au$.  The constant term integrates to
$\log(a/(2\pi))$, while the $\log u$ term is zero by $u\mapsto1/u$.
This proves \eqref{a03gc:gc:exp-L}.  The remaining identities follow by
substitution, using the exact recurrence
$\psi(z+1)=\psi(z)+1/z$ \cite[5.5.2]{a03gc:dlmf}; here
$1/z=2/(a+1/2)$.  Positivity follows independently from \eqref{a03gc:gc:D}.

For the incoming test $a=3/2$ one has $z=1$, $w=3/2$.
The special values
$\psi(1)=-\gamma_E$ and $\psi(1/2)=-\gamma_E-2\log2$,
and the recurrence at $1/2$, give
\cite[5.4.12--13 and 5.5.2]{a03gc:dlmf}
\begin{align}
 P_{f_{3/2}}&=\gamma_E-\log2,&
 2H_{f_{3/2}}&=2-\gamma_E-\log3,\notag\\
 \mathscr L_{f_{3/2}}&=\log\frac3{4\pi},&
 W_{f_{3/2}}&=1-\gamma_E-\log\pi,\notag\\
 2H_{f_{3/2}}+\mathscr L_{f_{3/2}}
   &=2-\gamma_E-\log\pi-2\log2.
 \label{a03gc:gc:counterexample-values}
\end{align}
Consequently
\begin{equation}
 \boxed{D(f_{3/2})=2\log2-1>0.}
 \label{a03gc:gc:counterexample}
\end{equation}
There is also a direct evaluation independent of digamma values:
\begin{equation}
 D(f_{3/2})=2\int_0^\infty\frac{e^{-3x}}{1+e^{-x}}dx
 =2\int_0^1\frac{u^2}{1+u}du
 =2\left(\frac12-1+\log2\right).
 \label{a03gc:gc:direct-example}
\end{equation}
Strict positivity also follows from
$\log2=\int_1^2dt/t>1/2$, without a numerical approximation.
Thus the contradiction to the archived equality occurs inside its actual
recorded common domain, and does not depend on the repaired inclusion for
more slowly decaying functions.

\section{Exact supersession and provenance}

Equation \eqref{a03gc:gc:defect} supersedes the claimed zero-defect equality at
the following two archived labels:
\texttt{eq:guinand-weil-arch-map} and
\texttt{eq:guinand-weil-residual-complete}.
The proof that the latter follows from the two displayed local scalar
formulas cannot validate both scalar transcriptions with the printed
definitions: subtracting them gives the contradicted equality.  The
calculation here locates and evaluates the discrepancy exactly, while
preserving both displayed functionals and all their original factors.
It does not determine from this discrepancy alone which published formula
was miscopied.  That would require comparing the historical formulas
themselves, rather than using the local assertion that they were literal.

The source snapshot is
\texttt{06\_Foundations\_and\_Original\_Six\_Readers.tex}, SHA--256
\begin{center}\small\ttfamily
e653a19b9bfc7b4788550e59de42367d52ed666598e46540d040e06d7362cad0.
\end{center}
Its relevant original lines are 201784--201795 (Guinand domain and
transform), 201798--201843 (extension), 201953--202030 (expanded false
comparison), 202639--202660 (Weil conditions and regulator), and
202842--202906 (printed comparison).  The incoming correction is
\texttt{06\_PRINTED\_ARCHIMEDEAN\_CORRECTION.md}, SHA--256
\begin{center}\small\ttfamily
6d0f28606fe98d8ab49f25b54d94377e1481ce83a511ce8189bfebfba629dcee.
\end{center}
The source files are preserved unchanged.  The three functional identities
and the admissible exponential example are established here with complete
proofs; no zero-location conclusion, mixed-support control, or arithmetic
allocation follows merely from this correction.



% END PRESERVED INPUT 06_GUINAND_ARCHIMEDEAN_EXACT_CORRECTION.tex
\endgroup

\clearpage
\begingroup
\def\C{\mathbb C}
\def\Tr{\operatorname{Tr}}
\def\dd{\,d}
\def\im{\operatorname{im}}

\setcounter{equation}{0}
\renewcommand{\theequation}{A03ARU.\arabic{equation}}
\section*{Backward receivers and the original canonical primary-angle profile}
\addcontentsline{toc}{section}{Backward receivers and the original canonical primary-angle profile}
% BEGIN PRESERVED INPUT 07_RECEIVER_UPDATES.tex SHA256 f83a976e6ea8bf4d240906a0fb723a8178270544db89d9fce08f8357fde567f0


\section{The complete canonical scalar and its primary-angle profile}
Keep the full packet and original density of AG1--3 and PT1--5.
Thus $e=1+k(m_0-1)$, $q=e(k+1)^2$, $k\equiv1\pmod4$,
and $A$ is multiplication by $S$ on the original $E=\C[S]/\chi$.
Let $G_N$ be the original attained arithmetic quotient metric, $N\ge q-1$.
The canonical centered weight is
\[
 H_N=G_N^{-1/2}(A^*G_N+G_NA-kG_N)G_N^{-1/2}.
\]
The exact two-boundary-row formula PT23 gives rank at most two;
reflection in the original full polynomial gives the opposite signs,
and PT7 makes both nonzero. Thus its nonzero eigenvalues are
$+\epsilon_N,-\epsilon_N$. Let $u_{N,+},u_{N,-}$ be
unit eigenvectors and let $\Pi_N$ be the Euclidean orthogonal projection onto
$G_N^{1/2}E_+$, where $E_+$ is the complete positive-real-displacement
primary subspace of $A$. All primary orders remain in this definition.
PT7 and PT24 give
\begin{equation}\tag{A3R1}
 L_{h,k}=\frac{\delta q(k+1)}2
 =\Tr(\Pi_NH_N)
 =\epsilon_N\bigl(\|\Pi_Nu_{N,+}\|^2-\|\Pi_Nu_{N,-}\|^2\bigr).
\end{equation}
Set
\begin{equation}\tag{A3R2}
 \mathfrak d_N=\|(I-\Pi_N)u_{N,+}\|^2+\|\Pi_Nu_{N,-}\|^2,
 \qquad \mathfrak t_N=1-\mathfrak d_N=\frac{L_{h,k}}{\epsilon_N}.
\end{equation}
These are the original primary alignment defects. The equality follows
by subtracting the first squared projection from one in (A3R1).
The positive-trace minimum proved in PT8--16 gives
$0<\mathfrak t_N\le1$ and $0\le \mathfrak d_N<1$. Neither eigenvector phase changes them.

The original two-window action, before any asymptotic expansion, is
\begin{equation}\tag{A3R3}
 J_k^{\rm action}=\sum_{N=q-1}^{2q-1}w_N\log\epsilon_N^2,
 \qquad w_{q-1}=w_{2q-1}=1,\quad w_N=2\ (q\le N\le2q-2).
\end{equation}
There are exactly $2q$ weighted terms. This is CAI13--15, with the
original adjacent windows $[q-1,2q-2]$ and $[q,2q-1]$ both present.
Consequently its surplus over the positive-primary benchmark is exactly
\begin{equation}\tag{A3R4}
 J_k^{\rm action}-4q\log L_{h,k}
 =-2\sum_Nw_N\log(1-\mathfrak d_N),\qquad
 \left(\prod_N\mathfrak t_N^{w_N}\right)^{1/(2q)}
 =\exp\left[-\frac{J_k^{\rm action}-4q\log L_{h,k}}{4q}\right].
\end{equation}
This proves a metric and action morphism from the original full primary
projection to the computed scalar, rather than equating it with the
different original observation-kernel projection.

Use the exact constants of AG5, AG9, AG24 and AG26, and put
$C_A=C_\sigma+8\log2-4-(4m_0-2)C_\Gamma-4\log(\delta/2)$.
The newly evaluated AG31 now gives the original weighted geometric mean
to inverse-logarithmic precision:
\begin{align}\tag{A3R5}
 \frac1{2q}\sum_Nw_N\log \mathfrak t_N
 ={}&-\frac{C_Bq}{4}+\log(k+1)-\frac{C_A}{4}
 -\frac{\log2}{4\pi}\,I_h\frac{k^2}{q}
 +\frac{\eta_hM}{2}\frac{k(k+2)}q\\
 &-\frac{C_\Gamma}{8(\log q+c_\zeta)}+o_h(1/\log q).\notag
\end{align}
This is direct division of (AG31) by $-4q$ in the exact identity (A3R4).
The finite version retains the full SR remainder, the $O_h(\log q)$ Gamma
bound, and the CAI upper/lower interval AG28. The individual $\mathfrak t_N$ and
the original eight observation-flag allocations are not recovered by
a geometric mean. Their maps remain those of CRP/CRS and JRE.

PT41 gives a further exact receiver. Here $D_N$ is the off-diagonal action
block between $E_+$ and its actual $G_N$-orthogonal complement.
In the original primary coordinates the latter is the graph
$(-K_Nv,v)$ over $E_-$, with $K_N=P_N^{-1}X_N$ and the full positive
metric blocks of PT36; this is not an assertion that $E_+$ and $E_-$
are already orthogonal. Use the exact orthonormal frame of PT36;
then, if $\|\cdot\|_1$ denotes the sum of its singular values,
\begin{equation}\tag{A3R6}
 \epsilon_N^2\ge L_{h,k}^2+\|D_N\|_1^2,\qquad
 \sum_Nw_N\log\left(1+\frac{\|D_N\|_1^2}{L_{h,k}^2}\right)
 \le J_k^{\rm action}-4q\log L_{h,k}.
\end{equation}
Taking logarithms of each positive ratio in PT41 and summing the actual
weights proves the second inequality. It supplies an upper budget for
these couplings; a large total alone does not give a lower bound on any one
coupling. The explicit Sylvester inverse in PT35--41 retains all root
differences and nilpotent powers for the next such calculation.

The new angle notation $\mathfrak d_N$ is distinct from CAI's original
volume contraction. Write the latter explicitly as
$r_n=\det G_n/\det G_{n-1}$, and keep its original recurrence coefficient
$b_n^{\rm rec}=\omega_n/\omega_{n-1}$. CAI12 and (A3R2) give the exact
connecting map
\begin{equation}\tag{A3R6a}
 \mathfrak t_{q-1}^2
 =\frac{L_{h,k}^2}{b_q^{\rm rec}(r_q^{-1}-1)},\qquad
 \mathfrak t_N^2
 =\frac{L_{h,k}^2}
 {b_{N+1}^{\rm rec}(1-r_N)(r_{N+1}^{-1}-1)}\quad(N\ge q).
\end{equation}
Every denominator is positive on the original full packet by the nonzero
canonical weight and CAI12. This connects the computed angle product to
the actual attained covariance increments without equating the two defects.

\section{The original section and every retained metric cross block}
For the coefficient programme retain its actual degree $a=k+4$, simple
quartet, $q=(a+1)^2$, $q'=(a-7)^2$, $m=q-q'-v$, and actual selected
invertible conductor row set $I$. Let $I_0$ be the consecutive row set.
The exact matrix identity COI9--11 is
\begin{equation}\tag{A3R7}
 E_{J_I}=J_vD_{0I}+E_{J_0}M_{0I},\qquad
 L_I=F_CD_{0I}+L_0M_{0I},\qquad M_{0I}\text{ invertible}.
\end{equation}
Here $L_I$ is the original arithmetic residue injection, $F_C$ is its
conductor injection, and $L_0$ is the computed consecutive injection.
For every positive arithmetic form $G$ on the full packet, this gives
\begin{align}\tag{A3R8}
 L_I^*GL_I={}&M_{0I}^*L_0^*GL_0M_{0I}
 +D_{0I}^*F_C^*GF_CD_{0I}\\
 &+M_{0I}^*L_0^*GF_CD_{0I}
 +D_{0I}^*F_C^*GL_0M_{0I}.\notag
\end{align}
The proof is multiplication of (A3R7) and its adjoint. The same expansion
holds for either original boundary difference and for the full action form.
In contrast, division by the conductor translates the native minimizing
fibre by an admitted lower exponential. That translation is a bijection,
so its unchanged minimum gives exactly
\begin{equation}\tag{A3R9}
 D_{I,a,N-v}^{(s)}=M_{0I}^*D_{0,a,N-v}^{(s)}M_{0I},\qquad s=1,a.
\end{equation}
These two equations are the actual relation between the presentations.
The absence of the conductor cross terms in (A3R9) is justified by its
specific quotient, and does not remove them from (A3R8).

With four copies set $\mathscr G_N=(I_4\otimes L_I)^*
\operatorname{diag}(G_N,4)(I_4\otimes L_I)$ and
$\mathscr D_N=\operatorname{diag}(D_{I,a,N-v}^{(s)},4)$.
Let $\mathcal R$ have signs $(-,-,+,+)$ at $N=q-1,q,2q-1,2q$ and set
$T_N=\mathscr G_N^{-1/2}\mathscr D_N\mathscr G_N^{-1/2}$.
The full native total is $\mathcal T=\mathcal R\log\det D_{I,a,N-v}^{(s)}$.
Taking determinants before their signed sum proves
\begin{equation}\tag{A3R10}
 \mathcal R\log\det T_N=4\mathcal T+F_{\rm ar},\qquad
 F_{\rm ar}=-\mathcal R\log\det\mathscr G_N\ge0.
\end{equation}
Monotonicity follows by minimizing over the nested original relation spaces,
and persists under the fixed injection. It proves the last sign separately
for the two original endpoint links. For $I=I_0$, CAL1--18 bounds
$F_{\rm ar}\le-8m\log(1-\varepsilon_a)$ with its explicitly estimated
exponentially small $\varepsilon_a$. For every other $I$, the exact full
arithmetic value is retained in (A3R8) and NMR's complete projection formula.
No root-column ordering selects a row-subset ordering by itself.

If $b_-\mathscr G_N\preceq\mathscr D_N\preceq b_+\mathscr G_N$
at all four endpoints, then all $4m$ eigenvalues of $T_N$ lie in
$[b_-,b_+]$. Thus each signed pair costs at most
$4m\log(b_+/b_-)$, and (A3R10) gives
\begin{equation}\tag{A3R11}
 \log(b_+/b_-)\ge\frac{4\mathcal T+F_{\rm ar}}{8m}.
\end{equation}
This retains the full relative metric term at every original section.

\section{Exact negative-trace and period receivers}
RBV1--36 evaluates an explicit degree-$8a$ polynomial $P_a$ in the
original full packet. If the retained section is $I_0$, its exact preimage
is the Laurent coefficient vector of COI7. For any other $I$,
solve the full square frame $[J_v,E_{J_I}]$ for the moment vector of $P_a$;
this returns both its conductor component and its native component.
Substitution in (A3R8) gives the full energy with both cross terms.
This is the receiver for both boundary windows; it does not assert that
the native component alone represents $P_a$.

The complete tensor-residue comparison on the same full packet is
\begin{equation}\tag{A3R12}
 \mathcal W(f,g)=\mathcal B(f,a_{\rm tr}g)+\mathcal W_Z(f,g),
 \qquad a_{\rm tr}=\sum_{c_\lambda\ne0}\frac{e_\lambda}{c_\lambda},
\end{equation}
where $c_\lambda$ is the actual tensor coefficient of the original
four-term polynomial using $v_h(\zeta)/h'(\zeta)$, and $e_\lambda$ is the
original CRT idempotent. The residual is the sum over $c_\lambda=0$.
This identity follows site by site: a nonzero $c_\lambda$ cancels its
displayed reciprocal, and a zero coefficient is supplied by the residual.
RBV31--36 retains the full unit, proves both terms nonpositive on $P_a$,
and obtains a strict negative tensor value from the perimeter sites.
The second argument is the displayed full-packet polynomial; (A3R7)--(A3R8)
give its exact further receiver whenever a coefficient presentation is used.

UPR1--31 constructs the fixed homogeneous module
\[
\mathbb C[X_1,X_2,X_3,X_4]^{125}\oplus
\mathbb C[X_1,X_2,X_3,X_4]^{625}\longrightarrow
\mathbb C[X_1,X_2,X_3,X_4]^{625},
\]
with $(f,g)\mapsto f+Qg$,
with $X_i=x_i^5$ and original shifted degrees.
Multiplication by the quadratic may produce $X_iX_j$ from a residue
having both exponents four, so the coefficient degree bound is two.
The fixed decreasing rewrite score still gives a homogeneous image inverse
of coefficient norm $\exp(O_{\rm actual}(k))$.
The original word Gram and Lagrange map then give the actual period-to-
arithmetic image quotient bound $\exp(O_{h,u}(q\log(q+2)))$.
The map is defined on the original image, with the original kernel left
in its quotient, so it does not set an observation-kernel restriction to zero.

\section{Correct the archived scalar identities in their original domains}
The MA1--26 proof completes the original Montgomery finite part on
$1<\sigma_0<2$, $x\ge1$, $t\in\mathbb R$, retaining its jump and both poles.
Its exact replacement equation and full proof are incorporated alongside
this receiver. GC proves that the two archived Guinand--Weil functionals
are related by the additional linear functional
\begin{equation}\tag{A3R13}
 D(f)=\int_0^\infty f(x)
 [2e^{-x/2}-\operatorname{sech}(x/2)]\dd x
 =\int_0^\infty f(x)\frac{2e^{-3x/2}}{1+e^{-x}}\dd x.
\end{equation}
The equality of kernels follows by writing
$\operatorname{sech}(x/2)=2e^{-x/2}/(1+e^{-x})$.
On $f=e^{-3x/2}$ the substitution $u=e^{-x}$ gives
$D(f)=2\int_0^1u^2/(1+u)\dd u=2\log2-1>0$.
Consequently the previous equality without $D(f)$ must be replaced at
both named source equations; it cannot be used as a zero-defect premise.
The corrected identity has the original common-domain inclusion for decay
exponent greater than $1/2$, while GC proves the functional identity for
positive exponential decay. The updated foundation derivatives preserve
both original source files and record the exact replacement patches.


% END PRESERVED INPUT 07_RECEIVER_UPDATES.tex
\endgroup

\endgroup
